\appendix

\section{Proof Details}\label{sec:proof-details}

\subsection{Realized-factor conditioning}\label{app:step-001}

\begin{proof}[Auxiliary facts in the present notation]


\medskip\noindent\textbf{Rectangular Gaussian operator-norm tail.}\par

We use the following statement in the present notation. If \(G\in\mathbb R^{n\times r}\) has independent
\(\mathcal N(0,1)\) entries, then for every \(u\ge0\),
\[
  \mathbb P\!\left(\|G\|_{\rm op}\ge \sqrt n+\sqrt r+u\right)
  \le e^{-u^2/2}.
\]
Equivalently, if \(\Xi=(\rho/\sqrt n)G\), so every entry of \(\Xi\) has
variance exactly \(\rho^2/n\), then
\[
  \mathbb P\!\left(
    \|\Xi\|_{\rm op}
    \ge \rho\left(1+\sqrt{r/n}+u/\sqrt n\right)
  \right)
  \le e^{-u^2/2}.
\]

For each mode \(M\in\{A,B,C\}\), let
\(\Xi_M=M-\bar M\). Under
Assumption~\ref{assump:gaussian_smoothing},
\(G_M=(\sqrt n/\rho)\Xi_M\) has independent standard-normal entries, with
\(\rho=r^{-q}\).

The entry independence and variance convention are exactly those of
Assumption~\ref{assump:gaussian_smoothing}. The result is applied with
\[
  u_r=\sqrt{2\log(3r^{20})},
\]
giving per-mode failure \((3r^{20})^{-1}\). No independence between the three
mode events is needed for the later union bound.


\medskip\noindent\textbf{Weyl singular-value perturbation inequality.}\par

We use the following statement in the present notation. For matrices \(X,E\) of the same size,
\[
  \sigma_{\min}(X+E)\ge \sigma_{\min}(X)-\|E\|_{\rm op},
  \qquad
  \|X+E\|_{\rm op}\le \|X\|_{\rm op}+\|E\|_{\rm op}.
\]

Take \(X=\bar M\), \(E=\Xi_M\), and
\(X+E=M\) for each \(M\in\{A,B,C\}\).

Lemma~\ref{lem:step-001-base-conditioning}
supplies the singular bounds for \(\bar M\), and
Lemma~\ref{lem:step-001-gaussian-perturbation} supplies the operator bound for
\(\Xi_M\). The resulting bounds are exactly those defining
\(\mathcal E_{\rm cond}\).

\end{proof}



\begin{lemma}[conditioning of the deterministic base matrices]
\label{lem:step-001-base-conditioning}
Under Assumption~\ref{assump:base_conditioning}, for every
\(\bar M\in\{\bar A,\bar B,\bar C\}\),
\[
  \sigma_{\min}(\bar M)\ge \kappa^{-2},
  \qquad
  \|\bar M\|_{\rm op}\le \kappa^2.
\]

\end{lemma}

\begin{proof}
Write, exactly as in the setting,
\[
  \bar M=\bar M^\circ D_M,
  \qquad
  D_M=\operatorname{diag}(\|\bar m_1\|_2,\ldots,
  \|\bar m_r\|_2).
\]
Assumption~\ref{assump:base_conditioning} makes \(D_M\) invertible and gives
\[
  \sigma_{\min}(\bar M^\circ)\ge\kappa^{-1},\quad
  \|\bar M^\circ\|_{\rm op}\le\kappa,\quad
  \sigma_{\min}(D_M)\ge\kappa^{-1},\quad
  \|D_M\|_{\rm op}\le\kappa.
\]
For every \(x\in\mathbb R^r\),
\[
  \|\bar Mx\|_2
  =\|\bar M^\circ D_Mx\|_2
  \ge \sigma_{\min}(\bar M^\circ)\|D_Mx\|_2
  \ge \kappa^{-2}\|x\|_2,
\]
which proves the lower singular bound. Submultiplicativity gives
\[
  \|\bar M\|_{\rm op}
  \le \|\bar M^\circ\|_{\rm op}\|D_M\|_{\rm op}
  \le \kappa^2.
\]
This proves the lemma. 

\end{proof}


\begin{lemma}[simultaneous Gaussian perturbation control]
\label{lem:step-001-gaussian-perturbation}
Under Assumptions~\ref{assump:dimension} and
\ref{assump:gaussian_smoothing}, define
\[
  r_{0,\rm cond}(\kappa,q)
  :=\max\left\{3,\left\lceil(6\kappa^2)^{1/q}\right\rceil\right\}.
\]
Choose the dimension constant in Assumption~\ref{assump:dimension} so that
\(C(\kappa,q)\ge1\). Then, for every integer
\(r\ge r_{0,\rm cond}(\kappa,q)\) and every
\(n\ge C(\kappa,q)r^4\log r\), the event
\[
  \mathcal E_{\rm pert}
  :=\left\{
    \max_{M\in\{A,B,C\}}\|M-\bar M\|_{\rm op}
    \le \frac1{2\kappa^2}
  \right\}
\]
satisfies
\[
  \mathbb P(\mathcal E_{\rm pert})\ge1-r^{-20}.
\]
The probability is uniform over all deterministic base triples satisfying the
setting assumptions and over all such \(n\).

\end{lemma}

\begin{proof}
Put \(\rho=r^{-q}\),
\(\Xi_M=M-\bar M\), and
\[
  u_r=\sqrt{2\log(3r^{20})}.
\]
The checked Gaussian operator-norm tail gives, for each of the three modes,
\[
  \mathbb P\!\left(
    \|\Xi_M\|_{\rm op}
    >\rho\left(1+\sqrt{r/n}+u_r/\sqrt n\right)
  \right)
  \le \frac1{3r^{20}}.                                      \tag{1}
\]
For \(r\ge3\),
\[
  u_r^2=2\log 3+40\log r
  \le42\log r
  \le r^4\log r
  \le n,                                                     \tag{2}
\]
where the penultimate inequality uses \(r^4\ge81>42\). Also
\(n\ge r^4\log r\ge r\). Hence
\[
  \sqrt{r/n}\le1,
  \qquad u_r/\sqrt n\le1.                                   \tag{3}
\]
Combining (1)--(3), outside a per-mode event of probability at most
\((3r^{20})^{-1}\),
\[
  \|\Xi_M\|_{\rm op}
  \le3\rho=3r^{-q}.
\]
The definition of \(r_{0,\rm cond}(\kappa,q)\) gives
\[
  3r^{-q}\le\frac1{2\kappa^2}.                              \tag{4}
\]
A union bound over \(M=A,B,C\), which does not require independence, yields
\[
  \mathbb P(\mathcal E_{\rm pert}^c)
  \le3\cdot\frac1{3r^{20}}=r^{-20}.                         \tag{5}
\]
The right sides of (1)--(5) do not depend on the deterministic base triple.
Increasing \(n\) only decreases \(\sqrt{r/n}\) and \(u_r/\sqrt n\), so the
argument has no upper-dimensional restriction. This proves the lemma.


\end{proof}


\begin{proposition}[realized-factor conditioning]
\label{prop:step-001-realized-conditioning}
Under Assumptions~\ref{assump:base_conditioning},
\ref{assump:dimension}, and \ref{assump:gaussian_smoothing}, let
\(\kappa_1=2\kappa^2\), choose \(C(\kappa,q)\ge1\), and suppose
\(r\ge r_{0,\rm cond}(\kappa,q)\). Then, simultaneously for
\(M\in\{A,B,C\}\), with probability at least \(1-r^{-20}\),
\[
  \sigma_{\min}(M)\ge\kappa_1^{-1},
  \qquad
  \|M\|_{\rm op}\le\kappa_1,
  \qquad
  \|M^\dagger\|_{\rm op}\le\kappa_1.
\]
In particular, the generated event \(\mathcal E_{\rm cond}\) holds with
probability at least \(1-r^{-20}\).

\end{proposition}

\begin{proof}
On \(\mathcal E_{\rm pert}\), the checked Weyl
inequality and Lemma~\ref{lem:step-001-base-conditioning} give, for every mode,
\[
  \sigma_{\min}(M)
  \ge\kappa^{-2}-\frac1{2\kappa^2}
  =\frac1{2\kappa^2}
  =\kappa_1^{-1}.                                           \tag{6}
\]
The operator-norm triangle inequality gives
\[
  \|M\|_{\rm op}
  \le\kappa^2+\frac1{2\kappa^2}
  \le2\kappa^2
  =\kappa_1,                                                \tag{7}
\]
where the second inequality uses \(\kappa\ge1\). Equation (6) makes every
realized mode matrix full column rank and therefore
\[
  \|M^\dagger\|_{\rm op}
  =\sigma_{\min}(M)^{-1}
  \le\kappa_1.                                              \tag{8}
\]
Equations (6)--(7) show that
\(\mathcal E_{\rm pert}\subseteq\mathcal E_{\rm cond}\). By
Lemma~\ref{lem:step-001-gaussian-perturbation},
\(\mathbb P(\mathcal E_{\rm pert})\ge1-r^{-20}\), so the same lower bound
holds for \(\mathbb P(\mathcal E_{\rm cond})\). This proves the proposition.

\end{proof}

\begin{proof}[Proof of the subsection conclusion]

Lemma~\ref{lem:step-001-base-conditioning} converts
Assumption~\ref{assump:base_conditioning} into the deterministic interval
\([\kappa^{-2},\kappa^2]\) for each base mode matrix.
Lemma~\ref{lem:step-001-gaussian-perturbation} uses the exact smoothing
variance \(r^{-2q}/n\), the explicit threshold
\(u_r=\sqrt{2\log(3r^{20})}\), and a union over the three modes to produce
\(\mathcal E_{\rm pert}\) with failure at most \(r^{-20}\), uniformly over
all admissible bases and all \(n\) above the lower dimension threshold.
Proposition~\ref{prop:step-001-realized-conditioning} then applies Weyl to
obtain, with \(\kappa_1=2\kappa^2\),
\[
  \sigma_{\min}(M)\ge\kappa_1^{-1},\qquad
  \|M\|_{\rm op}\le\kappa_1
  \quad(M=A,B,C),
\]
on an event of probability at least \(1-r^{-20}\). These are precisely the
conditions defining \(\mathcal E_{\rm cond}\), which proves the proposition.
No initialization-geometry, tangent, or trajectory claim is used
or proved here.

\end{proof}

\subsection{Balanced coefficient Gaussianization}\label{app:step-002}

\begin{proof}[Auxiliary facts in the present notation]


\medskip\noindent\textbf{realized-factor conditioning.}\par

The preceding Proposition~\ref{prop:step-001-realized-conditioning} states that
on the generated event \(\mathcal E_{\rm cond}\), every realized mode matrix
has full column rank and
\[
  \sigma_{\min}(M)\ge\kappa_1^{-1},\qquad
  \|M\|_{\rm op}\le\kappa_1,
  \qquad M\in\{A,B,C\}.
\]
Consequently \(M^\dagger\) exists and
\(\|M^\dagger\|_{\rm op}\le\kappa_1\).

The argument conditions on exactly the
event produced by that proposition. The compact SVDs and pseudoinverses below
are therefore legal for the same realized matrices \(A,B,C\).


\medskip\noindent\textbf{Gaussian polar decomposition.}\par

We use the following statement in the present notation. If \(g\sim\mathcal N(0,I_n)\), then
\(\chi=\|g\|_2>0\) almost surely and, on that event,
\(\omega=g/\chi\) is uniform on \(S^{n-1}\), independent of \(\chi\). For an
independent family of standard Gaussian vectors, the corresponding
radius-direction pairs are independent.

Apply the result to
\(g_i^x=\sqrt n\,\widetilde x_i\),
\(g_i^y=\sqrt n\,\widetilde y_i\), and
\(g_i^z=\sqrt n\,\widetilde z_i\). Assumption~\ref{assump:independent_initialization}
gives exactly the required iid standard Gaussian law. The result is also
derived directly from the radial Gaussian density in
Lemma~\ref{lem:step-002-balancing-scalars}.


\medskip\noindent\textbf{Orthogonal images of a standard Gaussian and compact SVD.}\par

We use the following statement in the present notation. If \(g\sim\mathcal N(0,I_n)\) and
\(U\in\mathbb R^{n\times r}\) satisfies \(U^{\mathsf T}U=I_r\), then
\(U^{\mathsf T}g\sim\mathcal N(0,I_r)\). Applying deterministic linear maps
separately to independent Gaussian vectors preserves their independence. If
\(M=U\Sigma V^{\mathsf T}\) is a full-column-rank compact SVD, with
\(U^{\mathsf T}U=V^{\mathsf T}V=I_r\) and positive diagonal \(\Sigma\), then
\[
  M^\dagger=V\Sigma^{-1}U^{\mathsf T}.
\]

Conditional on \(A,B,C\), the compact
SVD factors are deterministic. Apply the Gaussian statement to the distinct
raw vectors \(g_i^m\), and define
\(z_i^M=U_M^{\mathsf T}g_i^m\) and
\(H_M=V_M\Sigma_M^{-1}\). Full column rank is supplied by
Proposition~\ref{prop:step-001-realized-conditioning}. These are the exact
objects used in
Lemma~\ref{lem:step-002-coefficient-gaussianization}.

\end{proof}



\begin{lemma}[Gaussian radii and exact balancing scalars]
\label{lem:step-002-balancing-scalars}
Under Assumption~\ref{assump:independent_initialization}, define, for every
component \(i\),
\[
  g_i^x=\sqrt n\,\widetilde x_i,
  \qquad g_i^y=\sqrt n\,\widetilde y_i,
  \qquad g_i^z=\sqrt n\,\widetilde z_i,
  \qquad \chi_i^m=\|g_i^m\|_2.
\]
Then the \(g_i^m\) are independent \(\mathcal N(0,I_n)\) vectors over all
\(i\) and \(m\). Almost surely, every \(\chi_i^m\) is positive, the directions
\(\omega_i^m=g_i^m/\chi_i^m\) are independent uniform sphere vectors and are
independent of all radii, and the balancing map satisfies
\[
  x_{i,0}=s_i^x\widetilde x_i,\qquad
  y_{i,0}=s_i^y\widetilde y_i,\qquad
  z_{i,0}=s_i^z\widetilde z_i,
\]
where, on the positive-radius branch,
\[
  s_i^x=\left(\frac{\chi_i^y\chi_i^z}{(\chi_i^x)^2}\right)^{1/3},
  \quad
  s_i^y=\left(\frac{\chi_i^x\chi_i^z}{(\chi_i^y)^2}\right)^{1/3},
  \quad
  s_i^z=\left(\frac{\chi_i^x\chi_i^y}{(\chi_i^z)^2}\right)^{1/3}.
\]
On the setting's zero-factor branch, define
\(s_i^x=s_i^y=s_i^z=1\). With this extension every scalar is nonzero, the
displayed scalar representation agrees with \(\mathcal G\) on every branch,
and
\[
  s_i^x s_i^y s_i^z=1
\]
identically.

\end{lemma}

\begin{proof}
Assumption~\ref{assump:independent_initialization}
gives the stated joint standard Gaussian law after multiplication by
\(\sqrt n\). The density of one \(g_i^m\) is a function only of its Euclidean
radius. In polar coordinates its measure is proportional to
\[
  e^{-\chi^2/2}\chi^{n-1}\,d\chi\,d\sigma(\omega),
\]
where \(d\sigma\) is uniform surface measure. This factorization proves the
radius-direction independence and uniformity. The singleton \(\{0\}\) has
Gaussian measure zero; a finite union over the \(3k\) raw columns therefore
shows that all radii are positive almost surely.

On that probability-one branch, the three raw physical norms are
\(\chi_i^x/\sqrt n,\chi_i^y/\sqrt n,\chi_i^z/\sqrt n\), so their geometric
mean is
\[
  q_i=\frac{(\chi_i^x\chi_i^y\chi_i^z)^{1/3}}{\sqrt n}.
\]
The balancing definition gives
\(s_i^x=q_i/(\chi_i^x/\sqrt n)\), and similarly in the other modes, which is
exactly the displayed formula. In particular, balancing changes only radii:
\(x_{i,0}=q_i\omega_i^x\), and analogously for \(y,z\), so all raw directions
are preserved. Direct multiplication gives
\[
  (s_i^x s_i^y s_i^z)^3
  =\frac{(\chi_i^y\chi_i^z)(\chi_i^x\chi_i^z)
          (\chi_i^x\chi_i^y)}
         {(\chi_i^x\chi_i^y\chi_i^z)^2}
  =1.
\]
All three scalars are positive, hence their product is one rather than another
cube root of one. If any raw factor is zero, the setting leaves the entire raw
triple unchanged; assigning all three scalars the value one represents that
branch exactly and preserves the product identity. This proves the lemma.


\end{proof}


\begin{lemma}[compact-SVD Gaussianization of the balanced coefficients]
\label{lem:step-002-coefficient-gaussianization}
Under Assumption~\ref{assump:independent_initialization}, Proposition~\ref{prop:step-001-realized-conditioning}, and
Lemma~\ref{lem:step-002-balancing-scalars}, condition on any realized triple
\((A,B,C)\in\mathcal E_{\rm cond}\). For each
\(M\in\{A,B,C\}\), choose a compact SVD
\[
  M=U_M\Sigma_MV_M^{\mathsf T},
  \qquad U_M^{\mathsf T}U_M=V_M^{\mathsf T}V_M=I_r,
  \qquad \Sigma_M\succ0,
\]
and define
\[
  H_M=V_M\Sigma_M^{-1},
  \qquad z_i^M=U_M^{\mathsf T}g_i^m,
\]
using the mode pairing \((A,x),(B,y),(C,z)\). Then, under the conditional
initialization law, the full family
\(\{z_i^M:M\in\{A,B,C\},i\in[k]\}\) is independent with
\(z_i^M\sim\mathcal N(0,I_r)\), every \(H_M\) is invertible and satisfies
\[
  \sigma(H_M)\subset[\kappa_1^{-1},\kappa_1],
\]
and the raw, balanced, and normalized coefficient vectors obey the exact
identities
\[
  \widetilde\zeta_i^M=\frac1{\sqrt n}H_Mz_i^M,
  \qquad
  \zeta_i^M=\frac{s_i^m}{\sqrt n}H_Mz_i^M,
  \qquad
  \boxed{\bar\zeta_i^M=\frac{s_i^m}{\sqrt r}H_Mz_i^M}.
\]
No independence between \(s_i^m\) and \(z_i^M\) is asserted or needed.

\end{lemma}

\begin{proof}
Conditional on the realized factors, all SVD
matrices are deterministic, while the standardized raw initialization vectors
remain independent standard Gaussians because initialization is independent
of smoothing. Since \(U_M^{\mathsf T}U_M=I_r\),
\[
  \mathbb E[z_i^M(z_i^M)^{\mathsf T}\mid A,B,C]
  =U_M^{\mathsf T}U_M=I_r,
\]
and each \(z_i^M\) is Gaussian. Distinct \((M,i)\) use distinct independent
raw vectors, proving independence across all modes and components.

The compact-SVD pseudoinverse formula gives
\[
  M^\dagger=V_M\Sigma_M^{-1}U_M^{\mathsf T}.
\]
Because \(\widetilde m_i=g_i^m/\sqrt n\),
\[
  \widetilde\zeta_i^M
  =M^\dagger\widetilde m_i
  =\frac1{\sqrt n}V_M\Sigma_M^{-1}U_M^{\mathsf T}g_i^m
  =\frac1{\sqrt n}H_Mz_i^M.                 \tag{1}
\]
Lemma~\ref{lem:step-002-balancing-scalars} gives
\(m_{i,0}=s_i^m\widetilde m_i\) on every branch under its scalar extension.
Linearity of \(M^\dagger\) therefore yields
\[
  \zeta_i^M=s_i^m\widetilde\zeta_i^M
  =\frac{s_i^m}{\sqrt n}H_Mz_i^M.           \tag{2}
\]
Multiplication by the setting's coefficient normalization
\(\sqrt{n/r}\) gives the boxed formula, with no omitted factor.

The singular values of \(H_M=V_M\Sigma_M^{-1}\) are exactly the reciprocals
of the singular values of \(M\). On \(\mathcal E_{\rm cond}\),
\(\sigma(M)\subset[\kappa_1^{-1},\kappa_1]\), and this interval is invariant
under reciprocal, proving the claimed bounds for \(H_M\). Repeated singular
values, permutations, or SVD sign choices only apply orthogonal changes to
the standard Gaussian coordinates and do not alter their joint law or the
singular values of \(H_M\). Thus any fixed compact-SVD convention, including
any measurable convention used below, gives the same identities. Finally, the
scalars are functions of the full raw radii and may be dependent on the
projected Gaussian coordinates; none of the algebraic conclusions requires
independence. This proves the lemma.

\end{proof}


\begin{proposition}[exact tangent-block and coefficient-product invariance]
\label{prop:step-002-balancing-invariance}
Under Assumption~\ref{assump:independent_initialization}, Proposition~\ref{prop:step-001-realized-conditioning}, and
Lemmas~\ref{lem:step-002-balancing-scalars} and
\ref{lem:step-002-coefficient-gaussianization}, condition on any realized
triple in \(\mathcal E_{\rm cond}\). For each component define the three
elliptic-Gaussian baseline tangent blocks
\[
  \begin{aligned}
  \mathcal B_i^x
  &=\{u\otimes H_Bz_i^B\otimes H_Cz_i^C:u\in\mathbb R^r\},\\
  \mathcal B_i^y
  &=\{H_Az_i^A\otimes v\otimes H_Cz_i^C:v\in\mathbb R^r\},\\
  \mathcal B_i^z
  &=\{H_Az_i^A\otimes H_Bz_i^B\otimes w:w\in\mathbb R^r\}.
  \end{aligned}
\]
Then every tangent block in the setting, both raw and normalized, has exactly
the corresponding baseline range. More precisely,
\[
  \begin{aligned}
  \{u\otimes\beta_{i,0}\otimes\gamma_{i,0}:u\in\mathbb R^r\}
    &=\frac{s_i^ys_i^z}{n}\mathcal B_i^x=\mathcal B_i^x,\\
  \{\alpha_{i,0}\otimes v\otimes\gamma_{i,0}:v\in\mathbb R^r\}
    &=\frac{s_i^xs_i^z}{n}\mathcal B_i^y=\mathcal B_i^y,\\
  \{\alpha_{i,0}\otimes\beta_{i,0}\otimes w:w\in\mathbb R^r\}
    &=\frac{s_i^xs_i^y}{n}\mathcal B_i^z=\mathcal B_i^z,
  \end{aligned}                                             \tag{3}
\]
and
\[
  \begin{aligned}
  \{u\otimes\bar\beta_{i,0}\otimes\bar\gamma_{i,0}:u\in\mathbb R^r\}
    &=\frac{s_i^ys_i^z}{r}\mathcal B_i^x=\mathcal B_i^x,\\
  \{\bar\alpha_{i,0}\otimes v\otimes\bar\gamma_{i,0}:v\in\mathbb R^r\}
    &=\frac{s_i^xs_i^z}{r}\mathcal B_i^y=\mathcal B_i^y,\\
  \{\bar\alpha_{i,0}\otimes\bar\beta_{i,0}\otimes w:w\in\mathbb R^r\}
    &=\frac{s_i^xs_i^y}{r}\mathcal B_i^z=\mathcal B_i^z.
  \end{aligned}                                             \tag{4}
\]
Consequently,
\[
  \mathscr S_0^{\rm raw}=\mathscr S_0^{\rm norm}
  =\operatorname{span}_{i\in[k]}
    (\mathcal B_i^x\cup\mathcal B_i^y\cup\mathcal B_i^z)
  =\mathscr S_0.                                            \tag{5}
\]
Moreover balancing preserves the exact raw coefficient tensor componentwise:
\[
  \widehat D_0
  =\sum_{i=1}^k\alpha_{i,0}\otimes\beta_{i,0}\otimes\gamma_{i,0}
  =\sum_{i=1}^k\widetilde\zeta_i^A\otimes
                    \widetilde\zeta_i^B\otimes
                    \widetilde\zeta_i^C
  =\frac1{n^{3/2}}\sum_{i=1}^k
       H_Az_i^A\otimes H_Bz_i^B\otimes H_Cz_i^C.           \tag{6}
\]
The normalization is certificate-only:
\[
  \sum_{i=1}^k\bar\alpha_{i,0}\otimes\bar\beta_{i,0}
                    \otimes\bar\gamma_{i,0}
  =\left(\frac nr\right)^{3/2}\widehat D_0,                \tag{7}
\]
so neither \(\widehat D_0\) nor the fixed raw target \(D_r\) is replaced by a
normalized tensor.

\end{proposition}

\begin{proof}
Lemma~\ref{lem:step-002-coefficient-gaussianization}
gives
\[
  \beta_{i,0}=\frac{s_i^y}{\sqrt n}H_Bz_i^B,
  \qquad
  \gamma_{i,0}=\frac{s_i^z}{\sqrt n}H_Cz_i^C.
\]
Substitution into the first raw tangent family yields the first equality in
(3). The other two follow by the same substitution in their displayed modes.
The normalized coefficient formula replaces each denominator \(\sqrt n\) by
\(\sqrt r\), giving (4). By
Lemma~\ref{lem:step-002-balancing-scalars}, every pair product appearing in
(3)--(4) is nonzero; also \(n,r>0\). Multiplication of a linear subspace by a
nonzero scalar is a bijection of that subspace, proving the second equality in
each line. Taking the span over all components and all three mode blocks proves
(5). This argument remains valid when a displayed baseline block happens to
be the zero subspace.

For each component, the same Gaussianization and the exact product-one identity
give
\[
  \begin{aligned}
  \alpha_{i,0}\otimes\beta_{i,0}\otimes\gamma_{i,0}
  &=\frac{s_i^xs_i^ys_i^z}{n^{3/2}}
      H_Az_i^A\otimes H_Bz_i^B\otimes H_Cz_i^C\\
  &=\frac1{n^{3/2}}
      H_Az_i^A\otimes H_Bz_i^B\otimes H_Cz_i^C\\
  &=\widetilde\zeta_i^A\otimes
    \widetilde\zeta_i^B\otimes
    \widetilde\zeta_i^C.
  \end{aligned}
\]
Summing proves (6). Repeating the calculation with the normalized formula
gives the factor \(r^{-3/2}\), which compared with (6) is exactly
\((n/r)^{3/2}\), proving (7). The target \(D_r\) is fixed by the realized
factor model and is not an initialization output, so neither balancing nor the
coefficient normalization acts on it. Thus the tangent range and the raw
coefficient target remain exactly the objects specified in Section~\ref{sec:setup}.


\end{proof}

\begin{proof}[Proof of the subsection conclusion]

Lemma~\ref{lem:step-002-balancing-scalars} derives the exact raw
radius-direction law and the balancing multipliers from
Assumption~\ref{assump:independent_initialization}. It proves that the zero
branch is null, supplies an exact scalar extension on that branch, and proves
that all balancing scalars are nonzero with
\(s_i^xs_i^ys_i^z=1\).

Conditional on \(\mathcal E_{\rm cond}\),
Lemma~\ref{lem:step-002-coefficient-gaussianization} applies compact SVDs to
the raw initialization and proves, with all factors displayed,
\[
  \bar\zeta_i^M=\frac{s_i^m}{\sqrt r}H_Mz_i^M,
  \qquad z_i^M\stackrel{\rm iid}{\sim}\mathcal N(0,I_r),
  \qquad \sigma(H_M)\subset[\kappa_1^{-1},\kappa_1].
\]
The independence is simultaneous across every mode and component, conditional
on the realized factors. The lemma explicitly records that the balancing
scalars may depend on the Gaussian arrays, which creates no gap because the
remaining conclusions are exact scalar identities.

Proposition~\ref{prop:step-002-balancing-invariance} then shows that every
individual tangent block is multiplied only by a nonzero pair scalar, including
the fixed \(1/n\) or \(1/r\) normalization factor. Hence every block range,
and therefore \(\mathscr S_0\), is unchanged. The same proposition uses the
triple-product identity to prove exact componentwise and summed preservation
of \(\widehat D_0\), while keeping \(D_r\) and the raw residual convention
untouched. Together these results give the exact
\(H_M,z_i^M,s_i^m\) representation and its invariance properties.
No Gram concentration, quotient-range statement, Haar factorization, or Haar
disintegration is asserted here.

\end{proof}

\subsection{Normalized Khatri--Rao Gram control}\label{app:step-003}

\begin{proof}[Auxiliary facts in the present notation]


\medskip\noindent\textbf{realized conditioning and balanced Gaussianization.}\par

We use the following statement in the present notation. Proposition~\ref{prop:step-001-realized-conditioning} and
Lemma~\ref{lem:step-002-coefficient-gaussianization} imply that, conditional
on any realized triple in \(\mathcal E_{\rm cond}\), the matrices \(H_M\)
are fixed and satisfy
\[
  \sigma(H_M)\subset[\kappa_1^{-1},\kappa_1],
\]
while the vectors \(z_i^M\) are independent \(\mathcal N(0,I_r)\) over all
modes and components and
\[
  \bar\zeta_i^M=\frac{s_i^m}{\sqrt r}H_Mz_i^M.
\]
Lemma~\ref{lem:step-002-balancing-scalars} additionally gives, on
the positive-radius branch,
\[
  s_i^x=\left(\frac{\chi_i^y\chi_i^z}{(\chi_i^x)^2}\right)^{1/3},\quad
  s_i^y=\left(\frac{\chi_i^x\chi_i^z}{(\chi_i^y)^2}\right)^{1/3},\quad
  s_i^z=\left(\frac{\chi_i^x\chi_i^y}{(\chi_i^z)^2}\right)^{1/3},
\]
where \(\chi_i^m=\|g_i^m\|_2\) and
\(g_i^m=\sqrt n\,\widetilde m_i\sim\mathcal N(0,I_n)\).

These are exactly the realized
elliptic maps, projected Gaussian arrays, and balancing multipliers used in
the three setting-defined normalized pair matrices. The proof below retains
the possible dependence of \(s_i^m\) on \(z_i^M\); it uses only a pointwise
diagonal bound on a separately generated full-radius event.


\medskip\noindent\textbf{Gaussian chi-square concentration.}\par

We use the following statement in the present notation. If \(g\sim\mathcal N(0,I_d)\), then, for
\(0<\delta<1\),
\[
  \mathbb P\!\left(\left|\frac{\|g\|_2^2}{d}-1\right|>\delta\right)
  \le 2\exp\!\left(-\frac{d\delta^2}{8}\right).
\]

Apply this with \(d=n\),
\(\delta=1/2\), and the \(3k\) standardized raw initialization columns.
Assumptions~\ref{assump:dimension}, \ref{assump:rank_window}, and
\ref{assump:independent_initialization} supply the dimension, number of
columns, and Gaussian law.


\medskip\noindent\textbf{Rectangular Gaussian operator-norm tail.}\par

We use the following statement in the present notation. If \(Z\in\mathbb R^{r\times k}\) has iid standard
Gaussian entries, then, for every \(u\ge0\),
\[
  \mathbb P\!\left(\|Z\|_{\rm op}>\sqrt r+\sqrt k+u\right)
  \le e^{-u^2/2}.
\]

In
Lemma~\ref{lem:step-003-isotropic-khatri-rao}, this is applied to one of the
two standard Gaussian mode arrays with \(u=8\sqrt{\log r}\). The lower rank
window \(k>r\) and \(r\ge2^{40}\) imply
\(u\le\sqrt r<\sqrt k\), so the resulting operator bound is
\(\|Z\|_{\rm op}\le3\sqrt k\).


\medskip\noindent\textbf{Self-adjoint matrix Bernstein inequality.}\par

We use the following statement in the present notation. Let \(X_1,\ldots,X_N\) be independent centered
self-adjoint \(d\times d\) random matrices satisfying
\(\|X_a\|_{\rm op}\le R\) almost surely, and put
\[
  v=\left\|\sum_{a=1}^N\mathbb E X_a^2\right\|_{\rm op}.
\]
Then, for every \(t\ge0\),
\[
  \mathbb P\!\left(
    \left\|\sum_{a=1}^N X_a\right\|_{\rm op}\ge t
  \right)
  \le 2d\exp\!\left(-\frac{t^2}{2v+(2Rt)/3}\right).
\]

Conditional on the first Gaussian
mode array and on an entrywise truncation event for the second array, the
row contributions
\[
  r^{-2}\bigl(D_a A D_a-\mu\operatorname{diag}(A)\bigr)
\]
are independent, centered, self-adjoint \(k\times k\) matrices. Their
almost-sure norm and variance proxy are computed explicitly in
Lemma~\ref{lem:step-003-isotropic-khatri-rao} before Bernstein is applied.


\medskip\noindent\textbf{Kronecker and product singular-value comparisons.}\par

We use the following statement in the present notation. For square invertible \(P,Q\),
\[
  \sigma_{\min}(P\otimes Q)
    =\sigma_{\min}(P)\sigma_{\min}(Q),\qquad
  \|P\otimes Q\|_{\rm op}=\|P\|_{\rm op}\|Q\|_{\rm op}.
\]
If \(K\) has full column rank and \(D\) is invertible diagonal, then
\[
  \sigma_{\min}((P\otimes Q)KD)
  \ge\sigma_{\min}(P\otimes Q)\sigma_{\min}(K)\sigma_{\min}(D),
\]
with the analogous upper bound for the operator norm.

Proposition~\ref{prop:step-003-elliptic-transfer}
uses \(P=H_M\), \(Q=H_N\), the isotropic Khatri--Rao matrix, and the exact
diagonal of pairwise balancing products. Full column rank follows from the
lower isotropic Gram bound.

\end{proof}



\begin{lemma}[full-radius control and dependent balancing diagonals]
\label{lem:step-003-radius-diagonals}
Under Assumptions~\ref{assump:dimension}, \ref{assump:rank_window}, and
\ref{assump:independent_initialization}, enlarge the dimension constant so
that \(C(\kappa,q)\ge1\), and suppose \(r\ge2^{40}\). Define the proof-local
event
\[
  \mathcal R_{\rm rad}
  :=\bigcap_{i=1}^k\bigcap_{m\in\{x,y,z\}}
    \left\{\frac12\le\frac{(\chi_i^m)^2}{n}\le\frac32\right\}.
\]
Then, conditional on any realized smoothing outcome,
\[
  \mathbb P(\mathcal R_{\rm rad}^c\mid A,B,C)\le r^{-28}.
\]
On \(\mathcal R_{\rm rad}\), for every component \(i\),
\[
  3^{-1/3}\le s_i^ys_i^z\le3^{1/3},\qquad
  3^{-1/3}\le s_i^xs_i^z\le3^{1/3},\qquad
  3^{-1/3}\le s_i^xs_i^y\le3^{1/3}.
\]

\end{lemma}

\begin{proof}
Initialization is independent of smoothing, so,
after conditioning on \((A,B,C)\), the \(3k\) vectors \(g_i^m\) remain
independent \(\mathcal N(0,I_n)\). The checked chi-square bound with
\(\delta=1/2\) and a union bound give
\[
  \mathbb P(\mathcal R_{\rm rad}^c\mid A,B,C)
  \le 6k e^{-n/32}
  \le 6r^{5/4}e^{-r^4\log r/32}.
\]
For \(r\ge2^{40}\), one has \(6\le r^{3/4}\) and \(r^4/32\ge30\), hence
\[
  6r^{5/4}e^{-r^4\log r/32}
  \le r^2r^{-30}=r^{-28}.                                  \tag{1}
\]
This bound improves as \(n\) increases, so it is uniform with no upper
restriction on \(n\).

On \(\mathcal R_{\rm rad}\), all radii are positive. The exact balancing formulas give
\[
  (s_i^ys_i^z)^3
    =\frac{(\chi_i^x)^2}{\chi_i^y\chi_i^z},\quad
  (s_i^xs_i^z)^3
    =\frac{(\chi_i^y)^2}{\chi_i^x\chi_i^z},\quad
  (s_i^xs_i^y)^3
    =\frac{(\chi_i^z)^2}{\chi_i^x\chi_i^y}.                \tag{2}
\]
Every squared radius lies in \([n/2,3n/2]\), while every product of two
radii lies in the same interval \([n/2,3n/2]\). Each ratio in (2) therefore
lies in \([1/3,3]\). Taking its positive cube root proves the three displayed
bounds. This is a pointwise consequence on \(\mathcal R_{\rm rad}\); it does
not use or assert independence between the pair multipliers and the projected
Gaussian coefficient vectors. 

\end{proof}


\begin{lemma}[rectangular isotropic Gaussian Khatri--Rao concentration]
\label{lem:step-003-isotropic-khatri-rao}
Under Assumption~\ref{assump:rank_window}, let
\(Z=[z_1\ \cdots\ z_k]\) and \(W=[w_1\ \cdots\ w_k]\) be independent
\(r\times k\) matrices with iid \(\mathcal N(0,1)\) entries. Define
\[
  K_0=\frac1r
    [z_1\otimes w_1\ \cdots\ z_k\otimes w_k]
    \in\mathbb R^{r^2\times k}.
\]
If \(r\ge2^{40}\), then
\[
  \mathbb P\!\left(
    \operatorname{spec}(K_0^{\mathsf T}K_0)
      \not\subset[1/2,3/2]
  \right)\le r^{-26}.                                      \tag{3}
\]

\end{lemma}

\begin{proof}
Put \(A=Z^{\mathsf T}Z\) and
\(L_r=8\sqrt{\log r}\). First define
\[
  \mathcal Z
  :=\left\{
    \max_i\left|\frac{\|z_i\|_2^2}{r}-1\right|\le\frac14,
    \quad \|Z\|_{\rm op}\le3\sqrt k
  \right\}.                                                \tag{4}
\]
The checked chi-square bound with \(d=r\), \(\delta=1/4\), followed by a
union over the \(k\) columns, gives a first failure term
\(2ke^{-r/128}\). The checked Gaussian operator tail with \(u=L_r\) gives
a second term \(r^{-32}\). Since \(r\ge2^{40}\) implies
\(L_r\le\sqrt r<\sqrt k\), its good-event threshold is at most
\(3\sqrt k\). The elementary numerical inequalities
\[
  2r^{5/4}e^{-r/128}\le r^{-28},
  \qquad r\ge2^{40},                                       \tag{5}
\]
and \(r^{-32}\le r^{-28}\) yield
\[
  \mathbb P(\mathcal Z^c)\le r^{-27}.                     \tag{6}
\]
On \(\mathcal Z\),
\[
  \|A\|_{\rm op}\le9k,\qquad
  a_*:=\max_i A_{ii}\le\frac54r,\qquad
  \frac34I_k\preceq\frac{\operatorname{diag}(A)}r
    \preceq\frac54I_k.                                    \tag{7}
\]

Write the rows of \(W\) as \(w_{a\bullet}^{\mathsf T}\),
\(a\in[r]\), and let
\(D_a=\operatorname{diag}(w_{a1},\ldots,w_{ak})\). The Hadamard/diagonal
identity gives
\[
  K_0^{\mathsf T}K_0
  =\frac1{r^2}\bigl(A\circ W^{\mathsf T}W\bigr)
  =\frac1{r^2}\sum_{a=1}^r D_aAD_a.                        \tag{8}
\]
Introduce the entrywise truncation event
\[
  \mathcal T_W:=\left\{\max_{a,i}|w_{ai}|\le L_r\right\}.
\]
A scalar Gaussian union bound gives
\[
  \mathbb P(\mathcal T_W^c)
  \le2rk e^{-L_r^2/2}
  \le2r^{9/4}r^{-32}
  \le r^{-29}.                                             \tag{9}
\]
Because \(\mathcal T_W\) is a product of coordinatewise truncation events,
conditional on \(\mathcal T_W\) all entries remain independent with the
common law of a standard normal \(\xi\) conditioned on
\(|\xi|\le L_r\). Put
\[
  \mu=\mathbb E\xi^2,\qquad \nu=\mathbb E\xi^4
\]
under this conditional law. Symmetry gives \(\mathbb E\xi=0\). Since
\(L_r\ge4\), integration of the Gaussian tail gives
\[
  \frac78\le\mu\le1,\qquad \nu\le6.                     \tag{10}
\]
Indeed,
\(\mathbb E[g^2\mathbf1_{\{|g|>L\}}]
 =2(L\varphi(L)+\overline\Phi(L))\), so the conditional second-moment loss
is below \(1/8\) for \(L\ge4\); and the fourth numerator is at most
\(\mathbb Eg^4=3\), while
\(\mathbb P(|g|\le L)\ge1/2\).

Fix any \(Z\in\mathcal Z\), put \(B=\operatorname{diag}(A)\), and, under
the conditional law given \(\mathcal T_W\), define the independent centered
self-adjoint matrices
\[
  X_a=\frac1{r^2}(D_aAD_a-\mu B),\qquad a\in[r].            \tag{11}
\]
The almost-sure truncation and (7) give
\[
  \|X_a\|_{\rm op}
  \le\frac{(L_r^2+1)\|A\|_{\rm op}}{r^2}
  \le\frac{18L_r^2k}{r^2}=:R_B.                            \tag{12}
\]
For completeness, independence, symmetry, and (10) make
\(\mathbb E(D_aAD_a)^2\) diagonal. Its \(i\)-th diagonal entry is
\[
  \mu^2(A^2)_{ii}+(\nu-\mu^2)A_{ii}^2.
\]
After subtracting the square of the mean, the \(i\)-th diagonal entry of
\(\mathbb E(D_aAD_a-\mu B)^2\) is at most
\[
  (A^2)_{ii}+6A_{ii}^2
  \le \|A\|_{\rm op}A_{ii}+6A_{ii}^2
  \le7\|A\|_{\rm op}a_*.
\]
Consequently the matrix-Bernstein variance proxy satisfies
\[
  v_B:=\left\|\sum_{a=1}^r\mathbb EX_a^2\right\|_{\rm op}
  \le\frac{7\|A\|_{\rm op}a_*}{r^3}
  \le\frac{315}{4}\frac{k}{r^2}
  \le79\frac{k}{r^2}.                                     \tag{13}
\]

Apply the checked self-adjoint matrix Bernstein inequality at \(t=1/8\).
Using \(L_r^2=64\log r\), (12)--(13), and
\(k\le r^{5/4}\), its exponent is at least
\[
  \frac{(1/8)^2}{2v_B+(2R_B/8)/3}
  \ge\frac{r^2}{64(158+96\log r)k}
  \ge\frac{r^{3/4}}{64(158+96\log r)}.                    \tag{14}
\]
For every \(r\ge2^{40}\),
\[
  \frac{r^{3/4}}{64(158+96\log r)}\ge30\log r.            \tag{15}
\]
The inequality holds directly at \(2^{40}\), and the ratio of its left side
to \(\log r\) is increasing thereafter. Thus, uniformly for
\(Z\in\mathcal Z\),
\[
  \mathbb P\!\left(
    \left\|K_0^{\mathsf T}K_0-\frac{\mu B}{r}\right\|_{\rm op}
      >\frac18
    \ \middle|\ Z,\mathcal T_W
  \right)
  \le2kr^{-30}\le r^{-28}.                                \tag{16}
\]

On the complementary event in (16), (7) and (10) imply
\[
  \lambda_{\min}(K_0^{\mathsf T}K_0)
  \ge\frac78\frac34-\frac18
  =\frac{17}{32}\ge\frac12,                               \tag{17}
\]
and
\[
  \lambda_{\max}(K_0^{\mathsf T}K_0)
  \le\frac54+\frac18
  =\frac{11}{8}\le\frac32.                                \tag{18}
\]
Finally, (6), (9), and (16), with no independence needed between their
failures, give
\[
  \mathbb P\!\left(
    \operatorname{spec}(K_0^{\mathsf T}K_0)
      \not\subset[1/2,3/2]
  \right)
  \le r^{-27}+r^{-29}+r^{-28}
  \le r^{-26}.
\]
This proves the lemma. 

\end{proof}


\begin{proposition}[elliptic transfer with dependent balancing scalars]
\label{prop:step-003-elliptic-transfer}
Under Proposition~\ref{prop:step-001-realized-conditioning},
Lemmas~\ref{lem:step-002-balancing-scalars} and
\ref{lem:step-002-coefficient-gaussianization}, and
Lemma~\ref{lem:step-003-radius-diagonals}, condition on any realized triple
in \(\mathcal E_{\rm cond}\). For the three mode pairs define
\[
  \begin{aligned}
  K_0^{BC}&=\frac1r[z_1^B\otimes z_1^C\ \cdots\ z_k^B\otimes z_k^C],
  &D^{BC}&=\operatorname{diag}(s_1^ys_1^z,\ldots,s_k^ys_k^z),\\
  K_0^{AC}&=\frac1r[z_1^A\otimes z_1^C\ \cdots\ z_k^A\otimes z_k^C],
  &D^{AC}&=\operatorname{diag}(s_1^xs_1^z,\ldots,s_k^xs_k^z),\\
  K_0^{AB}&=\frac1r[z_1^A\otimes z_1^B\ \cdots\ z_k^A\otimes z_k^B],
  &D^{AB}&=\operatorname{diag}(s_1^xs_1^y,\ldots,s_k^xs_k^y).
  \end{aligned}                                             \tag{19}
\]
If \(\mathcal R_{\rm rad}\) holds and each corresponding isotropic Gram in
(19) has spectrum in \([1/2,3/2]\), then each exact normalized balanced pair
Gram has spectrum in
\[
  \left[
    \frac1{2\,3^{2/3}\kappa_1^4},
    \frac{3^{5/3}}2\kappa_1^4
  \right].                                                  \tag{20}
\]

\end{proposition}

\begin{proof}
The normalized coefficient formula and the mixed-product identity for tensor
products give
\[
  \begin{aligned}
  \bar K_0^{\bar\beta\bar\gamma}
    &=(H_B\otimes H_C)K_0^{BC}D^{BC},\\
  \bar K_0^{\bar\alpha\bar\gamma}
    &=(H_A\otimes H_C)K_0^{AC}D^{AC},\\
  \bar K_0^{\bar\alpha\bar\beta}
    &=(H_A\otimes H_B)K_0^{AB}D^{AB}.
  \end{aligned}\tag{21}
\]
The factor \(1/r\) in (19) is exactly the product of the two
\(1/\sqrt r\) normalized coefficient factors. Thus (21) uses the setting's
normalized convention without a raw-scale substitution.

On \(\mathcal E_{\rm cond}\), every Kronecker elliptic map in (21) has
singular values in \([\kappa_1^{-2},\kappa_1^2]\). On
\(\mathcal R_{\rm rad}\), Lemma~\ref{lem:step-003-radius-diagonals} gives
\[
  \sigma(D^{MN})\subset[3^{-1/3},3^{1/3}]
\]
for all three pairs. The assumed isotropic Gram event gives
\[
  \sigma(K_0^{MN})\subset[2^{-1/2},(3/2)^{1/2}].
\]
Applying the checked product singular-value inequalities to (21) yields
\[
  \sigma_{\min}(\bar K_0^{MN})
    \ge \kappa_1^{-2}2^{-1/2}3^{-1/3},
  \qquad
  \|\bar K_0^{MN}\|_{\rm op}
    \le \kappa_1^2(3/2)^{1/2}3^{1/3}.                      \tag{22}
\]
Squaring (22) gives exactly (20).

This transfer is deterministic on the intersection of the stated events.
The diagonal \(D^{MN}\) can be arbitrarily dependent on \(K_0^{MN}\): the
pointwise singular-value comparison uses only its realized lower and upper
diagonal entries. Hence no false scalar/Gaussian independence is introduced.
The endpoint singular values of the \(H_M\) and endpoint pair multipliers are
included in (22). 

\end{proof}


\begin{proposition}[simultaneous normalized pair-Gram event]
\label{prop:step-003-normalized-gram-event}
Under Assumptions~\ref{assump:dimension}, \ref{assump:rank_window}, and
\ref{assump:independent_initialization}, Proposition~\ref{prop:step-001-realized-conditioning}, and Lemmas~\ref{lem:step-002-balancing-scalars} and
\ref{lem:step-002-coefficient-gaussianization}, define
\[
  r_{0,{\rm gram}}(\kappa)
  :=\max\left\{
    2^{40},
    \left\lceil
      \bigl(2\,3^{2/3}\kappa_1^4\bigr)^{1/20}
    \right\rceil
  \right\},
  \qquad \kappa_1=2\kappa^2.                              \tag{23}
\]
Choose \(C(\kappa,q)\ge1\). Then, for every
\(r\ge r_{0,{\rm gram}}(\kappa)\), every
\(n\ge C(\kappa,q)r^4\log r\), every
\(r<k\le\lfloor r^{5/4}\rfloor\), and every realized factor triple in
\(\mathcal E_{\rm cond}\),
\[
  \mathbb P\!\left(
    \mathcal E_{\rm gram}^{\rm norm}
    \mid A,B,C
  \right)\ge1-r^{-20}.                                     \tag{24}
\]

\end{proposition}

\begin{proof}
Conditional on the realized factors, the three
mode arrays \(Z_M=[z_1^M\ \cdots\ z_k^M]\) are independent standard
Gaussian arrays by the Gaussianization lemma. Therefore
Lemma~\ref{lem:step-003-isotropic-khatri-rao} applies to each of the three
pairs in (19), giving failure at most \(r^{-26}\) per pair. The pairs share
mode arrays, but a union bound needs no independence. Together with
Lemma~\ref{lem:step-003-radius-diagonals},
\[
  \mathbb P\!\left(
    \mathcal R_{\rm rad}^c
    \ \text{or some isotropic pair event fails}
    \mid A,B,C
  \right)
  \le r^{-28}+3r^{-26}
  \le4r^{-26}
  \le r^{-20}.                                             \tag{25}
\]

On the complementary event, Proposition~\ref{prop:step-003-elliptic-transfer}
places every normalized pair-Gram eigenvalue in the constant interval (20).
The threshold (23) gives
\[
  r^{20}\ge2\,3^{2/3}\kappa_1^4.                          \tag{26}
\]
Consequently
\[
  \frac1{2\,3^{2/3}\kappa_1^4}\ge r^{-20},               \tag{27}
\]
and, because
\[
  \frac{3^{5/3}}2\kappa_1^4
  \le2\,3^{2/3}\kappa_1^4,
\]
also
\[
  \frac{3^{5/3}}2\kappa_1^4\le r^{20}.                   \tag{28}
\]
Thus all three exact normalized Grams lie in the setting-defined polynomial
window on an event whose total conditional failure is bounded by (25). This
is precisely \(\mathcal E_{\rm gram}^{\rm norm}\), proving (24).

The proof is uniform in the realized factor triple because it uses that triple
only through the common \(\kappa_1\) singular interval. It is uniform in
\(n\) above its lower threshold because the projected \(z_i^M\) law is exactly
standard Gaussian for every \(n\), while the full-radius failure in (1)
decreases with \(n\). The cases \(k=r+1\) and
\(k=\lfloor r^{5/4}\rfloor\) are both included, as are the endpoint elliptic
and radial bounds. 

\end{proof}

\begin{proof}[Proof of the subsection conclusion]

Lemma~\ref{lem:step-003-radius-diagonals} derives a full-radius
event from the primitive ambient Gaussian initialization and proves uniform
pointwise bounds for all three pairwise balancing diagonals. This is the only
place where full radii enter, and their dependence on the projected
coefficient Gaussians is retained rather than discarded.

Lemma~\ref{lem:step-003-isotropic-khatri-rao} proves both spectral sides for
an \(r^2\times k\) isotropic Gaussian Khatri--Rao matrix. Its row
decomposition (8) makes the conditional mean diagonal; entrywise Gaussian
truncation supplies an almost-sure Bernstein envelope; (12)--(13) verify the
matrix-Bernstein hypotheses; and (14)--(18) give the explicit lower and upper
window with failure at most \(r^{-26}\).

Proposition~\ref{prop:step-003-elliptic-transfer} then uses the exact
preceding formula
\[
  \bar K^{MN}=(H_M\otimes H_N)K_0^{MN}D^{MN}
\]
with the setting's \(1/r\) normalized scale. Pointwise singular-value
comparisons transfer the isotropic window through the elliptic maps and the
possibly dependent balancing diagonal, producing the constant interval (20)
for each exact normalized Gram.

Finally, Proposition~\ref{prop:step-003-normalized-gram-event} unions the
single radial event and the three pair events, converts their total failure to
at most \(r^{-20}\), and verifies by (26)--(28) that the constant interval is
contained in \([r^{-20},r^{20}]\). These named results prove exactly all three
normalized pair-Gram windows and nothing about the raw-scale bridge or
\(\mathcal E_{\rm size}\).

\end{proof}

\subsection{Balanced initialization radius}\label{app:step-004}

\begin{proof}[Auxiliary facts in the present notation]


\medskip\noindent\textbf{exact balancing law.}\par

Lemma~\ref{lem:step-002-balancing-scalars} states that for each component
\(i\), if the three raw norms are positive, the product-preserving balancing
map makes all three output norms equal to their geometric mean. If any raw
factor is zero, the map leaves the whole raw triple unchanged.

The raw columns here are exactly
\(\widetilde x_i,\widetilde y_i,\widetilde z_i\) from the setting. The lemma
was under Assumption~\ref{assump:independent_initialization}; no
coefficient, Gram, or factor-conditioning output from Proposition~\ref{prop:step-002-balancing-invariance} is used.
Proposition~\ref{prop:step-004-balanced-size-transfer} applies this law only
after the simultaneous raw-radius event has been proved.


\medskip\noindent\textbf{Gaussian-square moment generating function.}\par

We use the following statement in the present notation. If \(g\sim\mathcal N(0,I_n)\), then for every
\(0<\lambda<1/2\),
\[
  \mathbb E e^{\lambda\|g\|_2^2}=(1-2\lambda)^{-n/2}.
\]
Consequently, Chernoff's inequality gives
\[
  \mathbb P(\|g\|_2^2\ge 4n)
  \le \exp\!\left[-\frac{3-\log 4}{2}n\right]
  \le e^{-n/2}.
\]

For each raw column set
\(g_i^m=\sqrt n\,\widetilde m_i\). Assumption~\ref{assump:independent_initialization}
gives \(g_i^m\sim\mathcal N(0,I_n)\), even after conditioning on the
independent smoothing. Lemma~\ref{lem:step-004-raw-radius-tail} derives the
displayed moment generating function and tail directly, so no external
citation is used as proof authority.

\end{proof}



\begin{lemma}[raw Gaussian radius tail]
\label{lem:step-004-raw-radius-tail}
Under Assumption~\ref{assump:independent_initialization}, fix any realized
factor triple \((A,B,C)\) and any component-mode pair
\((i,m)\in[k]\times\{x,y,z\}\). Then
\[
  \mathbb P\!\left(\|\widetilde m_i\|_2>2\mid A,B,C\right)
  \le \exp\!\left[-\frac{3-\log4}{2}n\right]
  \le e^{-n/2}.
\]

\end{lemma}

\begin{proof}
Independence of initialization from smoothing
implies that, conditionally on the realized factors,
\(g_i^m:=\sqrt n\,\widetilde m_i\) remains a standard Gaussian vector in
\(\mathbb R^n\). For one standard normal scalar \(G\), direct Gaussian
integration gives
\[
  \mathbb E e^{\lambda G^2}=(1-2\lambda)^{-1/2},
  \qquad 0<\lambda<\frac12.
\]
Independence of the \(n\) coordinates therefore yields
\[
  \mathbb E e^{\lambda\|g_i^m\|_2^2}
  =(1-2\lambda)^{-n/2}.
\]
Markov's inequality with \(\lambda=3/8\) gives
\[
  \begin{aligned}
  \mathbb P\!\left(\|\widetilde m_i\|_2>2\mid A,B,C\right)
  &=\mathbb P\!\left(\|g_i^m\|_2^2>4n\mid A,B,C\right)\\
  &\le e^{-(3/8)4n}(1-3/4)^{-n/2}\\
  &=e^{-3n/2}4^{n/2}
   =\exp\!\left[-\frac{3-\log4}{2}n\right].
  \end{aligned}
\]
Since \(\log4<2\), the exponent coefficient
\((3-\log4)/2\) is larger than \(1/2\), proving the second inequality. The
calculation is uniform in the realized factors because their values do not
enter the conditional Gaussian law. 

\end{proof}


\begin{proposition}[simultaneous control of all raw initialization radii]
\label{prop:step-004-uniform-raw-size}
Under Assumptions~\ref{assump:dimension}, \ref{assump:rank_window}, and
\ref{assump:independent_initialization}, and
Lemma~\ref{lem:step-004-raw-radius-tail}, define the event
\[
  \mathcal E_{\rm raw,size}
  :=\left\{\max_{i\in[k],\,m\in\{x,y,z\}}
       \|\widetilde m_i\|_2\le2\right\}.
\]
Choose the dimension-constant contribution
\(C_{\rm size}:=1\) and the rank threshold \(r_{\rm size}:=3\). If
\[
  C(\kappa,q)\ge C_{\rm size},\qquad r\ge r_{\rm size},
\]
then, conditionally on every realized factor triple,
\[
  \mathbb P\!\left(\mathcal E_{\rm raw,size}^{\mathsf c}
    \mid A,B,C\right)\le r^{-20}.
\]
This conclusion holds throughout the full rank window, including the boundary
\(k=\lfloor r^{5/4}\rfloor\).

\end{proposition}

\begin{proof}
A union bound over the \(3k\) raw columns and
Lemma~\ref{lem:step-004-raw-radius-tail} give
\[
  \mathbb P\!\left(\mathcal E_{\rm raw,size}^{\mathsf c}
      \mid A,B,C\right)
  \le 3k e^{-n/2}.
\]
No independence is needed for this union bound, although the columns are in
fact independent. Assumptions~\ref{assump:dimension} and
\ref{assump:rank_window}, together with \(C(\kappa,q)\ge1\), imply
\[
  3k e^{-n/2}
  \le 3r^{5/4}
       \exp\!\left(-\frac12r^4\log r\right).
\]
For \(r\ge3\),
\[
  \frac{r^4}{2}
  \ge \frac{81}{2}
  \ge \frac{85}{4}+\frac{\log3}{\log r},
\]
because \(\log3/\log r\le1\). Multiplying by \(\log r>0\) and rearranging
gives the explicit absorption inequality
\[
  \log3+\frac54\log r-\frac{r^4}{2}\log r
  \le -20\log r.
\]
Exponentiation proves
\[
  3r^{5/4}\exp\!\left(-\frac12r^4\log r\right)
  \le r^{-20}.
\]
The use of \(k\le r^{5/4}\) is valid at
\(k=\lfloor r^{5/4}\rfloor\), so the maximal-rank boundary requires no
separate loss. The lower inequality \(r<k\) is part of the declared window
but is not needed for this upper-tail estimate. 

\end{proof}


\begin{proposition}[geometric-mean transfer to balanced columns]
\label{prop:step-004-balanced-size-transfer}
\leavevmode\par\noindent
Under Assumption~\ref{assump:independent_initialization},
Lemma~\ref{lem:step-002-balancing-scalars}, and
Proposition~\ref{prop:step-004-uniform-raw-size}, the deterministic inclusion
\[
  \mathcal E_{\rm raw,size}\subseteq\mathcal E_{\rm size}
\]
holds. Hence, whenever \(C(\kappa,q)\ge1\) and \(r\ge3\), conditionally on
every realized factor triple,
\[
  \mathbb P\!\left(\mathcal E_{\rm size}^{\mathsf c}
    \mid A,B,C\right)\le r^{-20}.
\]

\end{proposition}

\begin{proof}
Fix a component \(i\) on
\(\mathcal E_{\rm raw,size}\). First suppose all three raw norms are
positive. The balancing lemma gives
\[
  \|x_{i,0}\|_2=\|y_{i,0}\|_2=\|z_{i,0}\|_2
  =\left(\|\widetilde x_i\|_2
          \|\widetilde y_i\|_2
          \|\widetilde z_i\|_2\right)^{1/3}.
\]
Every raw norm on the right is at most \(2\), so
\[
  \left(\|\widetilde x_i\|_2
          \|\widetilde y_i\|_2
          \|\widetilde z_i\|_2\right)^{1/3}
  \le (2\cdot2\cdot2)^{1/3}=2.
\]

If at least one raw factor is exactly zero, the setting-defined zero branch
leaves the entire raw triple unchanged. On
\(\mathcal E_{\rm raw,size}\), each unchanged raw norm is already at most
\(2\), including any nonzero factors in that triple. Thus the same balanced
size conclusion holds on the zero branch without division by a raw norm and
without removing a null event.

The argument applies to every \(i\in[k]\), proving the event inclusion.
Consequently
\[
  \mathbb P\!\left(\mathcal E_{\rm size}^{\mathsf c}\mid A,B,C\right)
  \le
  \mathbb P\!\left(\mathcal E_{\rm raw,size}^{\mathsf c}\mid A,B,C\right)
  \le r^{-20}
\]
by Proposition~\ref{prop:step-004-uniform-raw-size}. 

\end{proof}

\begin{proof}[Proof of the subsection conclusion]

Lemma~\ref{lem:step-004-raw-radius-tail} derives from
Assumption~\ref{assump:independent_initialization} the explicit conditional
one-column bound \(e^{-n/2}\). Proposition~\ref{prop:step-004-uniform-raw-size}
unions that bound over exactly the \(3k\) raw initialization columns and uses
the full admissible boundary \(k\le\lfloor r^{5/4}\rfloor\), the explicit
dimension contribution \(C_{\rm size}=1\), and the explicit threshold
\(r_{\rm size}=3\) to obtain
\[
  \mathbb P\!\left(\mathcal E_{\rm raw,size}^{\mathsf c}
    \mid A,B,C\right)\le r^{-20}.
\]

Proposition~\ref{prop:step-004-balanced-size-transfer}, using the exact
balancing law from Proposition~\ref{prop:step-002-balancing-invariance}, proves on the positive-radius branch
that each balanced norm is the geometric mean of three raw norms bounded by
\(2\), and separately proves the same conclusion on the setting-defined
zero-factor no-op branch. Therefore
\(\mathcal E_{\rm raw,size}\subseteq\mathcal E_{\rm size}\), and hence
\[
  \boxed{
  \mathbb P\!\left(\mathcal E_{\rm size}^{\mathsf c}
    \mid A,B,C\right)\le r^{-20}}
\]
uniformly for every realized factor triple in \(\mathcal E_{\rm cond}\),
every \(r\ge3\), every \(n\ge C(\kappa,q)r^4\log r\) with
\(C(\kappa,q)\ge1\), and every
\(r<k\le\lfloor r^{5/4}\rfloor\). This proves the claimed conditional
failure bound for the balanced initial-size event. No Gram or path statement
has been used or proved.

\end{proof}

\subsection{Raw and normalized tangent geometry}\label{app:step-005}

\begin{proof}[Auxiliary facts in the present notation]


\medskip\noindent\textbf{Balanced coefficient representation and raw invariance.}\par

Lemma~\ref{lem:step-002-coefficient-gaussianization} and
Proposition~\ref{prop:step-002-balancing-invariance} state, conditional on a
realized triple in \(\mathcal E_{\rm cond}\), that the
raw balanced coefficient vectors
\(\alpha_{i,0},\beta_{i,0},\gamma_{i,0}\) are well defined and that the
setting normalization is exactly
\[
  \bar\alpha_{i,0}=\sqrt{\frac nr}\,\alpha_{i,0},\qquad
  \bar\beta_{i,0}=\sqrt{\frac nr}\,\beta_{i,0},\qquad
  \bar\gamma_{i,0}=\sqrt{\frac nr}\,\gamma_{i,0}.
\]
Lemma~\ref{lem:step-002-coefficient-gaussianization} proves nonvanishing on
the regular branch and handles the exceptional zero-vector branch by an exact
scalar extension. Proposition~\ref{prop:step-002-balancing-invariance}
preserves
\(\widehat D_0=\sum_i\alpha_{i,0}\otimes\beta_{i,0}\otimes\gamma_{i,0}\),
and does not normalize \(D_r\).

The results below use exactly these raw and normalized coefficient vectors.
The regular-branch nonvanishing is used only
to identify each local kernel as exactly two-dimensional. All range and
quotient conclusions are also proved directly on exceptional branches, so no
branch is removed from the algebra. The synthesis-map equality below is
derived directly, independently of the preceding blockwise span identity.

\end{proof}



\begin{lemma}[Exact scaling of all three pair Grams]
\label{lem:step-005-pair-gram-scaling}
Under the setting definitions, Lemma~\ref{lem:step-002-coefficient-gaussianization}, and Proposition~\ref{prop:step-002-balancing-invariance}, define
\[
\begin{aligned}
  K_{\rm raw}^{\beta\gamma}
    &=[\,\beta_{1,0}\otimes\gamma_{1,0}\ \cdots\
          \beta_{k,0}\otimes\gamma_{k,0}\,],\\
  K_{\rm raw}^{\alpha\gamma}
    &=[\,\alpha_{1,0}\otimes\gamma_{1,0}\ \cdots\
          \alpha_{k,0}\otimes\gamma_{k,0}\,],\\
  K_{\rm raw}^{\alpha\beta}
    &=[\,\alpha_{1,0}\otimes\beta_{1,0}\ \cdots\
          \alpha_{k,0}\otimes\beta_{k,0}\,],
\end{aligned}
\]
and define the normalized matrices by replacing every coefficient vector by
its barred version. For every pair
\(pq\in\{\beta\gamma,\alpha\gamma,\alpha\beta\}\), let
\(G_{\rm raw}^{pq}=(K_{\rm raw}^{pq})^{\mathsf T}K_{\rm raw}^{pq}\) and
\(G_{\rm norm}^{pq}=(K_{\rm norm}^{pq})^{\mathsf T}K_{\rm norm}^{pq}\).
Then, pointwise on every branch,
\[
  K_{\rm raw}^{pq}=\frac rn K_{\rm norm}^{pq},
  \qquad
  \boxed{G_{\rm raw}^{pq}=\left(\frac rn\right)^2G_{\rm norm}^{pq}}
  \quad
  \bigl(pq=\beta\gamma,\alpha\gamma,\alpha\beta\bigr).
\]

\end{lemma}

\begin{proof}
Put \(c=\sqrt{n/r}>0\). The setting definitions
give \(\bar\alpha_{i,0}=c\alpha_{i,0}\),
\(\bar\beta_{i,0}=c\beta_{i,0}\), and
\(\bar\gamma_{i,0}=c\gamma_{i,0}\). Thus, for example,
\[
  \beta_{i,0}\otimes\gamma_{i,0}
  =c^{-2}(\bar\beta_{i,0}\otimes\bar\gamma_{i,0})
  =\frac rn(\bar\beta_{i,0}\otimes\bar\gamma_{i,0}).
\]
The identical calculation applies to the other two pairs, proving the matrix
identity column by column. Taking transpose times the matrix gives
\[
  (K_{\rm raw}^{pq})^{\mathsf T}K_{\rm raw}^{pq}
  =\left(\frac rn\right)^2
    (K_{\rm norm}^{pq})^{\mathsf T}K_{\rm norm}^{pq}.
\]
No division by a coefficient vector occurs, so the calculation remains valid
when a pair column is zero. This proves the lemma. 

\end{proof}


\begin{proposition}[Exact scaling and range of the tangent synthesis maps]
\label{prop:step-005-synthesis-scaling}
Under the setting definitions, Lemma~\ref{lem:step-002-coefficient-gaussianization}, and Proposition~\ref{prop:step-002-balancing-invariance}, let
\[
  \mathcal P=\bigoplus_{i=1}^k
    (\mathbb R^r\oplus\mathbb R^r\oplus\mathbb R^r)
\]
and write \(p=((u_i,v_i,w_i))_{i=1}^k\in\mathcal P\). Define the raw and
normalized tangent synthesis maps into
\(\mathcal H=(\mathbb R^r)^{\otimes3}\) by
\[
\begin{aligned}
  \mathcal T_{\rm raw}p
  &=\sum_{i=1}^k\bigl(
      u_i\otimes\beta_{i,0}\otimes\gamma_{i,0}
      +\alpha_{i,0}\otimes v_i\otimes\gamma_{i,0}
      +\alpha_{i,0}\otimes\beta_{i,0}\otimes w_i\bigr),\\
  \mathcal T_{\rm norm}p
  &=\sum_{i=1}^k\bigl(
      u_i\otimes\bar\beta_{i,0}\otimes\bar\gamma_{i,0}
      +\bar\alpha_{i,0}\otimes v_i\otimes\bar\gamma_{i,0}
      +\bar\alpha_{i,0}\otimes\bar\beta_{i,0}\otimes w_i\bigr).
\end{aligned}
\]
Then, pointwise on every branch,
\[
  \boxed{\mathcal T_{\rm raw}=\frac rn\mathcal T_{\rm norm}},
  \qquad
  \operatorname{range}(\mathcal T_{\rm raw})
    =\mathscr S_0^{\rm raw},
  \qquad
  \operatorname{range}(\mathcal T_{\rm norm})
    =\mathscr S_0^{\rm norm}.
\]
Consequently,
\[
  \boxed{\mathscr S_0^{\rm raw}=\mathscr S_0^{\rm norm}=\mathscr S_0},
  \qquad
  \ker(\mathcal T_{\rm raw})=\ker(\mathcal T_{\rm norm}).
\]

\end{proposition}

\begin{proof}
With \(c=\sqrt{n/r}\), every normalized tangent
summand contains exactly two barred coefficient vectors and therefore equals
\(c^2=n/r\) times the corresponding raw summand. Summing over components
gives \(\mathcal T_{\rm norm}=(n/r)\mathcal T_{\rm raw}\), equivalently the
boxed map identity. The scalar \(r/n\) is nonzero.

Every image of \(\mathcal T_{\rm raw}\) is a finite sum of generators in the
definition of \(\mathscr S_0^{\rm raw}\), so its range is contained in that
span. Conversely, each displayed raw generator is obtained by setting exactly
one of the variables \(u_i,v_i,w_i\) and setting all other variables to zero.
Hence the range is the full displayed span. The normalized case is identical.
Multiplication of a linear map by a nonzero scalar changes neither its range
nor its kernel, proving both final equalities. Zero coefficient vectors merely
remove zero generators from both maps and do not affect any argument. This
proves the proposition. 

\end{proof}


\begin{lemma}[Exact component gauge kernels and canonical gauge images]
\label{lem:step-005-gauge-kernel}
Under the setting definitions, Lemma~\ref{lem:step-002-coefficient-gaussianization}, and Proposition~\ref{prop:step-002-balancing-invariance}, let
\[
  a_i=\alpha_{i,0},\qquad b_i=\beta_{i,0},\qquad c_i=\gamma_{i,0},
\]
and let \(\mathcal T_{{\rm raw},i}\) denote the \(i\)-th summand map. Define
the canonical gauge map and its image by
\[
\begin{aligned}
  \Gamma_i(\lambda,\mu)
  &=\lambda(a_i,-b_i,0)+\mu(a_i,0,-c_i),\\
  \mathfrak n_i&=\operatorname{image}(\Gamma_i).
\end{aligned}
\]
Then \(\mathfrak n_i\subseteq\ker(\mathcal T_{{\rm raw},i})\) on every
branch. On the almost-sure branch supplied by
Lemma~\ref{lem:step-002-coefficient-gaussianization}, where
\(a_i,b_i,c_i\ne0\),
\[
  \ker(\mathcal T_{{\rm raw},i})
  =\{(\lambda a_i,\mu b_i,\nu c_i):\lambda+\mu+\nu=0\}
  =\mathfrak n_i,\qquad \dim\mathfrak n_i=2.
\]
On all branches the exact classification is:

\begin{center}
\small
\begin{adjustbox}{max width=\textwidth}
\begin{tabular}{lllll}
\hline
Branch & Local synthesis map & Local kernel & Kernel dimension & \(\dim\mathfrak n_i\) \\
\hline
\(a_i,b_i,c_i\ne0\) & Full three-term map & \(\{(\lambda a_i,\mu b_i,\nu c_i):\lambda+\mu+\nu=0\}\) & \(2\) & \(2\) \\
\(a_i=0,\ b_i,c_i\ne0\) & \(u\otimes b_i\otimes c_i\) & \(\{0\}\oplus\mathbb R^r\oplus\mathbb R^r\) & \(2r\) & \(2\) \\
\(b_i=0,\ a_i,c_i\ne0\) & \(a_i\otimes v\otimes c_i\) & \(\mathbb R^r\oplus\{0\}\oplus\mathbb R^r\) & \(2r\) & \(2\) \\
\(c_i=0,\ a_i,b_i\ne0\) & \(a_i\otimes b_i\otimes w\) & \(\mathbb R^r\oplus\mathbb R^r\oplus\{0\}\) & \(2r\) & \(2\) \\
\(a_i\ne0,\ b_i=c_i=0\) & \(0\) & \((\mathbb R^r)^3\) & \(3r\) & \(1\) \\
\(b_i\ne0,\ a_i=c_i=0\) & \(0\) & \((\mathbb R^r)^3\) & \(3r\) & \(1\) \\
\(c_i\ne0,\ a_i=b_i=0\) & \(0\) & \((\mathbb R^r)^3\) & \(3r\) & \(1\) \\
\(a_i=b_i=c_i=0\) & \(0\) & \((\mathbb R^r)^3\) & \(3r\) & \(0\) \\
\hline
\end{tabular}
\end{adjustbox}
\end{center}

The normalized gauge map is
\(\Gamma_i^{\rm norm}=\sqrt{n/r}\,\Gamma_i\), so its image is the same
\(\mathfrak n_i\), and the normalized and raw local kernels agree.

\end{lemma}

\begin{proof}
Direct substitution gives
\[
\begin{aligned}
  \mathcal T_{{\rm raw},i}(a_i,-b_i,0)
  &=a_i\otimes b_i\otimes c_i-a_i\otimes b_i\otimes c_i=0,\\
  \mathcal T_{{\rm raw},i}(a_i,0,-c_i)
  &=a_i\otimes b_i\otimes c_i-a_i\otimes b_i\otimes c_i=0.
\end{aligned}
\]
This proves the all-branch inclusion without a nonvanishing assumption.

Now suppose \(a_i,b_i,c_i\ne0\) and
\(\mathcal T_{{\rm raw},i}(u,v,w)=0\). Choose linear functionals
\(b_i^*,c_i^*\) with \(b_i^*(b_i)=c_i^*(c_i)=1\), and let
\(P_{a_i^\perp}\) be orthogonal projection onto \(a_i^\perp\). Applying
\(P_{a_i^\perp}\otimes b_i^*\otimes c_i^*\) gives
\(u=\lambda a_i\). Repeating in the other modes gives
\(v=\mu b_i\) and \(w=\nu c_i\). Substitution yields
\[
  (\lambda+\mu+\nu)a_i\otimes b_i\otimes c_i=0.
\]
The rank-one tensor is nonzero, so \(\lambda+\mu+\nu=0\). Conversely, every
such triple is killed by the map, and
\[
  (\lambda a_i,\mu b_i,\nu c_i)
  =-\mu(a_i,-b_i,0)-\nu(a_i,0,-c_i)
\]
proves equality with \(\mathfrak n_i\). The two
generators are independent because all three coefficient vectors are nonzero.

If exactly one fixed vector is zero, the table's displayed one-term map
survives and vanishes exactly when its free vector is zero; the two columns of
\(\Gamma_i\) remain independent. If exactly two fixed vectors are zero, the
local map is zero and the gauge image is the one-dimensional span of the
surviving block vector. If all three vanish, both maps vanish. These
observations prove the remaining table rows. The normalized conclusions follow
from the common nonzero normalization scalar and
Proposition~\ref{prop:step-005-synthesis-scaling}. 

\end{proof}


\begin{proposition}[Gauge quotient, unreduced range, and raw-target compatibility]
\label{prop:step-005-quotient-range}
Under the setting definitions, Lemma~\ref{lem:step-002-coefficient-gaussianization}, Proposition~\ref{prop:step-002-balancing-invariance},
Proposition~\ref{prop:step-005-synthesis-scaling}, and
Lemma~\ref{lem:step-005-gauge-kernel}, embed each \(\mathfrak n_i\) in the
\(i\)-th parameter block and define
\[
  \mathfrak N_{\rm can}=\bigoplus_{i=1}^k\mathfrak n_i\subseteq\mathcal P,
  \qquad \mathcal Q_{\rm can}=\mathcal P/\mathfrak N_{\rm can}.
\]
Then both synthesis maps factor uniquely through the quotient map
\(\pi_{\rm can}:\mathcal P\to\mathcal Q_{\rm can}\):
\[
  \widetilde{\mathcal T}_{\rm raw}([p])=\mathcal T_{\rm raw}p,
  \qquad
  \widetilde{\mathcal T}_{\rm norm}([p])=\mathcal T_{\rm norm}p.
\]
These induced maps satisfy
\[
\begin{aligned}
  \operatorname{range}(\widetilde{\mathcal T}_{\rm raw})
  &=\operatorname{range}(\mathcal T_{\rm raw})
   =\mathscr S_0,\\
  \operatorname{range}(\widetilde{\mathcal T}_{\rm norm})
  &=\operatorname{range}(\mathcal T_{\rm norm})
   =\mathscr S_0,\\
  \widetilde{\mathcal T}_{\rm raw}
  &=\frac rn\widetilde{\mathcal T}_{\rm norm}.
\end{aligned}
\]
On the almost-sure all-nonzero branch,
\(\dim\mathfrak N_{\rm can}=2k\), so this is exactly the two-gauge-directions-per-component quotient. On every zero branch it is the
quotient by the branch-dependent, possibly lower-dimensional canonical gauge
image. If
\(\mathcal P=\mathfrak N_{\rm can}\oplus\mathcal C_{\rm can}\), call a basis
of \(\mathcal C_{\rm can}\) a reduced parameter frame. Its synthesized
columns have image range \(\mathscr S_0\), but no injectivity or linear
independence is claimed. Furthermore,
\[
  \boxed{\widehat D_0\in\mathscr S_0}.
\]
All these are raw coefficient-space statements: \(D_r\) is unchanged,
\(D_r-\widehat D_0\) is the raw coefficient residual in the Frobenius inner
product on \(\mathcal H\), and the physical loss remains the setting's
\(\|T-S(X,Y,Z)\|_F^2\) with no normalization.

\end{proposition}

\begin{proof}
Lemma~\ref{lem:step-005-gauge-kernel} gives
\(\mathfrak n_i\subseteq\ker(\mathcal T_{{\rm raw},i})\) for each component.
Because the global map is the sum of its local maps,
\(\mathfrak N_{\rm can}\subseteq\ker(\mathcal T_{\rm raw})\).
Proposition~\ref{prop:step-005-synthesis-scaling} then gives
\(\mathfrak N_{\rm can}\subseteq\ker(\mathcal T_{\rm norm})\). Hence
the two formulas on cosets are well defined: replacing \(p\) by \(p+g\),
\(g\in\mathfrak N_{\rm can}\), does not change either image.

The quotient map \(\pi_{\rm can}\) is surjective, so every quotient image is an
unreduced image and every unreduced image is the image of the coset of its
parameter. Therefore
\[
  \operatorname{range}(\widetilde{\mathcal T}_{\rm raw})
  =\operatorname{range}(\mathcal T_{\rm raw}),
  \qquad
  \operatorname{range}(\widetilde{\mathcal T}_{\rm norm})
  =\operatorname{range}(\mathcal T_{\rm norm}).
\]
The remaining range and scaling identities follow from
Proposition~\ref{prop:step-005-synthesis-scaling}. If
\(\mathcal P=\mathfrak N_{\rm can}\oplus\mathcal C_{\rm can}\), then for
\(p=g+c\),
\(\mathcal T_{\rm raw}p=\mathcal T_{\rm raw}c\); hence restriction to any
such complement, and thus any reduced parameter frame, has the same image
range as the unreduced generator frame. Additional kernel vectors may remain
in \(\mathcal C_{\rm can}\), so no basis or injectivity claim is made. An
actual tangent-space basis would instead use the full common kernel
\(\ker(\mathcal T_{\rm raw})=\ker(\mathcal T_{\rm norm})\).

For membership of the represented raw coefficient tensor, take
\(p_{\widehat D}=((\alpha_{i,0},0,0))_{i=1}^k\). Then
\[
  \mathcal T_{\rm raw}p_{\widehat D}
  =\sum_{i=1}^k
    \alpha_{i,0}\otimes\beta_{i,0}\otimes\gamma_{i,0}
  =\widehat D_0,
\]
so \(\widehat D_0\in\operatorname{range}(\mathcal T_{\rm raw})=\mathscr S_0\),
including on every zero branch. Neither the quotient nor the coefficient
normalization acts on \(D_r\), and no map here acts on the physical
tensor space. This proves every raw-target and metric convention in the
statement. 

The global kernel of \(\mathcal T_{\rm raw}\) can contain additional
cross-component linear dependencies. No equality
\(\ker(\mathcal T_{\rm raw})=\mathfrak N_{\rm can}\) is claimed or needed.
On the almost-sure nonzero branch the canonical image removes exactly two
local gauge directions per component; on zero branches its dimension is the
one in Lemma~\ref{lem:step-005-gauge-kernel}. In every case, inclusion in the
full kernel is the precise condition that preserves the range.

\end{proof}

\begin{proof}[Proof of the subsection conclusion]

Lemma~\ref{lem:step-005-pair-gram-scaling} applies the setting normalization
column by column to all three Khatri--Rao pair matrices and proves exactly
\[
  G_{\rm raw}^{pq}=\left(\frac rn\right)^2G_{\rm norm}^{pq},
  \qquad pq\in\{\beta\gamma,\alpha\gamma,\alpha\beta\}.
\]
It makes no claim about eigenvalues, concentration, or a Gram event.

Proposition~\ref{prop:step-005-synthesis-scaling} defines the unreduced raw
and normalized tangent synthesis maps and proves the exact nonzero scalar
identity between them. Their ranges are, by direct two-sided comparison with
the setting generators, precisely \(\mathscr S_0^{\rm raw}\) and
\(\mathscr S_0^{\rm norm}\). This yields
\(\mathscr S_0^{\rm raw}=\mathscr S_0^{\rm norm}=\mathscr S_0\) without a
surrogate span or metric change.

Lemma~\ref{lem:step-005-gauge-kernel} identifies the exact two-dimensional
local scaling kernel on the previously established almost-sure nonzero branch and
gives the exact kernel and canonical gauge-image dimensions on every zero
branch.
Proposition~\ref{prop:step-005-quotient-range} then factors both synthesis
maps through the canonical-image quotient. On the nonzero branch this removes
exactly two gauge directions per component; on every branch the quotient,
unreduced generator frame, and complement-based reduced parameter frame have
the same image range \(\mathscr S_0\), with no injectivity claim. The same
proposition also proves the exact membership
\(\widehat D_0\in\mathscr S_0\).

Thus the range and membership claims are proved. In particular, the exact raw,
normalized, and quotient ranges all equal \(\mathscr S_0\), while
\(\widehat D_0\in\mathscr S_0\) and the raw target, residual, and Frobenius
conventions remain unchanged. No probability, Gram concentration,
full-rank Gaussian-array, Haar, projection-energy, or deficit conclusion is
asserted.

\end{proof}

\subsection{Product-Haar orbit factorization}\label{app:step-006}

\begin{proof}[Auxiliary facts in the present notation]


\medskip\noindent\textbf{Gaussianized coefficient arrays and scalar invariance.}\par

Lemma~\ref{lem:step-002-coefficient-gaussianization} and
Proposition~\ref{prop:step-002-balancing-invariance} state that, conditional
on a realized triple in \(\mathcal E_{\rm cond}\), the arrays
\(Z_M=[z_1^M\ \cdots\ z_k^M]\) are independent over modes and have iid
\(\mathcal N(0,1)\) entries, the matrices \(H_M\) are deterministic and
invertible. Write \(H_a:=H_A\), \(H_b:=H_B\), and \(H_c:=H_C\). Then
\[
\begin{aligned}
  \mathscr S_0
  =\operatorname{span}_{i\in[k]}\{&
    u\otimes H_bz_i^B\otimes H_cz_i^C,
    H_az_i^A\otimes v\otimes H_cz_i^C,\\
   &H_az_i^A\otimes H_bz_i^B\otimes w:
    u,v,w\in\mathbb R^r\}.
\end{aligned}\tag{1}
\]
Both the raw factors \(1/\sqrt n\) and the normalized factors \(1/\sqrt r\),
together with the balancing pair scalars, multiply entire tangent blocks by
nonzero numbers and hence do not change equation (1).

The proof below applies QR only to the standard Gaussian arrays \(Z_M\) from
Lemma~\ref{lem:step-002-coefficient-gaussianization}. It never asserts
that the balancing scalars are independent of their orientations. Their
removal from the range is exactly the algebraic conclusion above.


\medskip\noindent\textbf{exact raw and quotient range.}\par

Proposition~\ref{prop:step-005-synthesis-scaling} and
Proposition~\ref{prop:step-005-quotient-range} state that the raw and
normalized tangent synthesis maps, their canonical quotient maps, and any
complement-based reduced parameter frame all have image range
\(\mathscr S_0\). The target \(D_r\), represented tensor
\(\widehat D_0\), coefficient residual, and physical Frobenius metric remain
raw.

Equation (1) is therefore the exact
raw coefficient tangent range used below, not a normalized or
quotient surrogate. The internal subspace \(E\) below is allowed to use the
unreduced \(3kr\)-parameter synthesis map because the quotient
result proves equality of image ranges.


\medskip\noindent\textbf{uniqueness of Haar measure.}\par

We use the following statement in the present notation. On a compact group \(G\), every finite left-invariant
Borel measure is a nonnegative scalar multiple of normalized Haar measure.
Consequently, if random variables \((U,R)\in G\times\mathcal R\) satisfy
\((U,R)\stackrel d=(VU,R)\) for every deterministic \(V\in G\), then \(U\)
is Haar and independent of \(R\).

In
Lemma~\ref{lem:step-006-gaussian-qr-shape}, positive-diagonal QR is equivariant
under left multiplication, while a standard Gaussian matrix is left
orthogonally invariant. For a Borel shape event \(B\), the finite measure
\(A\mapsto\mathbb P(U\in A,R\in B)\) is therefore left invariant. Haar
uniqueness gives exactly the product factorization needed for independence.
The same fact on \(SO(r)\) identifies the orientation after the determinant
bit is removed.

\end{proof}



\begin{lemma}[Gaussian block QR separates orientation from the full rotated shape]
\label{lem:step-006-gaussian-qr-shape}
Under Assumptions~\ref{assump:independent_initialization} and
\ref{assump:rank_window}, where the latter gives \(r<k\), hence \(k\ge r\),
and Lemma~\ref{lem:step-002-coefficient-gaussianization}, let
\(Z\in\mathbb R^{r\times k}\) have iid \(\mathcal N(0,1)\) entries. Write
\[
  Z=[G\ W],\qquad G\in\mathbb R^{r\times r},
  \qquad W\in\mathbb R^{r\times(k-r)}.
\]
There are Borel functions \(O(Z)\in O(r)\) and
\(R(Z)\in\mathbb R^{r\times k}\) such that
\[
  Z=O(Z)R(Z)                                                   \tag{2}
\]
for every \(Z\), and, under the Gaussian law,
\[
  O(Z)\sim\operatorname{Haar}(O(r)),
  \qquad O(Z)\ \perp\!\!\!\perp\ R(Z).                       \tag{3}
\]
On the full-probability event \(\Omega_{\rm qr}=\{\det G\ne0\}\), if
\[
  G=O_0R_0
\]
is the unique QR decomposition with \(R_0\) upper triangular and strictly
positive diagonal, then the construction is explicitly
\[
  O(Z)=O_0,\qquad
  Y=O_0^{\mathsf T}W,\qquad
  R(Z)=[R_0\ Y].                                               \tag{4}
\]
Thus the orientation is independent not only of the triangular factor but of
the complete shape containing all rotated remaining Gaussian columns.

\end{lemma}

\begin{proof}
The determinant is a nonzero polynomial in the
entries of \(G\), and the Gaussian law has a density, so
\[
  \mathbb P(\Omega_{\rm qr}^c)=0.                              \tag{5}
\]
On \(\Omega_{\rm qr}\), positive-diagonal QR is unique. Moreover,
\(R_0\) is the unique upper-triangular positive-diagonal Cholesky factor of
\(G^{\mathsf T}G\), and
\[
  O_0=GR_0^{-1}.                                               \tag{6}
\]
The Cholesky map is continuous on the positive-definite cone, so \(O_0,R_0\)
are Borel functions of \(G\) there. On \(\Omega_{\rm qr}^c\), set
\(O(Z)=I_r\), and on all branches set
\[
  R(Z)=O(Z)^{\mathsf T}Z.
\]
This is a Borel extension and proves equation (2) pointwise. On
\(\Omega_{\rm qr}\) it agrees with equation (4).

For a deterministic \(U\in O(r)\), the positive-diagonal QR decomposition of
\(UG\) is
\[
  UG=(UO_0)R_0.                                                \tag{7}
\]
Because \(UG\stackrel d=G\), equation (7) implies
\[
  (O_0,R_0)\stackrel d=(UO_0,R_0).                            \tag{8}
\]
For every Borel set \(B\) in the triangular-factor space, the finite measure
\[
  \mu_B(A)=\mathbb P(O_0\in A,R_0\in B)
\]
on \(O(r)\) is left invariant by equation (8). Haar uniqueness therefore gives
\[
  \mu_B(A)=\operatorname{Haar}_{O(r)}(A)\,
            \mathbb P(R_0\in B),                             \tag{9}
\]
so \(O_0\) is Haar and independent of \(R_0\).

It remains to include the unused columns in the shape. The Gaussian block
\(W\) is independent of \(G\), hence of \((O_0,R_0)\). Conditional on
\((O_0,R_0)=(o,t)\), orthogonal invariance of every column of \(W\) gives
\[
  o^{\mathsf T}W\stackrel d=W,                                \tag{10}
\]
and this conditional law does not depend on \(o\) or \(t\). More explicitly,
for bounded Borel \(f\) and \(g\), define
\(h(t)=\mathbb E[g(t,W)]\). Then
\[
\begin{aligned}
  \mathbb E[f(O_0)g(R_0,O_0^{\mathsf T}W)]
  &=\mathbb E[f(O_0)h(R_0)]\\
  &=\mathbb E[f(O_0)]\,\mathbb E[h(R_0)]\\
  &=\mathbb E[f(O_0)]\,
    \mathbb E[g(R_0,O_0^{\mathsf T}W)].
\end{aligned}\tag{11}
\]
Thus \(O_0\) is independent of \((R_0,Y)\), proving equations (3)--(4). The
measurable extension changes the variables only on the null event in equation
(5), so the Haar and independence laws remain exact. Since invertibility of
\(G\) implies full row rank of \(Z\), every rank-deficient branch is contained
in the same null event. The argument also covers \(k=r+1\), where \(W\) has
exactly one column.


\end{proof}


\begin{proposition}[Reflection-bit absorption gives an independent Haar orientation in \(SO(r)\)]
\label{prop:step-006-reflection-absorption}
Under Lemma~\ref{lem:step-006-gaussian-qr-shape}, fix
\[
  J=\operatorname{diag}(-1,1,\ldots,1)\in O(r),
  \qquad \det J=-1,
\]
and write \(Z=O_0R\) for the lemma's factorization. Define the measurable
determinant bit, orientation, and reflected shape by
\[
  \varepsilon=\mathbf 1_{\{\det O_0=-1\}},\qquad
  Q=O_0J^{\varepsilon},\qquad
  \widetilde R=J^{\varepsilon}R.                              \tag{12}
\]
Then
\[
  Q\in SO(r),\qquad Z=Q\widetilde R,                           \tag{13}
\]
\(Q\) is Haar on \(SO(r)\), and
\[
  Q\ \perp\!\!\!\perp\ \widetilde R.                       \tag{14}
\]
The bit \(\varepsilon\) is stored entirely in \(\widetilde R\).

\end{proposition}

\begin{proof}
Since \(J^2=I_r\), equation (12) gives
\[
  \det Q=(\det O_0)(-1)^{\varepsilon}=1,
  \qquad
  Q\widetilde R=O_0J^{\varepsilon}J^{\varepsilon}R=O_0R=Z,
\]
which proves equation (13) on every branch.

Normalized Haar measure on \(O(r)\) assigns mass \(1/2\) to each determinant
component because right multiplication by \(J\) interchanges them. On the
component \(\det O_0=1\), \(Q=O_0\); on the component \(\det O_0=-1\), right
multiplication by \(J\) maps that component bijectively and measure
preservingly onto \(SO(r)\). Hence
\[
  Q\sim\operatorname{Haar}(SO(r)),\qquad
  Q\ \perp\!\!\!\perp\ \varepsilon.                         \tag{15}
\]
Lemma~\ref{lem:step-006-gaussian-qr-shape} gives
\(O_0\perp\!\!\!\perp R\). The measurable bijection
\(O_0\leftrightarrow(Q,\varepsilon)\) therefore gives
\((Q,\varepsilon)\perp\!\!\!\perp R\). Combining this with equation (15),
for bounded Borel \(f,g\),
\[
\begin{aligned}
  \mathbb E[f(Q)g(J^{\varepsilon}R)]
  &=\mathbb E\!\left[g(J^{\varepsilon}R)
      \mathbb E[f(Q)\mid\varepsilon,R]\right]\\
  &=\mathbb E[f(Q)]\,\mathbb E[g(J^{\varepsilon}R)],
\end{aligned}
\]
which proves equation (14). Thus absorbing the reflection does not leave a
determinant bit correlated with the exposed \(SO(r)\) orientation.


\end{proof}


\begin{proposition}[Product-Haar orientation--shape factorization]
\label{prop:step-006-product-shape}
\leavevmode\par\noindent
Under Assumptions~\ref{assump:independent_initialization} and
\ref{assump:rank_window}, where the latter gives \(r<k\), hence \(k\ge r\),
Lemma~\ref{lem:step-002-coefficient-gaussianization}, and
Proposition~\ref{prop:step-006-reflection-absorption}, there are measurable
factorizations
\[
  Z_A=Q_aR_A,\qquad Z_B=Q_bR_B,\qquad Z_C=Q_cR_C,              \tag{16}
\]
such that, conditional on the realized factors,
\[
  (Q_a,Q_b,Q_c)\sim
  \operatorname{Haar}(SO(r))^{\otimes3},
  \qquad
  (Q_a,Q_b,Q_c)\ \perp\!\!\!\perp\ (R_A,R_B,R_C).           \tag{17}
\]
Writing \(r_i^M\) for column \(i\) of \(R_M\), define
\[
\begin{aligned}
  E=\operatorname{span}_{i\in[k]}\{&
    u\otimes r_i^B\otimes r_i^C,
    r_i^A\otimes v\otimes r_i^C,
    r_i^A\otimes r_i^B\otimes w:\
   &u,v,w\in\mathbb R^r\}.
\end{aligned}\tag{18}
\]
Then \(E\) is a measurable random subspace,
\[
  (Q_a,Q_b,Q_c)\ \perp\!\!\!\perp\ E,
  \qquad \boxed{\dim(E)\le3kr}.                              \tag{19}
\]

\end{proposition}

\begin{proof}
Conditional on the realized factor triple, the
three arrays \(Z_A,Z_B,Z_C\) are independent standard Gaussian arrays by the
Gaussianization lemma. Apply
Proposition~\ref{prop:step-006-reflection-absorption} separately to each
array. Each pair \((Q_m,R_M)\) is a measurable function of its own mode array,
so the three pairs are independent. Within each pair the orientation is Haar
and independent of its shape. The joint law consequently factors as
\[
  \bigotimes_{m\in\{a,b,c\}}
    \left(\operatorname{Haar}_{SO(r)}\otimes\mathcal L(R_M)\right),
\]
which is exactly equation (17), not merely pairwise independence.

For measurability of equation (18), define the internal synthesis map
\[
\begin{aligned}
  \mathcal T_R((u_i,v_i,w_i)_{i=1}^k)
  =\sum_{i=1}^k(&u_i\otimes r_i^B\otimes r_i^C
    +r_i^A\otimes v_i\otimes r_i^C\\
    &+r_i^A\otimes r_i^B\otimes w_i).
\end{aligned}\tag{20}
\]
In standard bases, the matrix \(V_R\) of \(\mathcal T_R\) has entries that
are polynomial in the entries of \(R_A,R_B,R_C\), and
\(E=\operatorname{range}(V_R)\). Its orthogonal projector is the Borel
function
\[
  P_E=V_RV_R^\dagger,
  \qquad
  V_R^\dagger=\lim_{j\to\infty}
    (V_R^{\mathsf T}V_R+j^{-1}I)^{-1}V_R^{\mathsf T},          \tag{21}
\]
where the limit follows directly from a singular-value decomposition. Thus
the random subspace is measurable. Since it is a function only of the three
shapes, equation (17) proves the independence in equation (19).

Finally, the domain of \(\mathcal T_R\) has dimension \(3kr\), so
\[
  \dim(E)=\operatorname{rank}(\mathcal T_R)\le3kr.
\]
This remains true if some shape column vanishes or if cross-component
dependencies enlarge the synthesis kernel. It makes no injectivity or exact
dimension claim. At \(k=r+1\), each QR shape contains one rotated remaining
column, and the same construction and bound apply without modification.


\end{proof}


\begin{proposition}[Exact raw tangent equivariance]
\label{prop:step-006-exact-equivariance}
Under Lemma~\ref{lem:step-002-coefficient-gaussianization}, Proposition~\ref{prop:step-002-balancing-invariance}, Propositions~\ref{prop:step-005-synthesis-scaling} and
\ref{prop:step-005-quotient-range}, and
Proposition~\ref{prop:step-006-product-shape}, define
\[
  L=H_a\otimes H_b\otimes H_c,
  \qquad Q=Q_a\otimes Q_b\otimes Q_c.                         \tag{22}
\]
Then \(L\) and \(Q\) are invertible coefficient-tensor operators and
\[
  \boxed{
  \mathscr S_0
  =(H_a\otimes H_b\otimes H_c)
   (Q_a\otimes Q_b\otimes Q_c)E
  =LQE.}                                                       \tag{23}
\]
This is equality with the exact raw tangent range from the setting and the
quotient convention.

\end{proposition}

\begin{proof}
By equation (16), for every component and mode,
\[
  z_i^A=Q_ar_i^A,\qquad z_i^B=Q_br_i^B,
  \qquad z_i^C=Q_cr_i^C.                                     \tag{24}
\]
Let \(E_i^a,E_i^b,E_i^c\) denote the three component subspaces displayed in
equation (18). For the first one,
\[
\begin{aligned}
  LQ E_i^a
  &=\{H_aQ_au\otimes H_bQ_br_i^B\otimes H_cQ_cr_i^C:
      u\in\mathbb R^r\}\\
  &=\{u'\otimes H_bz_i^B\otimes H_cz_i^C:
      u'\in\mathbb R^r\},
\end{aligned}\tag{25}
\]
because \(H_aQ_a\) is invertible. The same argument in the other two free
modes gives
\[
\begin{aligned}
  LQ E_i^b
    &=\{H_az_i^A\otimes v'\otimes H_cz_i^C:v'\in\mathbb R^r\},\\
  LQ E_i^c
    &=\{H_az_i^A\otimes H_bz_i^B\otimes w':w'\in\mathbb R^r\}.
\end{aligned}\tag{26}
\]
Taking the span over \(i\) and comparing equations (25)--(26) with the
exact range formula (1) proves equation (23).

For completeness, the actual raw first-mode tangent block contains the
factor \(s_i^ys_i^z/n\), while its normalized counterpart contains
\(s_i^ys_i^z/r\); the analogous factors occur in the other two modes. The
balancing-invariance proposition proves that every such factor is
nonzero, including its exact scalar convention on the no-op zero branch, so
each multiplies a whole block by a bijective scalar. Thus no balancing scalar,
normalization scalar, or independence assertion about those scalars is hidden
in equation (23). Proposition~\ref{prop:step-005-quotient-range} then
identifies this range with the raw, normalized, and quotient ranges. The fixed coefficient
target \(D_r\), \(\widehat D_0\), and every physical-space object are not acted
on or redefined in this proof. 

\end{proof}

\begin{proof}[Proof of the subsection conclusion]

Lemma~\ref{lem:step-006-gaussian-qr-shape} applies positive-diagonal QR to
the first \(r\) columns of each standard Gaussian mode array. It proves
explicitly that the QR orientation is Haar on \(O(r)\) and
independent of the complete rotated shape
\([R_0\ O_0^{\mathsf T}W]\), not only of the triangular factor. The singular
first-block and rank-deficient branches are null; the piecewise definition
still gives a measurable exact factorization on those branches.

Proposition~\ref{prop:step-006-reflection-absorption} writes each Haar
\(O(r)\) orientation as \(Q_mJ^{\varepsilon_m}\), with \(Q_m\) Haar on
\(SO(r)\), and replaces the shape by \(J^{\varepsilon_m}R_M\). It proves both
the exact identity \(Z_M=Q_mR_M\) after this replacement and independence of
the exposed orientation from the reflected shape. Therefore no untracked
determinant bit remains outside the internal shape.

Proposition~\ref{prop:step-006-product-shape} uses independence of the three
Gaussian mode arrays to obtain the product law
  \(\operatorname{Haar}(SO(r))^{\otimes3}\), jointly independent of all three
shapes. It defines the measurable internal tangent subspace \(E\) from those
shapes and proves
\[
  (Q_a,Q_b,Q_c)\perp\!\!\!\perp E,
  \qquad \dim(E)\le3kr.
\]
The proof explicitly includes the boundary \(k=r+1\), null rank branch, and
possible rank loss of the internal synthesis map.

Finally, the exact identities \(z_i^M=Q_mr_i^M\), together with
Propositions~\ref{prop:step-002-balancing-invariance} and
\ref{prop:step-005-quotient-range}, give the blockwise identity
\[
  \mathscr S_0
  =(H_a\otimes H_b\otimes H_c)
   (Q_a\otimes Q_b\otimes Q_c)E.
\]
These results establish the claimed product-Haar factorization. No
projection of \(D_r\), Haar mean, concentration inequality, leverage bound,
or tangent-deficit conclusion is asserted.

\end{proof}

\subsection{Fixed-target product-Haar projection}\label{app:step-007}

\begin{proof}[Auxiliary facts in the present notation]


\medskip\noindent\textbf{independent product-Haar orientations and internal shape.}\par

Proposition~\ref{prop:step-006-product-shape} states that, conditional on the
realized factor triple, the orientations satisfy
\[
  (Q_a,Q_b,Q_c)\sim\operatorname{Haar}(SO(r))^{\otimes3}
\]
and are jointly independent of the measurable internal subspace \(E\), whose
dimension obeys \(d\le3kr\). Proposition~\ref{prop:step-006-exact-equivariance} states that this is the
same product orientation and subspace appearing in \(\mathscr S_0=LQE\).

Condition on all internal shapes.
Then \(E\) is deterministic and the conditional law of the three
orientations remains product Haar. The estimates proved below are uniform in
the resulting fixed \(E,d,X\), so the conditional application spends no
additional failure probability and uses no data-dependent supremum.


\medskip\noindent\textbf{curvature of a bi-invariant compact Lie group.}\par

In the metric convention used here, the relevant statement is as follows. Let a compact Lie group
carry the bi-invariant metric induced at the identity by an
\(\operatorname{Ad}\)-invariant inner product. For left-invariant tangent
vectors \(K,L,M\),
\[
  \nabla_KL=\frac12[K,L],
  \qquad
  R(K,L)M=-\frac14[[K,L],M].                                \tag{1}
\]
For \(SO(r)\) with
\[
  \langle K,L\rangle_F=-\operatorname{tr}(KL)
  \quad(K,L\in\mathfrak{so}(r)),                            \tag{2}
\]
the Killing form is
\[
  \operatorname{tr}(\operatorname{ad}_K\operatorname{ad}_L)
  =(r-2)\operatorname{tr}(KL),                              \tag{3}
\]
and consequently
\[
  \operatorname{Ric}_{SO(r)}=\frac{r-2}{4}g.                \tag{4}
\]

Equation (2) is exactly the unnormalized Hilbert--Schmidt metric used here:
there is no
factor \(1/2\), \(1/r\), or dimension normalization. The calculation of
equations (1)--(4), including the normalization in equation (3), is displayed
again inside Proposition~\ref{prop:step-007-product-tail}. The product metric
therefore has the same Ricci lower bound \((r-2)/4\).


\medskip\noindent\textbf{Bakry--Emery log-Sobolev criterion.}\par

In the metric convention used here, the relevant statement is as follows. If a connected complete
Riemannian manifold with normalized Riemannian volume \(\mu\) satisfies
\(\operatorname{Ric}\ge\rho g\) for some \(\rho>0\), then every smooth
real-valued \(f\) satisfies
\[
  \operatorname{Ent}_\mu(f^2)
  \le \frac2\rho\int\|\nabla f\|_g^2\,d\mu.                \tag{5}
\]
Here
\(\operatorname{Ent}_\mu(Y)=\int Y\log Y\,d\mu
-(\int Y\,d\mu)\log(\int Y\,d\mu)\).

For \(r\ge3\),
\(SO(r)^3\) is compact and connected, its normalized Riemannian volume is
product Haar, and equation (4) gives
\(\rho=(r-2)/4>0\). Product tensorization is already contained in the product
Ricci calculation, so equation (5) applies directly to the three-factor
metric, with no extra factor of three. Applying equation (5) to
\(f=\exp(\lambda(h-\mathbb Eh)/2)\) yields the Herbst bound derived explicitly
in Proposition~\ref{prop:step-007-product-tail}. No theorem number or
paper-specific source label is asserted.

\end{proof}



\begin{lemma}[Sequential \(SO(r)\) twirling gives the exact projection mean]
\label{lem:step-007-sequential-twirl}
Under Assumption~\ref{assump:dimension}, using only \(r\ge3\), let
\(E\subset\mathcal H=(\mathbb R^r)^{\otimes3}\) be a fixed subspace of
dimension \(d\), let \(X\in\mathcal H\setminus\{0\}\), and let
\(Q_a,Q_b,Q_c\) be independent Haar matrices in \(SO(r)\). Then
\[
  \mathbb E\frac{\|P_{(Q_a\otimes Q_b\otimes Q_c)E}X\|_F^2}
  {\|X\|_F^2}
  =\frac d{r^3}.                                             \tag{6}
\]

\end{lemma}

\begin{proof}
We first prove the one-mode identity in the exact
real representation used here. Let \(V=\mathbb R^r\), let \(W\) be any
finite-dimensional real inner-product space, and let
\(A\in\operatorname{End}(V\otimes W)\). Define
\[
  \mathcal T_V(A)
  =\int_{SO(r)}(U\otimes I_W)A(U^{\mathsf T}\otimes I_W)\,dU. \tag{7}
\]
Haar invariance shows that \(\mathcal T_V(A)\) commutes with
\(S\otimes I_W\) for every \(S\in SO(r)\).

For completeness, the commutant of the natural real \(SO(r)\)-action is
scalar when \(r\ge3\). Indeed, if \(C\in\operatorname{End}(V)\) commutes
with \(SO(r)\), then for every unit vector \(u\), the vector \(Cu\) is fixed
by the subgroup that fixes \(u\) and rotates \(u^\perp\). The fixed space of
that subgroup is \(\operatorname{span}\{u\}\), so every unit vector is an
eigenvector of \(C\). Applying linearity to two independent vectors and their
sum shows that all eigenvalues coincide; hence \(C\) is scalar. Taking matrix
coefficients of \(\mathcal T_V(A)\) between arbitrary vectors of \(W\) now
gives
\[
  \mathcal T_V(A)=I_V\otimes B
\]
for some \(B\in\operatorname{End}(W)\). Partial trace over \(V\) is unchanged
by the conjugations in equation (7), whereas
\(\operatorname{Tr}_V(I_V\otimes B)=rB\). Therefore
\[
  \boxed{\mathcal T_V(A)
  =\frac{I_V}{r}\otimes\operatorname{Tr}_V(A).}             \tag{8}
\]

Let \(P_E\) be the orthogonal projector onto \(E\). Since the tensor product
of orthogonal maps is orthogonal,
\[
  P_{(Q_a\otimes Q_b\otimes Q_c)E}
  =(Q_a\otimes Q_b\otimes Q_c)P_E
   (Q_a^{\mathsf T}\otimes Q_b^{\mathsf T}\otimes Q_c^{\mathsf T}). \tag{9}
\]
Apply equation (8) successively in modes \(a,b,c\). Independence permits the
three iterated Haar integrals, and the commuting mode actions give
\[
\begin{aligned}
  \mathbb E[QP_EQ^{\mathsf T}]
  &=\frac{I_r}{r}\otimes\frac{I_r}{r}\otimes\frac{I_r}{r}
    \operatorname{Tr}(P_E)\\
  &=\frac d{r^3}I_{\mathcal H}.
\end{aligned}\tag{10}
\]
Thus
\[
  \mathbb Eh
  =\frac{\langle X,\mathbb E[QP_EQ^{\mathsf T}]X\rangle}
         {\|X\|_F^2}
  =\frac d{r^3},
\]
which proves equation (6).

If \(d=0\), then \(P_E=0\) and \(h=0\) identically, in agreement with the
formula. The proof uses only \(\operatorname{Tr}(P_E)=d\), so it remains
valid at every actual endpoint, including \(d=3kr\) whenever a subspace of
that dimension is supplied. No separate inequality \(3kr<r^3\) is assumed;
the existence of \(E\subset\mathcal H\) already enforces \(d\le r^3\).


\end{proof}


\begin{lemma}[Metric-exact Lipschitz bound for product rotations]
\label{lem:step-007-lipschitz}
Fix \(E\subset\mathcal H\) and \(X\in\mathcal H\setminus\{0\}\). Equip each
\(SO(r)\) factor with the bi-invariant metric whose identity inner product is
\[
  \langle K,L\rangle_F=-\operatorname{tr}(KL),
  \qquad K,L\in\mathfrak{so}(r),
\]
and equip \(SO(r)^3\) with the product metric. Then
\[
  \boxed{\operatorname{Lip}(h)\le2\sqrt3}.                  \tag{11}
\]
Equivalently, a product tangent vector with skew generators
\((K_a,K_b,K_c)\) has squared speed
\[
  \sum_{m\in\{a,b,c\}}\|K_m\|_F^2,
\]
with no normalization factor, and the directional derivative of \(h\) is at
most \(2\sqrt3\) times that speed.

\end{lemma}

\begin{proof}
By left invariance, a product geodesic through
\((Q_a,Q_b,Q_c)\) can be written
\[
  Q_m(s)=e^{sK_m}Q_m,
  \qquad K_m\in\mathfrak{so}(r),
\]
and its speed is exactly
\[
  \left(\|K_a\|_F^2+\|K_b\|_F^2+\|K_c\|_F^2\right)^{1/2}.  \tag{12}
\]
The three tensor-mode generators commute, so the induced tensor action is
\[
\begin{aligned}
  Q(s)&=e^{s\mathcal K}Q,\\
  \mathcal K
  &=K_a\otimes I_r\otimes I_r
    +I_r\otimes K_b\otimes I_r
    +I_r\otimes I_r\otimes K_c.
\end{aligned}\tag{13}
\]
This gives the explicit tensor-generator estimate
\[
\begin{aligned}
  \|\mathcal K\|_{\rm op}
  &\le \|K_a\|_{\rm op}+\|K_b\|_{\rm op}+\|K_c\|_{\rm op}\\
  &\le \|K_a\|_F+\|K_b\|_F+\|K_c\|_F\\
  &\le \sqrt3
  \left(\|K_a\|_F^2+\|K_b\|_F^2+\|K_c\|_F^2\right)^{1/2}.
\end{aligned}\tag{14}
\]

Write \(P=P_{QE}\). Along the geodesic,
\[
  P(s)=e^{s\mathcal K}Pe^{-s\mathcal K},
  \qquad P'(0)=[\mathcal K,P].                               \tag{15}
\]
Consequently,
\[
\begin{aligned}
  |h'(0)|
  &=\frac{|\langle X,[\mathcal K,P]X\rangle|}{\|X\|_F^2}\\
  &\le\|[\mathcal K,P]\|_{\rm op}
  \le2\|\mathcal K\|_{\rm op}\\
  &\le2\sqrt3
  \left(\|K_a\|_F^2+\|K_b\|_F^2+\|K_c\|_F^2\right)^{1/2}.
\end{aligned}\tag{16}
\]
Here \(\|P\|_{\rm op}\le1\), including the case \(E=\{0\}\). The bound
holds at every point and along every geodesic. Integrating equation (16)
along a minimizing geodesic proves equation (11) for exactly the metric in
the statement. 

\end{proof}


\begin{proposition}[Ricci--log-Sobolev--Herbst tail in the unnormalized product metric]
\label{prop:step-007-product-tail}
Under Assumption~\ref{assump:dimension}, using only \(r\ge3\), and under the
hypotheses of Lemma~\ref{lem:step-007-sequential-twirl}, one has for every
\(t>0\)
\[
  \mathbb P(h-\mathbb Eh\ge t)
  \le \exp\!\left(-\frac{(r-2)t^2}{96}\right)
  \le \exp\!\left(-\frac{rt^2}{288}\right).                \tag{17}
\]
In particular, for \(0<t\le1\), equation (Haar) holds with the explicit
universal choice
\[
  \boxed{c_H=\frac1{288}}.                                  \tag{18}
\]

\end{proposition}

\begin{proof}
We first verify the metric normalization rather
than importing a scale-free concentration slogan. On
\(\mathfrak{so}(r)\), use the inner product in equation (2). The matrices
\[
  F_{ij}=\frac{E_{ij}-E_{ji}}{\sqrt2},
  \qquad 1\le i<j\le r,
\]
form an orthonormal basis. The Koszul formula for left-invariant fields and
the \(\operatorname{Ad}\)-invariance of equation (2) give
\[
  \nabla_KL=\frac12[K,L],
  \qquad
  R(K,L)M=-\frac14[[K,L],M].                                \tag{19}
\]
Hence
\[
  \operatorname{Ric}(K,K)
  =\frac14\sum_{i<j}\|[K,F_{ij}]\|_F^2
  =-\frac14\operatorname{tr}(\operatorname{ad}_K^2).       \tag{20}
\]
A direct summation in the displayed matrix-unit basis gives, for
\(K,L\in\mathfrak{so}(r)\),
\[
  \operatorname{tr}(\operatorname{ad}_K\operatorname{ad}_L)
  =(r-2)\operatorname{tr}(KL).                              \tag{21}
\]
Since \(\operatorname{tr}(K^2)=-\|K\|_F^2\), equations
(20)--(21) yield the exact identity
\[
  \operatorname{Ric}_{SO(r)}(K,K)
  =\frac{r-2}{4}\|K\|_F^2.                                 \tag{22}
\]
This calculation is tied to the unnormalized basis \(F_{ij}\); replacing
the metric by \(-\operatorname{tr}(KL)/2\) or by a dimension-normalized
metric would change equation (22), but no such replacement is made here.

For the product metric and
\(\mathbf K=(K_a,K_b,K_c)\), Ricci curvature adds over factors:
\[
\begin{aligned}
  \operatorname{Ric}_{SO(r)^3}(\mathbf K,\mathbf K)
  &=\sum_{m\in\{a,b,c\}}
    \frac{r-2}{4}\|K_m\|_F^2\\
  &=\rho\|\mathbf K\|_g^2,
  \qquad \rho=\frac{r-2}{4}.
\end{aligned}\tag{23}
\]
Thus the standard Bakry--Emery criterion, equation (5), applies to normalized
product Haar measure. Let
\[
  Y=h-\mathbb Eh,
  \qquad
  \psi(\lambda)=\log\mathbb E e^{\lambda Y},
  \qquad \lambda>0.
\]
Apply equation (5) to \(f=e^{\lambda Y/2}\). By
Lemma~\ref{lem:step-007-lipschitz}, \(\|\nabla h\|_g\le L\) with
\(L=2\sqrt3\), so
\[
\begin{aligned}
  \operatorname{Ent}(e^{\lambda Y})
  &\le \frac2\rho\frac{\lambda^2}{4}
    \mathbb E[e^{\lambda Y}\|\nabla h\|_g^2]\\
  &\le \frac{\lambda^2L^2}{2\rho}\mathbb E e^{\lambda Y}.
\end{aligned}\tag{24}
\]
On the other hand,
\[
  \frac{\operatorname{Ent}(e^{\lambda Y})}
       {\mathbb E e^{\lambda Y}}
  =\lambda\psi'(\lambda)-\psi(\lambda).                   \tag{25}
\]
Dividing equations (24)--(25) by \(\lambda^2\), integrating the derivative
of \(\psi(\lambda)/\lambda\) from \(0\) to \(\lambda\), and using
\(\mathbb EY=0\), gives
\[
  \psi(\lambda)\le\frac{\lambda^2L^2}{2\rho}.             \tag{26}
\]
Chernoff's bound and optimization at \(\lambda=\rho t/L^2\) now give
\[
\begin{aligned}
  \mathbb P(Y\ge t)
  &\le\inf_{\lambda>0}
    \exp\!\left(-\lambda t+\frac{\lambda^2L^2}{2\rho}\right)\\
  &=\exp\!\left(-\frac{\rho t^2}{2L^2}\right)
   =\exp\!\left(-\frac{(r-2)t^2}{96}\right),
\end{aligned}\tag{27}
\]
because \(L^2=12\) and \(\rho=(r-2)/4\). For \(r\ge3\),
\(r-2\ge r/3\), proving the second inequality in equation (17). Finally,
\[
  \exp(-rt^2/288)
  \le8\exp(-c_Hrt^2)
  \quad\text{with}\quad c_H=1/288,                         \tag{28}
\]
which proves equation (18) and (Haar). The derivation actually holds for all
\(t>0\); restricting to \(0<t\le1\) gives the stated range.

When \(d=0\), \(h=0\) and the left side of equation (17) is zero for every
\(t>0\), so the conclusion also holds without invoking curvature. No case
\(X=0\) is introduced: the row quantifies over nonzero \(X\), which
is necessary for the displayed normalization. 

\end{proof}

\begin{proof}[Proof of the subsection conclusion]

Proposition~\ref{prop:step-006-product-shape} permits conditioning
on the internal shapes while retaining independent Haar
\(Q_a,Q_b,Q_c\in SO(r)\). For every resulting fixed \(E\) and every fixed
nonzero \(X\), Lemma~\ref{lem:step-007-sequential-twirl} proves
\[
  \mathbb Eh=\frac d{r^3}
\]
by three exact current-notation twirls. Lemma~\ref{lem:step-007-lipschitz}
computes the tensor generator and proves
\(\operatorname{Lip}(h)\le2\sqrt3\) in the unnormalized product
Hilbert--Schmidt geodesic metric. Proposition~\ref{prop:step-007-product-tail}
computes \(\operatorname{Ric}=(r-2)g/4\) in that same metric, applies the
metric-matched log-Sobolev inequality and Herbst argument, and obtains
\[
  \mathbb P\!\left(h\ge\frac d{r^3}+t\right)
  \le\exp(-rt^2/288)
  \le8\exp(-c_Hrt^2),
  \qquad c_H=1/288,
\]
for \(0<t\le1\).

These conclusions are uniform in fixed \(E,d,X\), include \(d=0\) and every
actual upper endpoint such as \(d=3kr\), and use only \(r\ge3\) from the
large-\(r\) setting. The proof introduces no operator norm over targets, net,
data-dependent supremum, \(O(r)\) surrogate, transformed target, or claim
about \(LQE\). Thus Proposition~\ref{prop:step-007-product-tail} holds for
the fixed objects and with the scope stated there.

\end{proof}

\subsection{Elliptic transfer and raw leverage}\label{app:step-008}

\begin{proof}[Auxiliary facts in the present notation]


\medskip\noindent\textbf{realized conditioning and elliptic singular intervals.}\par

Proposition~\ref{prop:step-001-realized-conditioning} proves the generated
event \(\mathcal E_{\rm cond}\). On this event, every realized mode matrix has
singular values in \([\kappa_1^{-1},\kappa_1]\). The elliptic maps \(H_m\)
appearing in Proposition~\ref{prop:step-006-exact-equivariance} are the
reciprocal-singular-value maps of
Lemma~\ref{lem:step-002-coefficient-gaussianization}, so
\[
  \sigma(H_m)\subset[\kappa_1^{-1},\kappa_1],
  \qquad m\in\{a,b,c\}.
  \tag{1}
\]

Proposition~\ref{prop:step-008-elliptic-transfer}
applies equation (1) to
\(L=H_a\otimes H_b\otimes H_c\). The reciprocal interval is unchanged
because its endpoints are reciprocal. No additional conditioning event is
introduced.


\medskip\noindent\textbf{exact product-Haar tangent factorization.}\par

Proposition~\ref{prop:step-006-product-shape} and
Proposition~\ref{prop:step-006-exact-equivariance} give, conditional on every
realized factor triple in \(\mathcal E_{\rm cond}\),
\[
  \mathscr S_0=LQE,
  \qquad d=\dim(E)\le3kr,
  \qquad
  (Q_a,Q_b,Q_c)\sim\operatorname{Haar}(SO(r))^{\otimes3},
\]
with the orientations jointly independent of \(E\).

Condition first on the realized
factor triple and then on the three internal shapes. This fixes \(L,E,d\) and
\(X=L^{\mathsf T}D_r\), while the conditional law of \(Q\) remains product
Haar. The identity is with the exact raw tangent range, so the projection
formula below has no span-transfer residual.


\medskip\noindent\textbf{fixed-target product-Haar mean and tail.}\par

Lemma~\ref{lem:step-007-sequential-twirl} and
Proposition~\ref{prop:step-007-product-tail} prove equation (Haar) uniformly
over every fixed \(E,d\) and fixed nonzero \(X\), with \(c_H=1/288\).

The argument uses the single fixed tensor
\[
  X=L^{\mathsf T}D_r.
\]
It is nonzero because \(L\) is invertible and \(D_r\ne0\). The deviation is
\(t=\tau_\kappa\), and \(0<\tau_\kappa\le1\) is verified before equation
(Haar) is applied. Uniformity in the conditioned shape permits integration
over the shapes without a union bound or confidence loss.

\end{proof}



\begin{lemma}[Exact oblique-basis projection transfer]
\label{lem:step-008-oblique-projection}
Under Proposition~\ref{prop:step-006-exact-equivariance} \(\mathscr S_0=LQE\), let
\(E\subset(\mathbb R^r)^{\otimes3}\) have dimension \(d\), let \(L\) be
invertible, and, when \(d>0\), let \(U\) have orthonormal columns spanning
\(QE\). For every coefficient tensor \(x\), define
\[
  X=L^{\mathsf T}x,
  \qquad b=U^{\mathsf T}X.
\]
If \(d>0\), then
\[
  P_{LQE}
  =LU\bigl(U^{\mathsf T}L^{\mathsf T}LU\bigr)^{-1}
    U^{\mathsf T}L^{\mathsf T},                            \tag{2}
\]
and
\[
  \|P_{LQE}x\|_F^2
  =b^{\mathsf T}\bigl(U^{\mathsf T}L^{\mathsf T}LU\bigr)^{-1}b
  \le\sigma_{\min}(L)^{-2}\|P_{QE}X\|_F^2.              \tag{3}
\]
If \(d=0\), then \(QE=LQE=\{0\}\), and both sides of the inequality in
equation (3) are zero.

\end{lemma}

\begin{proof}
Suppose first that \(d>0\). The columns of
\(LU\) span \(LQE\). Since \(L\) is invertible and \(U\) has full column
rank, \(LU\) has full column rank. The orthogonal projector onto the range of
a full-column-rank matrix \(A\) is
\(A(A^{\mathsf T}A)^{-1}A^{\mathsf T}\); substituting \(A=LU\) proves
equation (2) directly.

Because an orthogonal projector is symmetric and idempotent,
\[
\begin{aligned}
  \|P_{LQE}x\|_F^2
  &=\langle x,P_{LQE}x\rangle_F\\
  &=b^{\mathsf T}\bigl(U^{\mathsf T}L^{\mathsf T}LU\bigr)^{-1}b.
\end{aligned}\tag{4}
\]
For every \(y\in\mathbb R^d\), orthonormality of \(U\) gives
\[
  y^{\mathsf T}U^{\mathsf T}L^{\mathsf T}LUy
  =\|LUy\|_F^2
  \ge\sigma_{\min}(L)^2\|Uy\|_F^2
  =\sigma_{\min}(L)^2\|y\|_2^2.
\]
Hence
\[
  \bigl(U^{\mathsf T}L^{\mathsf T}LU\bigr)^{-1}
  \preceq\sigma_{\min}(L)^{-2}I_d.                         \tag{5}
\]
Moreover,
\[
  \|b\|_2^2=\|U^{\mathsf T}X\|_2^2
  =\|UU^{\mathsf T}X\|_F^2
  =\|P_{QE}X\|_F^2.                                       \tag{6}
\]
Combining equations (4)--(6) proves equation (3). When \(d=0\), both
subspaces are the zero subspace, so the stated zero-case conclusion is
immediate. 

\end{proof}


\begin{proposition}[Elliptic transfer with the exact condition-number loss]
\label{prop:step-008-elliptic-transfer}
Under the generated event \(\mathcal E_{\rm cond}\), the elliptic
singular intervals, Proposition~\ref{prop:step-006-exact-equivariance} \(\mathscr S_0=LQE\), and
Lemma~\ref{lem:step-008-oblique-projection},
\[
  \sigma_{\min}(L)\ge\kappa_1^{-3},
  \qquad
  \sigma_{\max}(L)\le\kappa_1^3.                          \tag{7}
\]
Consequently, for every nonzero coefficient tensor \(x\), with
\(X=L^{\mathsf T}x\),
\[
\begin{aligned}
  \|P_{\mathscr S_0}x\|_F^2
  &\le \sigma_{\min}(L)^{-2}
      h(Q;E,X)\|X\|_F^2\\
  &\le \kappa_1^{12}h(Q;E,X)\|x\|_F^2.
\end{aligned}\tag{8}
\]
In particular,
\[
  \|P_{\mathscr S_0}D_r\|_F^2
  \le\kappa_1^{12}h(Q;E,L^{\mathsf T}D_r)\|D_r\|_F^2.
                                                               \tag{9}
\]

\end{proposition}

\begin{proof}
The singular values of a tensor product are the
pairwise products of the singular values of its factors. Applying the
intervals in equation (1) to
\(L=H_a\otimes H_b\otimes H_c\) gives
\[
  \sigma_{\min}(L)
  =\prod_{m\in\{a,b,c\}}\sigma_{\min}(H_m)
  \ge\kappa_1^{-3},
\]
and
\[
  \sigma_{\max}(L)
  =\prod_{m\in\{a,b,c\}}\sigma_{\max}(H_m)
  \le\kappa_1^3,
\]
proving equation (7). These are the simultaneous worst cases in which all
three lower singular values equal \(\kappa_1^{-1}\), or all three upper
singular values equal \(\kappa_1\).

Lemma~\ref{lem:step-008-oblique-projection} and
\(\|X\|_F\le\sigma_{\max}(L)\|x\|_F\) give
\[
\begin{aligned}
  \|P_{\mathscr S_0}x\|_F^2
  &\le\sigma_{\min}(L)^{-2}\|P_{QE}X\|_F^2\\
  &=\sigma_{\min}(L)^{-2}h(Q;E,X)\|X\|_F^2\\
  &\le\sigma_{\min}(L)^{-2}\sigma_{\max}(L)^2
       h(Q;E,X)\|x\|_F^2\\
  &\le\kappa_1^{12}h(Q;E,X)\|x\|_F^2,
\end{aligned}\tag{10}
\]
because the lower and upper singular-value factors contribute
\(\kappa_1^6\) each. This proves equation (8), and \(x=D_r\) gives equation
(9). If \(d=0\), the same conclusion holds with both projection energies
equal to zero. 

\end{proof}


\begin{proposition}[Exact raw target leverage bound]
\label{prop:step-008-raw-leverage}
Under Assumption~\ref{assump:rank_window}, the dimension bound
\(d\le3kr\), the fixed-target Haar estimate, and
Proposition~\ref{prop:step-008-elliptic-transfer}, define
\[
  \tau_\kappa:=\frac1{4\kappa_1^{12}}.                     \tag{11}
\]
If
\[
  r\ge\left\lceil(12\kappa_1^{12})^{4/3}\right\rceil,     \tag{12}
\]
then, conditional on every realized factor triple in
\(\mathcal E_{\rm cond}\), equation (LEV) holds with conditional failure at
most
\[
  8\exp\!\left(-\frac{c_Hr}{16\kappa_1^{24}}\right).
                                                               \tag{13}
\]

\end{proposition}

\begin{proof}
Since \(\kappa\ge1\) and
\(\kappa_1=2\kappa^2\ge1\),
\[
  0<\tau_\kappa\le\frac14\le1,                            \tag{14}
\]
so it is an admissible deviation in equation (Haar). Assumption~\ref{assump:rank_window}
and the dimension bound give, including the maximal allowed rank,
\[
  \frac d{r^3}
  \le\frac{3kr}{r^3}
  =\frac{3k}{r^2}
  \le3r^{-3/4}.                                             \tag{15}
\]
The threshold in equation (12) is exactly the condition
\[
  r^{3/4}\ge12\kappa_1^{12},
\]
and therefore
\[
  3r^{-3/4}\le\frac1{4\kappa_1^{12}}=\tau_\kappa.          \tag{16}
\]

Condition now on the realized factors and the internal shapes. Then
\(E,d\) are fixed, \(Q\) remains product Haar, and
\(X=L^{\mathsf T}D_r\ne0\) is fixed. Applying equation (Haar) with
\(t=\tau_\kappa\), and using equations (15)--(16), shows that outside an
event of conditional probability at most
\[
  8e^{-c_Hr\tau_\kappa^2}
  =8\exp\!\left(-\frac{c_Hr}{16\kappa_1^{24}}\right),       \tag{17}
\]
one has
\[
  h(Q;E,L^{\mathsf T}D_r)
  \le\frac d{r^3}+\tau_\kappa
  \le2\tau_\kappa
  =\frac1{2\kappa_1^{12}}.                                 \tag{18}
\]
Proposition~\ref{prop:step-008-elliptic-transfer} then gives
\[
\begin{aligned}
  \|P_{\mathscr S_0}D_r\|_F^2
  &\le\kappa_1^{12}h(Q;E,L^{\mathsf T}D_r)\|D_r\|_F^2\\
  &\le\frac12\|D_r\|_F^2
   =\frac r2,
\end{aligned}\tag{19}
\]
because the \(r\) summands in \(D_r\) are orthonormal and hence
\(\|D_r\|_F^2=r\). This is exactly (LEV) for the raw target and exact raw
tangent span.

The bound in equation (17) is uniform in the internal shapes. Integrating
over them preserves the same conditional failure probability for every
realized factor triple in \(\mathcal E_{\rm cond}\); no union bound or
operator supremum is used. The case \(d=0\) is deterministic and already
satisfies equation (19). 

\end{proof}


\begin{lemma}[Explicit conversion to the polynomial failure budget]
\label{lem:step-008-tail-conversion}
Let
\[
  a_\kappa:=\frac{c_H}{16\kappa_1^{24}},                   \tag{20}
\]
and define the explicit threshold
\[
  r_{0,\rm LEV}(\kappa)
  :=\left\lceil
    \max\left\{
      3,
      (12\kappa_1^{12})^{4/3},
      \left(\frac{20+\log8}{a_\kappa}\right)^2
    \right\}
  \right\rceil.                                            \tag{21}
\]
For every integer \(r\ge r_{0,\rm LEV}(\kappa)\),
\[
  8e^{-a_\kappa r}\le r^{-20}.                             \tag{22}
\]

\end{lemma}

\begin{proof}
For every \(r\ge1\),
\[
  \log r\le\sqrt r.                                       \tag{23}
\]
Indeed, \(f(r)=\sqrt r-\log r\) decreases on \([1,4]\), increases on
\([4,\infty)\), and \(f(4)=2-\log4>0\). The last term in the maximum in
equation (21) gives
\[
  \sqrt r\ge\frac{20+\log8}{a_\kappa}.
\]
Using equation (23) and \(\sqrt r\ge1\),
\[
\begin{aligned}
  a_\kappa r
  &=a_\kappa\sqrt r\,\sqrt r\\
  &\ge(20+\log8)\sqrt r\\
  &\ge20\log r+\log8.
\end{aligned}\tag{24}
\]
Therefore
\[
  8e^{-a_\kappa r}
  \le8e^{-20\log r-\log8}
  =r^{-20},
\]
which is equation (22). The same threshold includes equation (12), so it
simultaneously discharges the mean-dominance and tail-conversion conditions.


\end{proof}

\begin{proof}[Proof of the subsection conclusion]

Lemma~\ref{lem:step-008-oblique-projection} gives the exact orthogonal
projection formula for the anisotropic subspace \(LQE\), including the
degenerate case \(d=0\), and proves
\[
  \|P_{LQE}x\|_F^2
  \le\sigma_{\min}(L)^{-2}\|P_{QE}L^{\mathsf T}x\|_F^2.
\]
Proposition~\ref{prop:step-008-elliptic-transfer} applies the singular intervals in all three modes and proves
\[
  \sigma_{\min}(L)\ge\kappa_1^{-3},
  \qquad
  \sigma_{\max}(L)\le\kappa_1^3,
\]
so the complete worst-case anisotropic loss is exactly
\(\kappa_1^{12}\).

Proposition~\ref{prop:step-008-raw-leverage} cites
Assumption~\ref{assump:rank_window} to use the maximal-rank inequality
\[
  \frac d{r^3}\le\frac{3k}{r^2}\le3r^{-3/4}
  \le\tau_\kappa,
  \qquad
  \tau_\kappa=\frac1{4\kappa_1^{12}},
\]
and applies the fixed-target Haar estimate at the single exact
tensor \(X=L^{\mathsf T}D_r\). It obtains
\(h\le2\tau_\kappa\), hence
\[
  \|P_{\mathscr S_0}D_r\|_F^2
  \le\kappa_1^{12}(2\tau_\kappa)\|D_r\|_F^2
  =\frac r2.
\]
Finally, Lemma~\ref{lem:step-008-tail-conversion} gives an explicit
\(r_{0,\rm LEV}(\kappa)\) for which the conditional failure in equation
(13) is at most \(r^{-20}\). Uniformity over the conditioned shapes permits
the same bound after integrating them out. These four named results prove
that (LEV) holds for the raw \(D_r\) and exact raw
\(\mathscr S_0\), with conditional failure at most \(r^{-20}\), without a
surrogate residual or target supremum.

\end{proof}

\subsection{Raw tangent-deficit witness}\label{app:step-009}

\begin{proof}[Auxiliary facts in the present notation]


\medskip\noindent\textbf{exact raw membership.}\par

Proposition~\ref{prop:step-005-quotient-range} proves, on every branch,
\[
  \widehat D_0\in\mathscr S_0.
  \tag{1}
\]
It also proves that \(D_r\), \(D_r-\widehat D_0\), and the coefficient
Frobenius geometry remain raw.

Equation (1) is used for the exact
cancellation of \(\widehat D_0\) against the normal witness. No normalized
target, quotient representative, or transformed residual is imported.


\medskip\noindent\textbf{exact raw leverage.}\par

Proposition~\ref{prop:step-008-raw-leverage}, together with Lemma~\ref{lem:step-008-tail-conversion}, proves
\[
  \|P_{\mathscr S_0}D_r\|_F^2\le \frac r2
\]
for the exact raw objects. For every realized factor triple in
\(\mathcal E_{\rm cond}\), this holds with conditional failure at most
\(r^{-20}\) for the sufficiently-large-\(r\) regime.

The present proof uses only this
projection-energy conclusion and its already-spent failure budget. It does not
reopen the Haar, anisotropic-transfer, rank-window, or tail argument.

\end{proof}



\begin{lemma}[Exact membership and nonzero raw normal energy]
\label{lem:step-009-normal-energy}
Under the setting definitions, Proposition~\ref{prop:step-005-quotient-range}, and Proposition~\ref{prop:step-008-raw-leverage}, if (LEV) holds, then
\[
  \widehat D_0\in\mathscr S_0,
  \qquad
  \|P_{\mathscr S_0^\perp}D_r\|_F^2\ge\frac r2>0.
  \tag{2}
\]
Thus the denominator in the setting-required definition of \(W_0\) is
nonzero.

\end{lemma}

\begin{proof}
The membership in equation (2) is exactly equation (1) from
Proposition~\ref{prop:step-005-quotient-range}. The tensors
\(e_j\otimes e_j\otimes e_j\), \(j\in[r]\), are orthonormal, so
\[
  \|D_r\|_F^2
  =\left\|\sum_{j=1}^r e_j\otimes e_j\otimes e_j\right\|_F^2
  =r.
  \tag{3}
\]
Orthogonal Pythagorean decomposition in the raw coefficient space gives
\[
\begin{aligned}
  \|P_{\mathscr S_0^\perp}D_r\|_F^2
  &=\|D_r\|_F^2-\|P_{\mathscr S_0}D_r\|_F^2\\
  &\ge r-\frac r2
   =\frac r2.
\end{aligned}\tag{4}
\]
Because \(r\ge1\), the right-hand side is strictly positive. Equality in
(LEV) merely makes equation (4) an equality and still leaves normal energy
\(r/2>0\). Conversely, a zero denominator would make the left-hand side of
equation (4) zero, contradicting \(r/2>0\). 

\end{proof}


\begin{proposition}[Raw unit normal witness and tangent deficit]
\label{prop:step-009-raw-witness}
Under the setting definitions and
Lemma~\ref{lem:step-009-normal-energy}, if (LEV) holds, define exactly
\[
  W_0:=\frac{P_{\mathscr S_0^\perp}D_r}
  {\|P_{\mathscr S_0^\perp}D_r\|_F}.
  \tag{5}
\]
Then
\[
  \|W_0\|_F=1,
  \qquad
  W_0\perp\mathscr S_0,
  \qquad
  \langle D_r-\widehat D_0,W_0\rangle
  \ge\delta_0\|D_r\|_F,
  \tag{6}
\]
where \(\delta_0=1/8\). Consequently,
\[
  (\mathrm{LEV})\quad\Longrightarrow\quad
  \mathcal E_{\rm deficit}.
  \tag{7}
\]

\end{proposition}

\begin{proof}
Lemma~\ref{lem:step-009-normal-energy} makes the
denominator in equation (5) strictly positive, so the impossible
zero-denominator branch does not occur under (LEV). By construction,
\(W_0\) has unit Frobenius norm and lies in \(\mathscr S_0^\perp\).
The exact membership \(\widehat D_0\in\mathscr S_0\) gives
\(\langle\widehat D_0,W_0\rangle=0\). Therefore
\[
\begin{aligned}
  \langle D_r-\widehat D_0,W_0\rangle
  &=\left\langle D_r,
    \frac{P_{\mathscr S_0^\perp}D_r}
    {\|P_{\mathscr S_0^\perp}D_r\|_F}\right\rangle_F\\
  &=\frac{\|P_{\mathscr S_0^\perp}D_r\|_F^2}
    {\|P_{\mathscr S_0^\perp}D_r\|_F}\\
  &=\|P_{\mathscr S_0^\perp}D_r\|_F\\
  &\ge\sqrt{\frac r2}
   =\frac1{\sqrt2}\|D_r\|_F
   \ge\frac18\|D_r\|_F
   =\delta_0\|D_r\|_F.
\end{aligned}\tag{8}
\]
The second equality uses the orthogonal decomposition of \(D_r\), equation
(3) gives \(\|D_r\|_F=\sqrt r\), and
\(1/\sqrt2\ge1/8\). Thus equation (6) is exactly the witness condition in
the setting-defined event, proving equation (7). 

For every realized factor triple in \(\mathcal E_{\rm cond}\), the conditional estimate from Proposition~\ref{prop:step-008-raw-leverage} and the deterministic inclusion in
equation (7) give
\[
\begin{aligned}
  \mathbb P_{\rm init}(\mathcal E_{\rm deficit}^{c}\mid A,B,C)
  &\le
  \mathbb P_{\rm init}\!\left(
    \|P_{\mathscr S_0}D_r\|_F^2>\frac r2\,\middle|\,A,B,C
  \right)\\
  &\le r^{-20}.
\end{aligned}\tag{9}
\]
This spends no new failure probability.

\end{proof}

\begin{proof}[Proof of the subsection conclusion]

Lemma~\ref{lem:step-009-normal-energy} combines the exact membership
\(\widehat D_0\in\mathscr S_0\), the raw identity
\(\|D_r\|_F^2=r\), Pythagoras, and (LEV) to prove
\[
  \|P_{\mathscr S_0^\perp}D_r\|_F^2\ge r/2>0.
\]
This handles equality in (LEV), uses \(r\ge1\), and excludes the only
potential zero-denominator branch.

Proposition~\ref{prop:step-009-raw-witness} then normalizes that exact raw
normal component, proves unit norm and orthogonality, cancels
\(\widehat D_0\) by its exact same-span membership, and obtains
\[
  \langle D_r-\widehat D_0,W_0\rangle
  =\|P_{\mathscr S_0^\perp}D_r\|_F
  \ge\sqrt{r/2}
  \ge\delta_0\|D_r\|_F.
\]
Hence \((\mathrm{LEV})\Rightarrow\mathcal E_{\rm deficit}\) deterministically,
and equation (9) inherits the conditional \(r^{-20}\) failure bound
without a union bound or additional budget. This proves the claimed raw
tangent-deficit bound, with the raw target, raw tangent span, raw represented tensor, raw
witness, and Frobenius metric unchanged.

\end{proof}

\subsection{Initialization-event probability}\label{app:step-010}

\begin{proof}[Auxiliary facts in the present notation]


\medskip\noindent\textbf{Previously established event bounds.}\par

Proposition~\ref{prop:step-001-realized-conditioning} supplies the
unconditional smoothing bound
\[
  \mathbb P(\mathcal E_{\rm cond}^{\mathsf c})\le r^{-20}.
\]
Propositions~\ref{prop:step-003-normalized-gram-event},
\ref{prop:step-004-balanced-size-transfer}, and
\ref{prop:step-009-raw-witness} supply, uniformly for every realized triple
in \(\mathcal E_{\rm cond}\), the three respective conditional failure
bounds
\[
  \max\left\{
  \mathbb P((\mathcal E_{\rm gram}^{\rm norm})^{\mathsf c}\mid A,B,C),
  \mathbb P(\mathcal E_{\rm size}^{\mathsf c}\mid A,B,C),
  \mathbb P(\mathcal E_{\rm deficit}^{\mathsf c}\mid A,B,C)
  \right\}
  \le r^{-20}.
\]
These are statements under the same remaining-initialization conditional law.
Their uniformity on \(\mathcal E_{\rm cond}\) is exactly what is used below.


\medskip\noindent\textbf{Tower property and conditional union bound.}\par

We use the following statement in the present notation. If \(\mathcal G\) is a sub-sigma-field,
\(A_0\in\mathcal G\), and \(B_1,B_2,B_3\) are events, then
\[
  \mathbb P\!\left(A_0\cap\bigcup_{j=1}^3B_j\right)
  =\mathbb E\!\left[
    \mathbf 1_{A_0}\,\mathbb P\!\left(\bigcup_{j=1}^3B_j
      \middle|\mathcal G\right)
  \right]
  \le\mathbb E\!\left[
    \mathbf 1_{A_0}\sum_{j=1}^3\mathbb P(B_j\mid\mathcal G)
  \right].
\]

Take
\(\mathcal G=\sigma(A,B,C)\), \(A_0=\mathcal E_{\rm cond}\), and let the
three \(B_j\) be the complements of
\(\mathcal E_{\rm gram}^{\rm norm}\), \(\mathcal E_{\rm size}\), and
\(\mathcal E_{\rm deficit}\). The event \(\mathcal E_{\rm cond}\) depends
only on the realized factor triple, so it belongs to \(\mathcal G\). The
three displayed bounds control the conditional terms pointwise on
\(\mathcal E_{\rm cond}\). This argument forms no conditional probability by
dividing by \(\mathbb P(\mathcal E_{\rm cond})\), so no null-event issue
arises.

\end{proof}



\begin{proposition}[Conditional assembly of the initialization event]
\label{prop:step-010-conditional-union}
Under Assumptions~\ref{assump:base_conditioning},
\ref{assump:dimension}, \ref{assump:rank_window},
\ref{assump:gaussian_smoothing}, and
\ref{assump:independent_initialization}, and under
Propositions~\ref{prop:step-001-realized-conditioning},
\ref{prop:step-003-normalized-gram-event},
\ref{prop:step-004-balanced-size-transfer}, and
\ref{prop:step-009-raw-witness}, the exact setting event satisfies
\[
  \mathbb P(\mathcal E_{\rm init\_norm}^{\mathsf c})\le4r^{-20}
\]
under the joint smoothing/initialization law conditional on the deterministic
base triple.

\end{proposition}

\begin{proof}
By the exact event identity,
\[
  \mathcal E_{\rm init\_norm}
  =\mathcal E_{\rm cond}\cap\mathcal E_{\rm gram}^{\rm norm}
   \cap\mathcal E_{\rm size}\cap\mathcal E_{\rm deficit},
\]
and hence
\[
\begin{aligned}
  \mathcal E_{\rm init\_norm}^{\mathsf c}
  ={}&\mathcal E_{\rm cond}^{\mathsf c}\ \uplus\
  \Bigl[\mathcal E_{\rm cond}\cap
    \bigl((\mathcal E_{\rm gram}^{\rm norm})^{\mathsf c}
      \cup\mathcal E_{\rm size}^{\mathsf c}
      \cup\mathcal E_{\rm deficit}^{\mathsf c}\bigr)\Bigr],
\end{aligned}                                               \tag{1}
\]
where \(\uplus\) denotes a disjoint union. Let
\(\mathcal G=\sigma(A,B,C)\). Applying the tower property and the
conditional union bound to the second term in equation (1) gives
\[
\begin{aligned}
&\mathbb P\!\left(\mathcal E_{\rm cond}\cap
    \bigl((\mathcal E_{\rm gram}^{\rm norm})^{\mathsf c}
      \cup\mathcal E_{\rm size}^{\mathsf c}
      \cup\mathcal E_{\rm deficit}^{\mathsf c}\bigr)\right)\\
&\quad=\mathbb E\!\left[\mathbf 1_{\mathcal E_{\rm cond}}
  \mathbb P\!\left(
    (\mathcal E_{\rm gram}^{\rm norm})^{\mathsf c}
      \cup\mathcal E_{\rm size}^{\mathsf c}
      \cup\mathcal E_{\rm deficit}^{\mathsf c}
      \middle|\mathcal G\right)\right]\\
&\quad\le\mathbb E\!\left[\mathbf 1_{\mathcal E_{\rm cond}}
  \left(
    \mathbb P((\mathcal E_{\rm gram}^{\rm norm})^{\mathsf c}
      \mid\mathcal G)
    +\mathbb P(\mathcal E_{\rm size}^{\mathsf c}\mid\mathcal G)
    +\mathbb P(\mathcal E_{\rm deficit}^{\mathsf c}\mid\mathcal G)
  \right)\right]\\
&\quad\le 3r^{-20}.
\end{aligned}\tag{2}
\]
The last line uses the three bounds for every realized triple in
\(\mathcal E_{\rm cond}\). It does not require independence among the
Gram, size, and deficit events. Combining equations (1)--(2) with the
smoothing bound yields
\[
  \mathbb P(\mathcal E_{\rm init\_norm}^{\mathsf c})
  \le \mathbb P(\mathcal E_{\rm cond}^{\mathsf c})+3r^{-20}
  \le r^{-20}+3r^{-20}=4r^{-20}.
\]
This proves the proposition. 

\end{proof}


\begin{proposition}[Public initialization confidence]
\label{prop:step-010-public-confidence}
Under Assumption~\ref{assump:dimension} and
Proposition~\ref{prop:step-010-conditional-union}, if \(r\ge2\), then
\[
  \mathbb P(\mathcal E_{\rm init\_norm})\ge1-r^{-10}.
\]

\end{proposition}

\begin{proof}
For \(r\ge2\),
\[
  4\le2^{10}\le r^{10},
\]
and therefore
\[
  4r^{-20}\le r^{10}r^{-20}=r^{-10}.                     \tag{3}
\]
Proposition~\ref{prop:step-010-conditional-union} and equation (3) imply
\[
\begin{aligned}
  \mathbb P(\mathcal E_{\rm init\_norm})
  &=1-\mathbb P(\mathcal E_{\rm init\_norm}^{\mathsf c})\\
  &\ge1-4r^{-20}\\
  &\ge1-r^{-10}.
\end{aligned}
\]
Assumption~\ref{assump:dimension} already permits enlarging the theorem's
large-\(r\) threshold, so requiring \(r\ge2\) adds no new theorem-facing
condition. 

\end{proof}

\begin{proof}[Proof of the subsection conclusion]

Proposition~\ref{prop:step-010-conditional-union} uses the exact event
identity, the unconditional failure for \(\mathcal E_{\rm cond}\),
and the three uniformly conditional failure bounds. The tower
property and a conditional union bound give
\[
  \mathbb P(\mathcal E_{\rm init\_norm}^{\mathsf c})
  \le r^{-20}+3r^{-20}=4r^{-20}
\]
without multiplying probabilities or asserting independence among the three
initialization events. Proposition~\ref{prop:step-010-public-confidence}
then proves explicitly that \(4r^{-20}\le r^{-10}\) for \(r\ge2\), and
therefore
\[
  \boxed{\mathbb P(\mathcal E_{\rm init\_norm})\ge1-r^{-10}}.
\]
This proves the claim under the exact joint probability space.

\end{proof}

\subsection{Finite-path convergence and factor radius}\label{app:step-011}

\begin{proof}[Auxiliary facts in the present notation]


\medskip\noindent\textbf{balanced initial-size event.}\par

Proposition~\ref{prop:step-004-balanced-size-transfer} proves the generated
event \(\mathcal E_{\rm size}\) for the actual balanced initialization. By
its setting definition, on this event
\[
  \|m_{i,0}\|_2\le2
  \quad\text{for every }i\in[k],\quad m\in\{x,y,z\}.
\]

The columns used below are exactly the columns of \(X_0,Y_0,Z_0\) from the
current setting. No raw initialization column or balancing scalar is
introduced in this argument.


\medskip\noindent\textbf{exact initialization event.}\par

Propositions~\ref{prop:step-010-conditional-union} and
\ref{prop:step-010-public-confidence} use the exact setting identity
\[
  \mathcal E_{\rm init\_norm}
  =\mathcal E_{\rm cond}\cap\mathcal E_{\rm gram}^{\rm norm}
   \cap\mathcal E_{\rm size}\cap\mathcal E_{\rm deficit}
\]
and establish \(\mathbb P(\mathcal E_{\rm init\_norm})\ge1-r^{-10}\).

The deterministic implication used
here is only
\(
  \mathcal E_{\rm init\_norm}\subseteq\mathcal E_{\rm size}
\),
with the exact events from the current setting. This inclusion does not
convert initialization confidence into a trajectory property.


\medskip\noindent\textbf{Finite-dimensional Euclidean completeness.}\par

We use the following statement in the present notation. Let \(\mathcal H\) be a finite-dimensional real Euclidean
space with norm \(\|\cdot\|_{\mathcal H}\). If for every \(\varepsilon>0\)
there is \(N\) such that
\(\|h_t-h_s\|_{\mathcal H}<\varepsilon\) for all \(s,t\ge N\), then there is
a finite \(h_\infty\in\mathcal H\) with
\(\|h_t-h_\infty\|_{\mathcal H}\to0\).

Here
\[
  \mathcal H=(\mathbb R^{n\times k})^3,
  \qquad
  \|(X,Y,Z)\|_{\mathcal H}
  =\bigl(\|X\|_F^2+\|Y\|_F^2+\|Z\|_F^2\bigr)^{1/2},
\]
so the setting metric satisfies
\(d_{\rm bal}(\theta,\theta')=\|\theta-\theta'\|_{\mathcal H}\).
Lemma~\ref{lem:step-011-finite-path-limit} verifies the Cauchy hypothesis from
the tail of the setting-defined nonnegative path series. This is a standard
finite-dimensional fact in exactly the current norm; no quotient-space or
additional completeness assumption is used.

\end{proof}



\begin{lemma}[finite total variation gives a finite balanced-factor limit]
\label{lem:step-011-finite-path-limit}
Under Assumption~\ref{assump:gd_step}, consider the setting-defined trajectory
\((\theta_t)_{t\ge0}\). On the sole local conditional hypothesis
\(\mathcal C_{\rm path}\), define
\[
  \ell_t:=d_{\rm bal}(\theta_{t+1},\theta_t)\ge0.
\]
Then, for every pair of integers \(0\le s<t\),
\[
  d_{\rm bal}(\theta_s,\theta_t)
  \le \sum_{u=s}^{t-1}\ell_u.                         \tag{1}
\]
The tails \(\sum_{u=s}^{\infty}\ell_u\) tend to zero as
\(s\to\infty\). Consequently, the balanced representatives form a Cauchy
sequence in \((\mathbb R^{n\times k})^3\) and converge in
\(d_{\rm bal}\) to a finite \(\theta_\infty\).

\end{lemma}

\begin{proof}
Assumption~\ref{assump:gd_step} is used only to
identify the recursively defined trajectory. The metric formula in the
setting is the Euclidean product-space distance, so repeated application of
its triangle inequality gives
\[
\begin{aligned}
  d_{\rm bal}(\theta_s,\theta_t)
  &\le d_{\rm bal}(\theta_s,\theta_{s+1})
      +d_{\rm bal}(\theta_{s+1},\theta_{s+2})
      +\cdots+d_{\rm bal}(\theta_{t-1},\theta_t)\\
  &=\sum_{u=s}^{t-1}\ell_u,
\end{aligned}
\]
which proves equation (1).

Let
\[
  S_N:=\sum_{u=0}^{N-1}\ell_u,
  \qquad S_0:=0.
\]
On \(\mathcal C_{\rm path}\), the nonnegative partial sums increase to the
finite value
\[
  E_{\rm path}=\sum_{u=0}^{\infty}\ell_u\le E_\star<\infty.
\]
Therefore the exact tail remainder
\[
  R_s:=E_{\rm path}-S_s=\sum_{u=s}^{\infty}\ell_u
\]
is nonnegative and satisfies \(R_s\downarrow0\). Given
\(\varepsilon>0\), choose \(N\) with \(R_N<\varepsilon\). For arbitrary
\(s,t\ge N\), assume without loss of generality that \(s<t\). Equation (1)
and nonnegativity give
\[
  d_{\rm bal}(\theta_s,\theta_t)
  \le\sum_{u=s}^{t-1}\ell_u
  \le R_s\le R_N<\varepsilon.
\]
Thus \((\theta_t)\) is Cauchy in the finite-dimensional Euclidean product
space
\(
  (\mathbb R^{n\times k})^3
\),
which is complete. Hence there is a finite
\(
  \theta_\infty=(X_\infty,Y_\infty,Z_\infty)
\)
such that
\[
  d_{\rm bal}(\theta_t,\theta_\infty)\longrightarrow0.
\]

If \(E_{\rm path}=0\), nonnegativity forces \(\ell_t=0\) for every \(t\).
Because \(d_{\rm bal}\) is a metric on the displayed product space,
\(\theta_{t+1}=\theta_t\) for every \(t\), so the trajectory is constant and
\(\theta_\infty=\theta_0\). At the other boundary
\(E_{\rm path}=E_\star\), the same finite-tail argument applies with equality
in the total budget and requires no strict slack. 

\end{proof}


\begin{proposition}[horizon-uniform path and endpoint radius]
\label{prop:step-011-path-radius}
\leavevmode\par\noindent
Under Assumption~\ref{assump:gd_step},
Lemma~\ref{lem:step-011-finite-path-limit}, and the sole local conditional
hypothesis \(\mathcal C_{\rm path}\), every \(t\ge0\) satisfies
\[
  d_{\rm bal}(\theta_t,\theta_0)
  \le\sum_{u=0}^{t-1}\ell_u
  \le E_{\rm path}
  \le E_\star
  \le1,                                                  \tag{2}
\]
where the first sum is empty when \(t=0\). The finite limit produced by
Lemma~\ref{lem:step-011-finite-path-limit} satisfies
\[
  d_{\rm bal}(\theta_\infty,\theta_0)
  \le E_{\rm path}
  \le E_\star.                                          \tag{3}
\]

\end{proposition}

\begin{proof}
Applying equation (1) from
Lemma~\ref{lem:step-011-finite-path-limit} with \(s=0\) gives
\[
  d_{\rm bal}(\theta_t,\theta_0)
  =d_{\rm bal}(\theta_0,\theta_t)
  \le\sum_{u=0}^{t-1}\ell_u
  \le\sum_{u=0}^{\infty}\ell_u
  =E_{\rm path}.
\]
The event \(\mathcal C_{\rm path}\) gives
\(E_{\rm path}\le E_\star\). From the setting definition
\[
  E_\star
  =\min\left\{1,
    \sqrt{\frac{\delta_0}{16C_{\rm CP}(\kappa,3)}}\right\},
\]
we also have \(E_\star\le1\). This proves equation (2).

Since \(\theta_t\to\theta_\infty\) in \(d_{\rm bal}\), continuity of the
Euclidean norm, equivalently the reverse triangle inequality, gives
\[
  d_{\rm bal}(\theta_\infty,\theta_0)
  =\lim_{t\to\infty}d_{\rm bal}(\theta_t,\theta_0).
\]
Taking the limit in the already uniform bound
\(d_{\rm bal}(\theta_t,\theta_0)\le E_{\rm path}\) proves equation (3).
Both conclusions remain non-strict when \(E_{\rm path}=E_\star\), while the
case \(E_{\rm path}=0\) reduces to the constant trajectory identified in
Lemma~\ref{lem:step-011-finite-path-limit}. 

\end{proof}


\begin{proposition}[uniform component bounds along the path and at the limit]
\label{prop:step-011-uniform-factor-radius}
Under Assumption~\ref{assump:gd_step}, Proposition~\ref{prop:step-004-balanced-size-transfer}, Propositions~\ref{prop:step-010-conditional-union} and
\ref{prop:step-010-public-confidence}, and
Proposition~\ref{prop:step-011-path-radius}, on
\(\mathcal E_{\rm init\_norm}\cap\mathcal C_{\rm path}\), for every
component \(i\in[k]\), every mode
\(M_t\in\{X_t,Y_t,Z_t\}\), and every \(t\ge0\),
\[
  \|m_{i,t}\|_2\le2+E_\star\le3.                       \tag{4}
\]
Writing \(M_\infty\in\{X_\infty,Y_\infty,Z_\infty\}\) for the
corresponding limit block and \(m_{i,\infty}\) for its \(i\)-th column,
\[
  \|m_{i,\infty}\|_2\le2+E_\star\le3.                 \tag{5}
\]

\end{proposition}

\begin{proof}
By the exact setting identity for \(\mathcal E_{\rm init\_norm}\),
\(\mathcal E_{\rm size}\) holds on
\(\mathcal E_{\rm init\_norm}\). Hence
\[
  \|m_{i,0}\|_2\le2
\]
for every component and mode.

Fix \(i\), one mode matrix \(M_t\), and \(t\ge0\). The Euclidean norm of one
column is bounded by the Frobenius norm of the full mode-matrix difference,
and that Frobenius norm is one summand of the product norm defining
\(d_{\rm bal}\). Therefore
\[
\begin{aligned}
  \|m_{i,t}\|_2
  &\le \|m_{i,0}\|_2+\|m_{i,t}-m_{i,0}\|_2\\
  &\le \|m_{i,0}\|_2+\|M_t-M_0\|_F\\
  &\le \|m_{i,0}\|_2+d_{\rm bal}(\theta_t,\theta_0)\\
  &\le2+E_\star\le3,
\end{aligned}
\]
where Proposition~\ref{prop:step-011-path-radius} supplies the fourth line and
the setting definition supplies \(E_\star\le1\). This proves equation (4)
simultaneously for every finite time.

The product-space convergence from
Lemma~\ref{lem:step-011-finite-path-limit} implies
\(M_t\to M_\infty\) in Frobenius norm and hence identifies
\(m_{i,\infty}\) as the Euclidean limit of the corresponding columns.
Using the endpoint bound from Proposition~\ref{prop:step-011-path-radius}
directly gives
\[
\begin{aligned}
  \|m_{i,\infty}\|_2
  &\le\|m_{i,0}\|_2+\|m_{i,\infty}-m_{i,0}\|_2\\
  &\le\|m_{i,0}\|_2+\|M_\infty-M_0\|_F\\
  &\le\|m_{i,0}\|_2+d_{\rm bal}(\theta_\infty,\theta_0)\\
  &\le2+E_\star\le3,
\end{aligned}
\]
which proves equation (5). At \(E_{\rm path}=0\), all inequalities reduce to
the unchanged initial columns and \(\theta_\infty=\theta_0\). At
\(E_{\rm path}=E_\star\), the non-strict bound still gives exactly
\(2+E_\star\le3\). 

\end{proof}

\begin{proof}[Proof of the subsection conclusion]

Work on the exact event
\(\mathcal E_{\rm init\_norm}\cap\mathcal C_{\rm path}\).
Assumption~\ref{assump:gd_step} identifies the trajectory but supplies no
convergence claim. Lemma~\ref{lem:step-011-finite-path-limit} uses only the
sole local conditional hypothesis \(\mathcal C_{\rm path}\) to prove the exact
tail-sum bound
\[
  d_{\rm bal}(\theta_s,\theta_t)
  \le\sum_{u=s}^{t-1}\ell_u,
\]
shows that the tails vanish, and obtains a finite limit
\(\theta_\infty\) by completeness of the fixed Euclidean product space.
Proposition~\ref{prop:step-011-path-radius} then proves the horizon-uniform and
endpoint bounds
\[
  d_{\rm bal}(\theta_t,\theta_0)\le E_\star,
  \qquad
  d_{\rm bal}(\theta_\infty,\theta_0)\le E_\star.
\]

The event identity and Proposition~\ref{prop:step-004-balanced-size-transfer}
give \(\mathcal E_{\rm size}\) on \(\mathcal E_{\rm init\_norm}\), and hence
the initial column bound \(2\). Proposition~\ref{prop:step-011-uniform-factor-radius}
combines that bound with the exact mode-matrix displacement dominated by
\(d_{\rm bal}\) to prove, for every component, mode, finite time, and the
limit,
\[
  \|m_{i,t}\|_2\le2+E_\star\le3,
  \qquad
  \|m_{i,\infty}\|_2\le2+E_\star\le3.
\]
Thus the exact conditional path-convergence and radius-\(3\) claim is proved
without upgrading \(\mathcal C_{\rm path}\) to an unconditional theorem and
without inserting convergence, trapping, or loss into the certificate.

\end{proof}

\subsection{Multilinear endpoint remainder}\label{app:step-012}

\begin{proof}[Auxiliary facts in the present notation]


\medskip\noindent\textbf{realized-factor conditioning.}\par

Proposition~\ref{prop:step-001-realized-conditioning} proves that on the
setting-defined event \(\mathcal E_{\rm cond}\), the realized matrices
\(A,B,C\) have full column rank and
\[
  \|M^\dagger\|_{\rm op}\le\kappa_1
  \quad\text{for }M\in\{A,B,C\},
  \qquad \kappa_1=2\kappa^2.
\]

These are exactly the three left
inverses in the setting map \(\Psi_{A,B,C}\). No inverse in a different
coordinate system and no surrogate coefficient map is used.


\medskip\noindent\textbf{finite endpoint, displacement, and radius.}\par

Lemma~\ref{lem:step-011-finite-path-limit} and
Propositions~\ref{prop:step-011-path-radius} and
\ref{prop:step-011-uniform-factor-radius} prove that on
\(\mathcal E_{\rm init\_norm}\cap\mathcal C_{\rm path}\), the actual
balanced representatives converge to a finite \(\theta_\infty\), and
\[
  d_{\rm bal}(\theta_\infty,\theta_0)
  \le E_{\rm path}\le E_\star\le1,
  \qquad
  \max_{i,m}\|m_{i,0}\|_2\le3.
\]
The result actually gives the sharper initial radius \(2\) and
radius \(3\) at every time and at the endpoint; the present step uses only
the displayed base radius \(3\).

The base point is exactly
\(\theta_0\), the comparison point is exactly \(\theta_\infty\), and the
metric is the setting's product Frobenius metric \(d_{\rm bal}\). Because
\(\mathcal E_{\rm init\_norm}\subseteq\mathcal E_{\rm cond}\), the
left-inverse bounds hold on the same event.


\medskip\noindent\textbf{Finite-sum Cauchy--Schwarz and rank-one Frobenius identity.}\par

We use the following statements in the present notation. For nonnegative finite sequences \((a_i)\) and
\((b_i)\),
\[
  \sum_i a_i b_i
  \le\left(\sum_i a_i^2\right)^{1/2}
      \left(\sum_i b_i^2\right)^{1/2}.
\]
For Euclidean vectors \(u,v,w\), direct expansion of the entries gives
\[
  \|u\otimes v\otimes w\|_F^2
  =\sum_{a,b,c}u_a^2v_b^2w_c^2
  =\|u\|_2^2\|v\|_2^2\|w\|_2^2.
\]

The sequences are column norms of
the physical increments \(\Delta X,\Delta Y,\Delta Z\), and the rank-one
tensors are the four exact remainder families in
Lemma~\ref{lem:step-012-exact-expansion}. All sums are finite over
\(i\in[k]\).

\end{proof}



\begin{lemma}[exact trilinear coefficient expansion]
\label{lem:step-012-exact-expansion}
Under Proposition~\ref{prop:step-001-realized-conditioning}, on
\(\mathcal E_{\rm cond}\), let
\(\theta=(X,Y,Z)\) and \(\theta'=(X',Y',Z')\), and write
\[
  \Delta X=X'-X,\qquad \Delta Y=Y'-Y,\qquad \Delta Z=Z'-Z.
\]
For every \(i\in[k]\), define the exact raw coefficient columns
\[
  \alpha_i=A^\dagger x_i,\quad
  \beta_i=B^\dagger y_i,\quad
  \gamma_i=C^\dagger z_i
\]
and their increments
\[
  \Delta\alpha_i=A^\dagger(x_i'-x_i),\quad
  \Delta\beta_i=B^\dagger(y_i'-y_i),\quad
  \Delta\gamma_i=C^\dagger(z_i'-z_i).
\]
Then
\[
\begin{aligned}
  D\Psi_{A,B,C}(\theta)[\theta'-\theta]
  =\sum_{i=1}^k\bigl(&
      \Delta\alpha_i\otimes\beta_i\otimes\gamma_i
      +\alpha_i\otimes\Delta\beta_i\otimes\gamma_i\\
      &+\alpha_i\otimes\beta_i\otimes\Delta\gamma_i
    \bigr),
\end{aligned}\tag{1}
\]
and, after subtracting the value and all three derivative terms per
component, the remainder is exactly
\[
\begin{aligned}
  &\Psi_{A,B,C}(\theta')-\Psi_{A,B,C}(\theta)
       -D\Psi_{A,B,C}(\theta)[\theta'-\theta]\\
  &\quad=\sum_{i=1}^k\bigl(
      \Delta\alpha_i\otimes\Delta\beta_i\otimes\gamma_i
      +\Delta\alpha_i\otimes\beta_i\otimes\Delta\gamma_i
      +\alpha_i\otimes\Delta\beta_i\otimes\Delta\gamma_i
      +\Delta\alpha_i\otimes\Delta\beta_i\otimes
        \Delta\gamma_i
    \bigr).
\end{aligned}\tag{2}
\]
Thus equation (2) contains exactly three quadratic tensors and one cubic
tensor for each component, and contains no constant or first-order term.

\end{lemma}

\begin{proof}
Linearity of the fixed left inverses gives
\[
  A^\dagger x_i'=\alpha_i+\Delta\alpha_i,
  \quad B^\dagger y_i'=\beta_i+\Delta\beta_i,
  \quad C^\dagger z_i'=\gamma_i+\Delta\gamma_i.
\]
Substitution into the exact setting map yields
\[
  \Psi_{A,B,C}(\theta')
  =\sum_{i=1}^k
    (\alpha_i+\Delta\alpha_i)
    \otimes(\beta_i+\Delta\beta_i)
    \otimes(\gamma_i+\Delta\gamma_i).
\]
Expanding each trilinear product gives eight tensors: the base tensor, the
three tensors linear in one increment, the three tensors quadratic in two
increments, and the tensor cubic in all three increments. The coefficient
of a scalar perturbation parameter in the three linear tensors is precisely
the Frechet derivative in equation (1). Subtracting the base value and
equation (1) leaves exactly equation (2). This is an algebraic identity, with
no limiting argument, generic smoothness estimate, or remainder notation.


\end{proof}


\begin{proposition}[quadratic raw Taylor bound from a base-column radius]
\label{prop:step-012-raw-remainder}
Under Proposition~\ref{prop:step-001-realized-conditioning} and
Lemma~\ref{lem:step-012-exact-expansion}, work on
\(\mathcal E_{\rm cond}\). Let \(R\ge0\) and suppose only that the columns
of the base point \(\theta=(X,Y,Z)\) satisfy
\[
  \max_{i\in[k]}
  \max\{\|x_i\|_2,\|y_i\|_2,\|z_i\|_2\}\le R.             \tag{3}
\]
For any \(\theta'=(X',Y',Z')\), put
\[
  d=d_{\rm bal}(\theta',\theta)
   =\bigl(\|\Delta X\|_F^2+\|\Delta Y\|_F^2
           +\|\Delta Z\|_F^2\bigr)^{1/2},
\]
and assume \(0\le d\le1\). Then the three quadratic families in equation
(2) each have Frobenius norm at most \(\kappa_1^3Rd^2\), the cubic family
has Frobenius norm at most \(\kappa_1^3d^3\le\kappa_1^3d^2\), and
\[
  \left\|\Psi_{A,B,C}(\theta')-\Psi_{A,B,C}(\theta)
  -D\Psi_{A,B,C}(\theta)[\theta'-\theta]\right\|_F
  \le \kappa_1^3(1+3R)d^2
  =C_{\rm CP}(\kappa,R)d^2.                               \tag{4}
\]
No column-radius condition on \(\theta'\) is required.

\end{proposition}

\begin{proof}
On \(\mathcal E_{\rm cond}\), Proposition~\ref{prop:step-001-realized-conditioning} gives
\[
  \|A^\dagger\|_{\rm op},\ \|B^\dagger\|_{\rm op},\
  \ \|C^\dagger\|_{\rm op}\le\kappa_1.                   \tag{5}
\]
Consequently, for every component,
\[
\begin{aligned}
  \|\Delta\alpha_i\|_2&\le\kappa_1\|\Delta x_i\|_2,
  &\|\alpha_i\|_2&\le\kappa_1R,\\
  \|\Delta\beta_i\|_2&\le\kappa_1\|\Delta y_i\|_2,
  &\|\beta_i\|_2&\le\kappa_1R,\\
  \|\Delta\gamma_i\|_2&\le\kappa_1\|\Delta z_i\|_2,
  &\|\gamma_i\|_2&\le\kappa_1R.
\end{aligned}\tag{6}
\]
Denote the three quadratic sums in their displayed order in equation (2) by
\(Q_{xy},Q_{xz},Q_{yz}\), and denote the cubic sum by \(C_{xyz}\). The
triangle inequality, the rank-one Frobenius identity, equation (6), and
Cauchy--Schwarz give each quadratic family explicitly:
\begin{align}
  \|Q_{xy}\|_F
  &\le\sum_{i=1}^k
       \|\Delta\alpha_i\|_2\|\Delta\beta_i\|_2\|\gamma_i\|_2\notag\\
  &\le\kappa_1^3R
       \sum_{i=1}^k\|\Delta x_i\|_2\|\Delta y_i\|_2
   \le\kappa_1^3R\|\Delta X\|_F\|\Delta Y\|_F
   \le\kappa_1^3Rd^2,                                    \tag{7}\\
  \|Q_{xz}\|_F
  &\le\sum_{i=1}^k
       \|\Delta\alpha_i\|_2\|\beta_i\|_2\|\Delta\gamma_i\|_2\notag\\
  &\le\kappa_1^3R
       \sum_{i=1}^k\|\Delta x_i\|_2\|\Delta z_i\|_2
   \le\kappa_1^3R\|\Delta X\|_F\|\Delta Z\|_F
   \le\kappa_1^3Rd^2,                                    \tag{8}\\
  \|Q_{yz}\|_F
  &\le\sum_{i=1}^k
       \|\alpha_i\|_2\|\Delta\beta_i\|_2\|\Delta\gamma_i\|_2\notag\\
  &\le\kappa_1^3R
       \sum_{i=1}^k\|\Delta y_i\|_2\|\Delta z_i\|_2
   \le\kappa_1^3R\|\Delta Y\|_F\|\Delta Z\|_F
   \le\kappa_1^3Rd^2.                                    \tag{9}
\end{align}
For the cubic family, let
\(a_i=\|\Delta x_i\|_2\),
\(b_i=\|\Delta y_i\|_2\), and
\(c_i=\|\Delta z_i\|_2\). Two finite-sum inequalities give
\[
\begin{aligned}
  \sum_{i=1}^k a_ib_ic_i
  &\le\left(\sum_i a_i^2\right)^{1/2}
       \left(\sum_i b_i^2c_i^2\right)^{1/2}\\
  &\le\left(\sum_i a_i^2\right)^{1/2}
       \left(\sum_i b_i^2\right)^{1/2}
       \left(\sum_i c_i^2\right)^{1/2},
\end{aligned}\tag{10}
\]
where the second line also follows directly from
\(\sum_i b_i^2c_i^2\le(\sum_i b_i^2)(\sum_i c_i^2)\).
Therefore
\[
\begin{aligned}
  \|C_{xyz}\|_F
  &\le\kappa_1^3\sum_{i=1}^k a_ib_ic_i\\
  &\le\kappa_1^3
       \|\Delta X\|_F\|\Delta Y\|_F\|\Delta Z\|_F
   \le\kappa_1^3d^3
   \le\kappa_1^3d^2,
\end{aligned}\tag{11}
\]
where the last inequality uses exactly \(d\le1\). Applying the triangle
inequality to the four exact families in equation (2) and using
equations (7)--(11) yields
\[
  3\kappa_1^3Rd^2+\kappa_1^3d^2
  =\kappa_1^3(1+3R)d^2,
\]
which is equation (4).

Only the unperturbed coefficients \(\alpha_i,\beta_i,\gamma_i\) occur in
the three quadratic terms, one per term. The comparison point
\(\theta'\) enters exclusively through \(\Delta X,\Delta Y,\Delta Z\).
This is why equation (3) is needed only at the base point and no radius for
\(\theta'\) is used.

If \(d=0\), all three physical increments and hence all coefficient
increments vanish, so equation (2) is exactly zero and equation (4) holds
with equality. If \(d=1\), the only degree reduction used above is
\(d^3\le d^2\), which is equality at that boundary; all other estimates
remain non-strict. Thus both boundary values are included. 

\end{proof}


\begin{proposition}[one-endpoint raw coefficient remainder]
\label{prop:step-012-endpoint-remainder}
Under Proposition~\ref{prop:step-001-realized-conditioning},
Lemma~\ref{lem:step-011-finite-path-limit} and
Propositions~\ref{prop:step-011-path-radius} and
\ref{prop:step-011-uniform-factor-radius}, and
Proposition~\ref{prop:step-012-raw-remainder}, on the exact conditional event
\(\mathcal E_{\rm init\_norm}\cap\mathcal C_{\rm path}\),
\[
\begin{aligned}
  &\left\|\Psi_{A,B,C}(\theta_\infty)
      -\Psi_{A,B,C}(\theta_0)
      -D\Psi_{A,B,C}(\theta_0)[\theta_\infty-\theta_0]
    \right\|_F\\
  &\qquad\le C_{\rm CP}(\kappa,3)
      d_{\rm bal}(\theta_\infty,\theta_0)^2
   \le C_{\rm CP}(\kappa,3)E_\star^2
   \le\frac{\delta_0}{16}.
\end{aligned}\tag{12}
\]
This is a raw coefficient-space Frobenius bound for the single endpoint
pair; it makes no physical-coordinate change and contains no sum over time.

\end{proposition}

\begin{proof}
On
\(\mathcal E_{\rm init\_norm}\cap\mathcal C_{\rm path}\), Proposition~\ref{prop:step-011-uniform-factor-radius} gives
\[
  \max_{i\in[k]}
  \max\{\|x_{i,0}\|_2,\|y_{i,0}\|_2,\|z_{i,0}\|_2\}
  \le3.                                                    \tag{13}
\]
This is the required base-point condition with \(R=3\). No endpoint column
bound is needed for Proposition~\ref{prop:step-012-raw-remainder}.
Proposition~\ref{prop:step-011-path-radius} gives
\[
  d_{\rm bal}(\theta_\infty,\theta_0)
  \le E_{\rm path}\le E_\star\le1.                       \tag{14}
\]
Moreover, the definition of \(\mathcal E_{\rm init\_norm}\) includes
\(\mathcal E_{\rm cond}\), so the left-inverse bounds hold for the
same realization. Proposition~\ref{prop:step-012-raw-remainder}, with
\(\theta=\theta_0\), \(\theta'=\theta_\infty\), \(R=3\), and equation
(14), gives the first two inequalities in equation (12). Finally, the exact
setting definition
\[
  E_\star
  =\min\left\{1,
       \sqrt{\frac{\delta_0}{16C_{\rm CP}(\kappa,3)}}
    \right\}
\]
implies
\[
  C_{\rm CP}(\kappa,3)E_\star^2\le\frac{\delta_0}{16},
\]
which proves the last inequality.

At the zero-path boundary, Proposition~\ref{prop:step-011-uniform-factor-radius} gives
\(\theta_\infty=\theta_0\), and every term on the left of equation (12)
vanishes exactly. At the unit-displacement boundary, the generic estimate
already includes \(d=1\); no strict radius or displacement inequality is
used. 

\end{proof}

\begin{proof}[Proof of the subsection conclusion]

Lemma~\ref{lem:step-012-exact-expansion} expands the exact setting map
\(\Psi_{A,B,C}\) in the raw coefficient coordinates induced by the realized
left inverses. After removing \(\Psi_{A,B,C}(\theta)\) and the exact Frechet
derivative, it leaves exactly three quadratic tensors and one cubic tensor
per component. Proposition~\ref{prop:step-012-raw-remainder} applies the
\(\kappa_1\) left-inverse bounds, the base physical column radius,
the rank-one Frobenius identity, and finite-sum Cauchy--Schwarz. Each
quadratic family is bounded by \(\kappa_1^3Rd^2\), while the cubic family is
bounded by \(\kappa_1^3d^3\le\kappa_1^3d^2\). Their sum is exactly
\[
  C_{\rm CP}(\kappa,R)d^2
  =\kappa_1^3(1+3R)d^2.
\]

On \(\mathcal E_{\rm init\_norm}\cap\mathcal C_{\rm path}\),
Proposition~\ref{prop:step-011-uniform-factor-radius} supplies the finite
endpoint, base radius \(R=3\), and
\(d_{\rm bal}(\theta_\infty,\theta_0)\le E_\star\le1\). Since this event is
contained in \(\mathcal E_{\rm cond}\),
Proposition~\ref{prop:step-001-realized-conditioning} supplies the required
left-inverse bounds. Proposition~\ref{prop:step-012-endpoint-remainder}
therefore establishes the exact one-endpoint raw coefficient remainder
\[
  \left\|\Psi_{A,B,C}(\theta_\infty)-\Psi_{A,B,C}(\theta_0)
      -D\Psi_{A,B,C}(\theta_0)[\theta_\infty-\theta_0]\right\|_F
  \le C_{\rm CP}(\kappa,3)E_\star^2
  \le\delta_0/16.
\]
This is the claimed Taylor constant. It is a single endpoint estimate, not a
time accumulation, and it remains in the exact raw coefficient Frobenius
norm.

\end{proof}

\subsection{Persistence of the raw coefficient margin}\label{app:step-013}

\begin{proof}[Auxiliary facts in the present notation]


\medskip\noindent\textbf{raw normal witness and initial margin.}\par

Proposition~\ref{prop:step-009-raw-witness} supplies, on the generated deficit
event contained in \(\mathcal E_{\rm init\_norm}\), the exact raw witness
\(W_0\) satisfying
\[
  \|W_0\|_F=1,
  \qquad W_0\perp\mathscr S_0,
  \qquad
  \langle D_r-\widehat D_0,W_0\rangle_F
  \ge\delta_0\sqrt r.
  \tag{1}
\]

These are the unnormalized target,
tangent span, represented initial coefficient tensor, witness, and Frobenius
inner product from the setting. The exact setting identity
\[
  \Psi_{A,B,C}(\theta_0)
  =\sum_{i=1}^k\alpha_{i,0}\otimes\beta_{i,0}\otimes\gamma_{i,0}
  =\widehat D_0                                             \tag{2}
\]
translates equation (1) directly to the base point of the endpoint expansion.


\medskip\noindent\textbf{finite path endpoint.}\par

Lemma~\ref{lem:step-011-finite-path-limit} and
Proposition~\ref{prop:step-011-path-radius} prove, on
\(\mathcal E_{\rm init\_norm}\cap\mathcal C_{\rm path}\), that the actual
balanced representatives converge to a finite \(\theta_\infty\) and that
\(d_{\rm bal}(\theta_\infty,\theta_0)\le E_\star\). If
\(E_{\rm path}=0\), then \(\theta_\infty=\theta_0\).

The endpoint and base point below are
exactly this \(\theta_\infty\) and the setting initialization \(\theta_0\).
No finite-time sequence or additional limit is introduced.


\medskip\noindent\textbf{derivative formula and endpoint remainder.}\par

Lemma~\ref{lem:step-012-exact-expansion} gives the exact Frechet derivative of
the raw coefficient map. Proposition~\ref{prop:step-012-endpoint-remainder} gives, for the same
endpoint pair on the same conditional event,
\[
  \left\|\Psi_{A,B,C}(\theta_\infty)-\Psi_{A,B,C}(\theta_0)
  -D\Psi_{A,B,C}(\theta_0)[\theta_\infty-\theta_0]\right\|_F
  \le\frac{\delta_0}{16}.                                \tag{3}
\]

Equation (3) is already in the exact
raw coefficient Frobenius norm paired with \(W_0\). It is one endpoint
estimate and has no time sum. Proposition~\ref{prop:step-012-endpoint-remainder}
verifies its base-radius, displacement, conditioning, and \(E_\star\)
conditions.


\medskip\noindent\textbf{Frobenius Cauchy--Schwarz.}\par

We use the following statement in the present notation. For coefficient tensors \(U,V\),
\[
  |\langle U,V\rangle_F|\le\|U\|_F\|V\|_F.
  \tag{4}
\]

Equation (4) is applied first with
\(U=R_\infty\), \(V=W_0\), and then with
\(U=D_r-\Psi_{A,B,C}(\theta_\infty)\), \(V=W_0\). Both tensors lie in the
same finite-dimensional raw coefficient Frobenius space, and equation (1)
gives \(\|W_0\|_F=1\).

\end{proof}



\begin{lemma}[raw tangent-range cancellation at initialization]
\label{lem:step-013-tangent-cancellation}
Under the exact setting definitions and
Propositions~\ref{prop:step-009-raw-witness},
\ref{prop:step-011-uniform-factor-radius}, and
\ref{prop:step-012-endpoint-remainder}, work on
\(\mathcal E_{\rm init\_norm}\cap\mathcal C_{\rm path}\). Let
\[
  \Delta\theta=\theta_\infty-\theta_0
  =(\Delta X,\Delta Y,\Delta Z).
\]
Then
\[
  D\Psi_{A,B,C}(\theta_0)[\Delta\theta]\in\mathscr S_0,
  \qquad
  \left\langle D\Psi_{A,B,C}(\theta_0)[\Delta\theta],W_0\right\rangle_F=0.
  \tag{5}
\]

\end{lemma}

\begin{proof}
Write \(\Delta x_i,\Delta y_i,\Delta z_i\) for
the columns of the three increment matrices. The exact derivative formula
from Lemma~\ref{lem:step-012-exact-expansion}, specialized at
\(\theta_0\), is
\[
\begin{aligned}
  D\Psi_{A,B,C}(\theta_0)[\Delta\theta]
  =\sum_{i=1}^k\bigl(&
    (A^\dagger\Delta x_i)\otimes\beta_{i,0}\otimes\gamma_{i,0}\\
    &+\alpha_{i,0}\otimes(B^\dagger\Delta y_i)\otimes\gamma_{i,0}\\
    &+\alpha_{i,0}\otimes\beta_{i,0}\otimes(C^\dagger\Delta z_i)
  \bigr).
\end{aligned}\tag{6}
\]
Each of \(A^\dagger\Delta x_i\), \(B^\dagger\Delta y_i\), and
\(C^\dagger\Delta z_i\) lies in \(\mathbb R^r\). By the exact raw tangent
definition in Section~\ref{sec:setup}, the three summands in equation (6) belong,
respectively, to the three generator families whose span is
\(\mathscr S_0\). Their finite sum therefore lies in \(\mathscr S_0\).
Proposition~\ref{prop:step-009-raw-witness} gives
\(W_0\perp\mathscr S_0\), which proves the second equality in equation (5).
No quotient representative, normalized target, or closure of the tangent
space is used. 

\end{proof}


\begin{proposition}[preserved positive raw coefficient margin]
\label{prop:step-013-preserved-raw-margin}
Under Propositions~\ref{prop:step-009-raw-witness},
\ref{prop:step-011-uniform-factor-radius}, and
\ref{prop:step-012-endpoint-remainder}, and
Lemma~\ref{lem:step-013-tangent-cancellation}, on
\(\mathcal E_{\rm init\_norm}\cap\mathcal C_{\rm path}\), define the exact
endpoint remainder
\[
  R_\infty:=\Psi_{A,B,C}(\theta_\infty)-\Psi_{A,B,C}(\theta_0)
  -D\Psi_{A,B,C}(\theta_0)[\Delta\theta].                \tag{7}
\]
Then
\[
  \left\langle D_r-\Psi_{A,B,C}(\theta_\infty),W_0\right\rangle_F
  \ge\frac{15}{16}\delta_0\sqrt r,                      \tag{8}
\]
and consequently
\[
  \left\|D_r-\Psi_{A,B,C}(\theta_\infty)\right\|_F
  \ge\frac{15}{16}\delta_0\sqrt r.                      \tag{9}
\]

\end{proposition}

\begin{proof}
Rearranging equation (7) gives the exact raw
coefficient identity
\[
  D_r-\Psi_{A,B,C}(\theta_\infty)
  =D_r-\Psi_{A,B,C}(\theta_0)
   -D\Psi_{A,B,C}(\theta_0)[\Delta\theta]-R_\infty.
  \tag{10}
\]
By equation (2), \(\Psi_{A,B,C}(\theta_0)=\widehat D_0\). Pairing equation
(10) with \(W_0\), using the exact cancellation in
Lemma~\ref{lem:step-013-tangent-cancellation}, the initial margin in equation
(1), and equation (4), gives
\[
\begin{aligned}
  \left\langle D_r-\Psi_{A,B,C}(\theta_\infty),W_0\right\rangle_F
  &=\langle D_r-\widehat D_0,W_0\rangle_F
    -\langle R_\infty,W_0\rangle_F\\
  &\ge\delta_0\sqrt r-|\langle R_\infty,W_0\rangle_F|\\
  &\ge\delta_0\sqrt r-\|R_\infty\|_F.
\end{aligned}\tag{11}
\]
No sign is assigned to the remainder pairing. The endpoint estimate
in equation (3) and \(r\ge1\) imply
\[
  \|R_\infty\|_F
  \le\frac{\delta_0}{16}
  \le\frac{\delta_0\sqrt r}{16}.                        \tag{12}
\]
Substituting equation (12) into equation (11) proves equation (8). Finally,
Frobenius Cauchy--Schwarz and \(\|W_0\|_F=1\) give
\[
  \left\|D_r-\Psi_{A,B,C}(\theta_\infty)\right\|_F
  \ge
  \left|\left\langle
    D_r-\Psi_{A,B,C}(\theta_\infty),W_0
  \right\rangle_F\right|
  \ge\frac{15}{16}\delta_0\sqrt r,
\]
because the pairing is positive by equation (8). This proves equation (9).

If \(E_{\rm path}=0\), Proposition~\ref{prop:step-011-uniform-factor-radius} gives
\(\theta_\infty=\theta_0\), so \(\Delta\theta=0\) and
\(R_\infty=0\). Equation (10) then retains the complete initial margin
\(\langle D_r-\widehat D_0,W_0\rangle_F\ge\delta_0\sqrt r\), rather than
only the conservative factor \(15/16\). If \(E_{\rm path}=E_\star\), the
non-strict bound in equation (3) still gives equation (12). If the
initial deficit inequality in equation (1) holds with equality, equation
(11) still yields the claimed floor after the worst-sign remainder. For
\(r=1\), the second inequality in equation (12) is equality, so the same
\(15/16\) conclusion holds without a large-r simplification. 

\end{proof}

\begin{proof}[Proof of the subsection conclusion]

Lemma~\ref{lem:step-013-tangent-cancellation} expands the exact raw
coefficient derivative at \(\theta_0\) and places each of its three families
inside a defining generator family of \(\mathscr S_0\). The raw
witness is orthogonal to this exact span, so the complete first-order term
cancels.

Proposition~\ref{prop:step-013-preserved-raw-margin} then uses the exact
identity \(\Psi_{A,B,C}(\theta_0)=\widehat D_0\), the initial raw
margin, and the single-endpoint remainder. It controls the remainder
pairing by absolute value and proves
\[
  \langle D_r-\Psi_{A,B,C}(\theta_\infty),W_0\rangle_F
  \ge\delta_0\sqrt r-\frac{\delta_0}{16}
  \ge\frac{15}{16}\delta_0\sqrt r.
\]
Unit norm of \(W_0\) converts this positive pairing into the same lower bound
for the exact raw coefficient residual norm. The argument is deterministic
on \(\mathcal E_{\rm init\_norm}\cap\mathcal C_{\rm path}\), uses one endpoint
remainder rather than a time sum, and makes no physical-coordinate transfer.

\end{proof}

\subsection{Same-target transfer to physical loss}\label{app:step-014}

\begin{proof}[Auxiliary facts in the present notation]


\medskip\noindent\textbf{realized-factor conditioning.}\par

Proposition~\ref{prop:step-001-realized-conditioning} supplies (1)
for the realized \(n\times r\) mode matrices in the physical operator norm.
The event inclusion
\(\mathcal E_{\rm init\_norm}\subseteq\mathcal E_{\rm cond}\) discharges its
use below. Its probability estimate is neither reused nor converted.


\medskip\noindent\textbf{preserved raw coefficient margin.}\par

Proposition~\ref{prop:step-013-preserved-raw-margin} supplies (2)
for the exact raw tensor \(D_r-\Psi_{A,B,C}(\theta_\infty)\) on the unchanged
conditional event. It is not a normalized, quotient, population, or
finite-time residual.


\medskip\noindent\textbf{Locally proved operator facts.}\par

Lemma~\ref{lem:step-014-mode-span-projection} proves the exact tensor-product
projection identity and Pythagorean decomposition in current notation.
Proposition~\ref{prop:step-014-positive-physical-loss} proves the required
tensor-product singular bounds directly from compact singular-value
decompositions and then applies them to (2). No external paper result is used.

\end{proof}



\begin{lemma}[exact mode-span projection of the physical residual]
\label{lem:step-014-mode-span-projection}
Under the basic setting definitions and Proposition~\ref{prop:step-001-realized-conditioning}, work on
\(\mathcal E_{\rm cond}\). For every \(\theta=(X,Y,Z)\), let
\[
  \mathcal P:=P_A\otimes P_B\otimes P_C,
  \qquad \mathcal L:=A\otimes B\otimes C.
\]
Then \(\mathcal P\) is the orthogonal projection onto
\(\operatorname{range}(A)\otimes\operatorname{range}(B)\otimes
\operatorname{range}(C)\),
\[
  \mathcal P(T-S(\theta))
  =\mathcal L(D_r-\Psi_{A,B,C}(\theta)),                   \tag{3}
\]
and
\[
  \|T-S(\theta)\|_F^2
  =\|\mathcal P(T-S(\theta))\|_F^2
   +\|(I-\mathcal P)(T-S(\theta))\|_F^2.                 \tag{4}
\]

\end{lemma}

\begin{proof}
By (1), every mode matrix has full column rank, so
\[
  P_M=M(M^{\mathsf T}M)^{-1}M^{\mathsf T}
\]
is self-adjoint and idempotent, satisfies \(P_MM=M\), and obeys
\(P_Mx=M(M^\dagger x)\). Tensoring self-adjointness and idempotence shows that
\(\mathcal P\) is the stated orthogonal projection.

Because the target factors are columns of \(A,B,C\),
\[
  \mathcal PT=T=\mathcal LD_r.                            \tag{5}
\]
For the model columns, with
\(\alpha_i=A^\dagger x_i\), \(\beta_i=B^\dagger y_i\), and
\(\gamma_i=C^\dagger z_i\),
\[
\begin{aligned}
  \mathcal P S(\theta)
  &=\sum_{i=1}^k P_Ax_i\otimes P_By_i\otimes P_Cz_i\\
  &=\sum_{i=1}^k A\alpha_i\otimes B\beta_i\otimes C\gamma_i
   =\mathcal L\Psi_{A,B,C}(\theta).
\end{aligned}\tag{6}
\]
Subtracting (6) from (5) proves (3). The tensor in parentheses on the right
of (3) lies in \((\mathbb R^r)^{\otimes3}\), while both sides lie in
\((\mathbb R^n)^{\otimes3}\), so the mode dimensions agree.

Finally, \(\mathcal P(T-S(\theta))\) and
\((I-\mathcal P)(T-S(\theta))\) are orthogonal because
\(\mathcal P(I-\mathcal P)=0\). Pythagoras gives (4). Thus the only discarded
quantity is the explicitly nonnegative second squared norm in (4).


\end{proof}


\begin{proposition}[positive relative physical loss]
\label{prop:step-014-positive-physical-loss}
\leavevmode\par\noindent
Under Propositions~\ref{prop:step-001-realized-conditioning} and
\ref{prop:step-013-preserved-raw-margin}, and
Lemma~\ref{lem:step-014-mode-span-projection}, work on
\(\mathcal E_{\rm init\_norm}\cap\mathcal C_{\rm path}\). Then, for every
\(H\in(\mathbb R^r)^{\otimes3}\),
\[
  \kappa_1^{-3}\|H\|_F
  \le\|(A\otimes B\otimes C)H\|_F
  \le\kappa_1^3\|H\|_F.                                  \tag{7}
\]
Moreover, (T2)--(T3) hold, with
\(\epsilon_0(\kappa)>0\) and \(\|T\|_F>0\).

\end{proposition}

\begin{proof}
Take compact singular-value decompositions
\(M=U_M\Sigma_MV_M^{\mathsf T}\) for \(M=A,B,C\). Then
\[
\begin{aligned}
  A\otimes B\otimes C
  ={}&(U_A\otimes U_B\otimes U_C)
      (\Sigma_A\otimes\Sigma_B\otimes\Sigma_C)\\
     &\quad\cdot(V_A\otimes V_B\otimes V_C)^{\mathsf T}.
\end{aligned}\tag{8}
\]
The first factor is an isometry into the physical tensor space, the last is
orthogonal, and the middle diagonal entries are
\(\sigma_i(A)\sigma_j(B)\sigma_\ell(C)\). Thus (1) gives (7), including the
exact lower product \(\kappa_1^{-3}\) and upper product \(\kappa_1^3\).

The \(r\) tensors \(e_j\otimes e_j\otimes e_j\) are orthonormal, so
\(\|D_r\|_F=\sqrt r\). Since
\(T=(A\otimes B\otimes C)D_r\), equation (7) gives
\[
  0<\kappa_1^{-3}\sqrt r
  \le\|T\|_F\le\kappa_1^3\sqrt r.                         \tag{9}
\]
In particular, (T2)'s target upper bound holds and the target is nonzero.

Apply (3)--(4) at \(\theta=\theta_\infty\), then (7), and finally the
raw margin (2):
\[
\begin{aligned}
  \|T-S(\theta_\infty)\|_F
  &\ge\|\mathcal P(T-S(\theta_\infty))\|_F\\
  &=\|(A\otimes B\otimes C)
      (D_r-\Psi_{A,B,C}(\theta_\infty))\|_F\\
  &\ge\kappa_1^{-3}
      \|D_r-\Psi_{A,B,C}(\theta_\infty)\|_F\\
  &\ge\kappa_1^{-3}\frac{15}{16}\delta_0\sqrt r.
\end{aligned}\tag{10}
\]
This proves the residual part of (T2). Combining (9) and (10) in the required
direction gives
\[
  \|T-S(\theta_\infty)\|_F
  \ge\frac{15}{16}\delta_0\kappa_1^{-6}\|T\|_F.          \tag{11}
\]
Both sides are nonnegative. Squaring (11) and using the exact definition of
\(F\) proves (T3). Since \(\delta_0=1/8\) and
\(\kappa_1=2\kappa^2>0\), \(\epsilon_0(\kappa)>0\); (9) then makes the
right side of (T3) strictly positive.

At equality in the allowed singular bounds, all displayed inequalities remain
valid with a positive lower product. If the orthogonal term in (4) vanishes,
the projected term still has the floor (10); if it is nonzero, it only
increases the physical norm. A zero raw residual is excluded by (2), and a
zero target is excluded by (9). Zero model columns cause no exception because
(5)--(6) use only linear projections. At \(E_{\rm path}=0\), Proposition~\ref{prop:step-013-preserved-raw-margin} retains the full raw margin \(\delta_0\sqrt r\), so the same
calculation preserves the stronger baseline constant
\(\delta_0^2\kappa_1^{-12}\). 

\end{proof}

\begin{proof}[Proof of the subsection conclusion]

Lemma~\ref{lem:step-014-mode-span-projection} proves (T1) directly for the
actual residual and shows by (4) that projection discards only a nonnegative
orthogonal squared norm. Proposition~\ref{prop:step-014-positive-physical-loss}
combines the raw floor with the exact tensor singular products,
proves both bounds in (T2), checks the normalization direction before
squaring, and proves (T3). The event scope, target, residual, Frobenius norm,
and conditional mode are unchanged.

\end{proof}

\subsection{Continuity and conditional probability}\label{app:step-015}

\begin{proof}[Auxiliary facts in the present notation]


\medskip\noindent\textbf{initialization confidence.}\par

Proposition~\ref{prop:step-010-public-confidence} proves, under the
joint smoothing/initialization law conditional on the deterministic base
triple,
\[
  \mathbb P(\mathcal E_{\rm init\_norm})\ge1-r^{-10}.
\]
It uses the exact event from the current setting. For the sufficiently large
\(r\) regime, enlarged if necessary so that \(r\ge2\), this lower bound is
strictly positive. No independence between
\(\mathcal E_{\rm init\_norm}\) and \(\mathcal C_{\rm path}\) is supplied or
used.


\medskip\noindent\textbf{finite path limit.}\par

Lemma~\ref{lem:step-011-finite-path-limit} proves, from the sole
finite-total-variation hypothesis \(\mathcal C_{\rm path}\), that the exact
balanced representatives satisfy
\[
  d_{\rm bal}(\theta_t,\theta_\infty)\longrightarrow0
\]
for a finite \(\theta_\infty\in(\mathbb R^{n\times k})^3\). The convergence
is in the same product Frobenius metric used below, not in a quotient or a
surrogate topology.


\medskip\noindent\textbf{positive endpoint loss.}\par

Proposition~\ref{prop:step-014-positive-physical-loss} proves on
\(\mathcal E_{\rm init\_norm}\cap\mathcal C_{\rm path}\) that
\[
  F(\theta_\infty)
  \ge\left(\frac{15}{16}\delta_0\right)^2
      \kappa_1^{-12}\|T\|_F^2>0.
\]
This is the setting's actual physical objective at the endpoint. It
contains no surrogate target, normalized loss, or omitted residual component.


\medskip\noindent\textbf{Finite-dimensional polynomial continuity.}\par

We use the following statement in the present notation. A map between finite-dimensional Euclidean spaces is
continuous when each output coordinate is a finite sum of finite products of
input coordinates and fixed real coefficients. Consequently, if
\(h_t\to h\) in the Euclidean norm, then the values of such a polynomial map
converge to its value at \(h\).

The parameter space is
\((\mathbb R^{n\times k})^3\), its Euclidean product distance is exactly
\(d_{\rm bal}\), and the output spaces are the finite-dimensional tensor
space \((\mathbb R^n)^{\otimes3}\) and \(\mathbb R\). Proposition~\ref{prop:step-015-polynomial-continuity}
displays every relevant coordinate polynomial explicitly.


\medskip\noindent\textbf{Conditional probability on a positive-probability event.}\par

We use the following statement in the present notation. If \(A,B\) are events and \(\mathbb P(A)>0\), then
\[
  \mathbb P(B\mid A):=\frac{\mathbb P(A\cap B)}{\mathbb P(A)},
  \qquad
  \mathbb P(A\cap B)=\mathbb P(A)\mathbb P(B\mid A).
\]
Moreover, event inclusion implies probability monotonicity.

Take
\(A=\mathcal E_{\rm init\_norm}\) and
\(B=\mathcal C_{\rm path}\). Positivity of \(\mathbb P(A)\) is proved
before the definition is used in
Proposition~\ref{prop:step-015-conditional-accounting}. Proposition~\ref{prop:step-015-event-inclusion}
supplies the inclusion needed for monotonicity.

\end{proof}



\begin{proposition}[polynomial continuity of the CP model and objective]
\label{prop:step-015-polynomial-continuity}
Under the basic setting definitions, fix any realization of the target
\(T\). The map
\[
  S:(\mathbb R^{n\times k})^3\longrightarrow
    (\mathbb R^n)^{\otimes3}
\]
is a homogeneous polynomial map of degree three, and
\(F:(\mathbb R^{n\times k})^3\to\mathbb R\) is a polynomial of degree at
most six. Both maps are continuous in the Euclidean product metric
\(d_{\rm bal}\). In particular, if
\(d_{\rm bal}(\theta_t,\theta_\infty)\to0\), then
\[
  \|S(\theta_t)-S(\theta_\infty)\|_F\to0,
  \qquad
  F(\theta_t)\to F(\theta_\infty).                       \tag{1}
\]

\end{proposition}

\begin{proof}
Write tensor coordinates as
\(1\le a,b,c\le n\). From the setting definition,
\[
  [S(X,Y,Z)]_{abc}
  =\sum_{i=1}^k X_{ai}Y_{bi}Z_{ci}.                     \tag{2}
\]
Thus every coordinate of \(S\) is a finite sum of degree-three monomials in
the entries of \((X,Y,Z)\). Hence \(S\) is a homogeneous polynomial map of
degree three.

For the fixed realized target, the objective has the coordinate expression
\[
  F(X,Y,Z)
  =\sum_{a=1}^n\sum_{b=1}^n\sum_{c=1}^n
    \left(T_{abc}-\sum_{i=1}^kX_{ai}Y_{bi}Z_{ci}\right)^2. \tag{3}
\]
Each summand is the square of a polynomial of degree at most three, so (3) is
a polynomial of degree at most six. Coordinate projections are continuous,
and finite sums and products of continuous real-valued functions are
continuous. Equations (2)--(3), together with the finite number of tensor
coordinates, therefore prove continuity of both maps.

Finally, the setting identity
\[
  d_{\rm bal}(\theta,\theta')
  =\bigl(\|X-X'\|_F^2+\|Y-Y'\|_F^2+\|Z-Z'\|_F^2\bigr)^{1/2}
\]
is exactly the Euclidean norm distance on the parameter space. Applying the
just-proved continuity to
\(d_{\rm bal}(\theta_t,\theta_\infty)\to0\) proves (1).


\end{proof}


\begin{proposition}[endpoint continuity and exact event inclusion]
\label{prop:step-015-event-inclusion}
\leavevmode\par\noindent
Under Assumption~\ref{assump:gd_step}, Lemma~\ref{lem:step-011-finite-path-limit}, Proposition~\ref{prop:step-014-positive-physical-loss}, and
Proposition~\ref{prop:step-015-polynomial-continuity}, every outcome in
\(\mathcal E_{\rm init\_norm}\cap\mathcal C_{\rm path}\) satisfies
\[
  d_{\rm bal}(\theta_t,\theta_\infty)\to0,
  \qquad
  \lim_{t\to\infty}F(\theta_t)
  =F(\theta_\infty)
  \ge\epsilon_0(\kappa)\|T\|_F^2>0.                    \tag{4}
\]
Consequently,
\[
  \mathcal E_{\rm init\_norm}\cap\mathcal C_{\rm path}
  \subseteq\mathcal F_+.                                \tag{5}
\]

\end{proposition}

\begin{proof}
Fix an arbitrary outcome in
\(\mathcal E_{\rm init\_norm}\cap\mathcal C_{\rm path}\). Lemma~\ref{lem:step-011-finite-path-limit} produces a finite
\(\theta_\infty\) and proves
\(d_{\rm bal}(\theta_t,\theta_\infty)\to0\). Applying
Proposition~\ref{prop:step-015-polynomial-continuity} to this exact
convergence yields
\[
  \lim_{t\to\infty}F(\theta_t)=F(\theta_\infty).        \tag{6}
\]
On the same outcome, Proposition~\ref{prop:step-014-positive-physical-loss} gives
\[
  F(\theta_\infty)
  \ge\epsilon_0(\kappa)\|T\|_F^2>0.                    \tag{7}
\]
Combining (6)--(7) proves (4).

The setting-defined event \(\mathcal F_+\) is precisely the event that the
balanced-GD representatives converge in \(d_{\rm bal}\) to a finite limit
and that their objective limit obeys the positive relative floor in (4).
Thus (4) is exactly membership in \(\mathcal F_+\), not membership in a
surrogate endpoint or finite-time event. Since the chosen outcome was
arbitrary, (5) follows. 

\end{proof}


\begin{proposition}[exact conditional-probability accounting]
\label{prop:step-015-conditional-accounting}
Under Assumption~\ref{assump:dimension}, Proposition~\ref{prop:step-010-public-confidence}, and
Proposition~\ref{prop:step-015-event-inclusion}, work under the setting's
joint smoothing/initialization law conditional on an arbitrary admissible
deterministic base triple. For \(r\ge2\),
\(\mathbb P(\mathcal E_{\rm init\_norm})>0\), the conditional probability
\(\mathbb P(\mathcal C_{\rm path}\mid
\mathcal E_{\rm init\_norm})\) is well defined by event conditioning, and
all three relations in (T) hold. They remain valid when this conditional
probability is zero.

\end{proposition}

\begin{proof}
Proposition~\ref{prop:step-010-public-confidence} gives
\[
  \mathbb P(\mathcal E_{\rm init\_norm})\ge1-r^{-10}.
\]
For \(r\ge2\), \(r^{-10}<1\), and therefore
\[
  \mathbb P(\mathcal E_{\rm init\_norm})
  \ge1-r^{-10}>0.                                       \tag{8}
\]
Assumption~\ref{assump:dimension} permits enlarging the sufficiently-large
\(r\) threshold to include \(r\ge2\), so (8) adds no theorem-facing
restriction. It also verifies the positive denominator before conditioning.

By Proposition~\ref{prop:step-015-event-inclusion} and monotonicity of
probability,
\[
  \mathbb P(\mathcal F_+)
  \ge\mathbb P(\mathcal E_{\rm init\_norm}
    \cap\mathcal C_{\rm path}).                         \tag{9}
\]
Using (8), the definition of conditional probability for the two exact
setting events gives the identity
\[
  \mathbb P(\mathcal E_{\rm init\_norm}
    \cap\mathcal C_{\rm path})
  =\mathbb P(\mathcal E_{\rm init\_norm})
    \mathbb P(\mathcal C_{\rm path}
      \mid\mathcal E_{\rm init\_norm}).                 \tag{10}
\]
The conditional probability in (10) is nonnegative. Multiplying the lower bound for \(\mathbb P(\mathcal E_{\rm init\_norm})\) by this
nonnegative factor preserves the inequality direction:
\[
\begin{aligned}
  \mathbb P(\mathcal E_{\rm init\_norm})
    \mathbb P(\mathcal C_{\rm path}
      \mid\mathcal E_{\rm init\_norm})
  \ge(1-r^{-10})
    \mathbb P(\mathcal C_{\rm path}
      \mid\mathcal E_{\rm init\_norm}).
\end{aligned}\tag{11}
\]
Equations (9)--(11) prove (T).

If
\(\mathbb P(\mathcal C_{\rm path}\mid
\mathcal E_{\rm init\_norm})=0\), equation (10) says that the intersection
has probability zero and (11) has zero on both sides. Thus the theorem still
asserts only \(\mathbb P(\mathcal F_+)\ge0\); it does not acquire an
unconditional positive failure probability or any lower bound on the path
event. No independence assertion is used anywhere. 

\end{proof}

\begin{proof}[Proof of the subsection conclusion]

Proposition~\ref{prop:step-015-polynomial-continuity} proves directly from
the CP coordinates that \(S\) and \(F\) are finite-dimensional polynomial
maps and converts the \(d_{\rm bal}\)-limit into
\[
  \lim_{t\to\infty}F(\theta_t)=F(\theta_\infty).
\]
Lemma~\ref{lem:step-011-finite-path-limit} supplies that finite
limit on the exact path event, while Proposition~\ref{prop:step-014-positive-physical-loss} supplies the actual
physical endpoint floor. Proposition~\ref{prop:step-015-event-inclusion}
combines them to obtain the exact inclusion
\[
  \mathcal E_{\rm init\_norm}\cap\mathcal C_{\rm path}
  \subseteq\mathcal F_+.
\]

Proposition~\ref{prop:step-010-public-confidence} and the explicit
large-\(r\) check in
Proposition~\ref{prop:step-015-conditional-accounting} show that the
conditioning event has positive probability. The same proposition then
uses only event monotonicity and the definition of conditional probability to
prove
\[
  \boxed{
  \mathbb P(\mathcal F_+)
  \ge(1-r^{-10})
    \mathbb P(\mathcal C_{\rm path}\mid
      \mathcal E_{\rm init\_norm})}.
\]
The probability space, all-time/asymptotic mode, physical objective, and
setting events are unchanged, and the unresolved conditional factor is
retained exactly even at its zero boundary.

\end{proof}

\subsection{Proof of the Main Theorem}\label{app:proof-main}

\begin{proof}[Proof of Theorem~\ref{thm:main}]
Proposition~\ref{prop:step-010-public-confidence} proves
$\mathbb P(\mathcal E_{\rm init\_norm})\geq1-r^{-10}$.
Lemma~\ref{lem:step-011-finite-path-limit} and
Propositions~\ref{prop:step-014-positive-physical-loss} and
\ref{prop:step-015-event-inclusion} prove, on
$\mathcal E_{\rm init\_norm}\cap\mathcal C_{\rm path}$,
the finite parameter limit and the positive asymptotic physical-loss floor
with $\epsilon_0(\kappa)=((15/16)\delta_0)^2\kappa_1^{-12}$.
Finally, Proposition~\ref{prop:step-015-conditional-accounting} proves the
exact event-conditioning identity and
\[
  \mathbb P(\mathcal F_+)
  \geq(1-r^{-10})
  \mathbb P(\mathcal C_{\rm path}\mid
    \mathcal E_{\rm init\_norm}).
\]
That proposition also verifies the positive conditioning denominator and
retains the conditional factor when it equals zero.  These are precisely all
claims in Theorem~\ref{thm:main}.
\end{proof}
